\documentclass[12pt,a4paper]{article}
\usepackage[utf8]{inputenc}
\usepackage[turkish]{babel}
\usepackage[T1]{fontenc}
\usepackage{amsmath, amssymb}
\usepackage[hidelinks]{hyperref}
\usepackage{geometry}
\usepackage{amssymb}

\usepackage[most]{tcolorbox}
\usepackage{xcolor}

\definecolor{testblue}{RGB}{30, 144, 255}
\definecolor{ansgreen}{RGB}{34, 139, 34}

\newtcolorbox{testcard}[1]{
    enhanced,
    colback=white,
    colframe=testblue!70!white,
    boxrule=0.6pt,
    arc=0pt,
    outer arc=0pt,
    title={\textbf{#1}},
    coltitle=white,
    colbacktitle=testblue!80!black,
    fonttitle=\large\bfseries,
    left=10pt, right=10pt, top=10pt, bottom=10pt,
    segmentation style={draw=testblue!30!white, solid, line width=1pt},
    lower separated=true
}

\geometry{margin=0.8in}

\title{\textbf{İş Sağlığı ve Güvenliği (İSG) Test Bankası}}
\author{Oğuz Taşdemir}
\date{\today}

\begin{document}

\maketitle
\newpage

\tableofcontents
\newpage

\section{Giriş}
Bu doküman, İş Sağlığı ve Güvenliği (İSG) konularını kapsayan detaylı ve açıklamalı soru bankası setidir. Her soru, sınav müfredatına ve hocanın ipuçlarına uygun olarak hazırlanmıştır.

\newpage
\section{Test Soruları}

\subsection{Ünite 1: İSG'ye Giriş ve Temel Kavramlar}
\begin{testcard}{Soru 1}
    İşyerinde bulunan tehlikelerin zarar verme olasılığı ile bu zararın derecesinin bileşkesine ne ad verilir?
    \begin{itemize}
        \item[A)] Tehlike
        \item[B)] Risk
        \item[C)] Kaza
        \item[D)] Olay
        \item[E)] Hiçbiri
    \end{itemize}
    \tcblower
    \textbf{Doğru Cevap: B} \\
    \textbf{Açıklama:} 6331 sayılı İSG Kanunu'na göre Tehlike potansiyeldir, Risk ise bu tehlikeden kaynaklanacak zarar veya kayıp oluşma olasılığı ile şiddetinin bileşkesidir.
\end{testcard}

\begin{testcard}{Soru 2}
    Aşağıdakilerden hangisi İSG'nin temel amaçları arasında birinci önceliğe sahiptir?
    \begin{itemize}
        \item[A)] İşletmenin güvenliğini sağlamak
        \item[B)] Çalışanları korumak
        \item[C)] Üretim ve ürün kalitesini artırmak
        \item[D)] Maliyetleri düşürmek
        \item[E)] Hiçbiri
    \end{itemize}
    \tcblower
    \textbf{Doğru Cevap: B} \\
    \textbf{Açıklama:} İSG'nin en önemli ve birinci öncelikli amacı çalışanların sağlığını, güvenliğini ve yaşam hakkını korumaktır.
\end{testcard}

\begin{testcard}{Soru 3}
    Aşağıdakilerden hangisi proaktif İSG yaklaşımının özelliklerinden biridir?
    \begin{itemize}
        \item[A)] Kaza gerçekleştikten sonra nedenlerini araştırmak
        \item[B)] Tehlikeleri önceden öngörüp riskler gerçekleşmeden önlem almak
        \item[C)] Sadece mevzuat cezalarından kaçınmak için çalışmak
        \item[D)] İş kazası raporlarını düzenlemek
        \item[E)] Hiçbiri
    \end{itemize}
    \tcblower
    \textbf{Doğru Cevap: B} \\
    \textbf{Açıklama:} Proaktif yaklaşım, tehlikeyi henüz kaza oluşmadan tahmin edip önleyici tedbirler alarak koruma sağlamayı hedefler.
\end{testcard}

\begin{testcard}{Soru 4}
    İSG bakımından reaktif yaklaşım neyi ifade eder?
    \begin{itemize}
        \item[A)] Kazalar veya meslek hastalıkları oluştuktan sonra önlem almayı
        \item[B)] İş öncesi risk analizi yapmayı
        \item[C)] Çalışan temsilcileriyle sürekli toplantı düzenlemeyi
        \item[D)] Tehlikeleri kaynağından yok etmeyi
        \item[E)] Hiçbiri
    \end{itemize}
    \tcblower
    \textbf{Doğru Cevap: A} \\
    \textbf{Açıklama:} Reaktif yaklaşım olay gerçekleştikten sonra tepki veren, kaza sonrası önlem almaya odaklanan geleneksel anlayıştır.
\end{testcard}

\begin{testcard}{Soru 5}
    İSG çalışmalarında 'ramak kala olay' ne demektir?
    \begin{itemize}
        \item[A)] Maddi hasarla sonuçlanan büyük kaza
        \item[B)] Yaralanma veya hasara yol açma potansiyeli olduğu halde zarar vermeyen olay
        \item[C)] İşçinin sürekli iş göremez hale gelmesi
        \item[D)] Meslek hastalığı teşhisi
        \item[E)] Hiçbiri
    \end{itemize}
    \tcblower
    \textbf{Doğru Cevap: B} \\
    \textbf{Açıklama:} Ramak kala olay, herhangi bir yaralanma veya hasara yol açmayan ancak bir adım sonrasında kazaya sebep olabilecek olaylardır.
\end{testcard}

\begin{testcard}{Soru 6}
    Aşağıdakilerden hangisi doğrudan bir tehlike kaynağıdır?
    \begin{itemize}
        \item[A)] İşçinin parmağının kesilmesi
        \item[B)] Açıkta duran yüksek gerilim kablosu
        \item[C)] İşitme kaybı yaşanması
        \item[D)] İş kazası tazminatı ödenmesi
        \item[E)] Hiçbiri
    \end{itemize}
    \tcblower
    \textbf{Doğru Cevap: B} \\
    \textbf{Açıklama:} Açıkta duran yüksek gerilim kablosu zarar verme potansiyeline sahip bir kaynak olup doğrudan bir tehlikedir. Diğerleri ise risktir veya kaza sonucudur.
\end{testcard}

\begin{testcard}{Soru 7}
    Risk kontrol hiyerarşisinde 'Eliminasyon' (Kaynaktan yok etme) neyi ifade eder?
    \begin{itemize}
        \item[A)] Tehlikeli olan makine veya prosesi tamamen ortadan kaldırmayı
        \item[B)] Tehlikeli maddeyi daha az tehlikeli olanla değiştirmeyi
        \item[C)] Çalışanlara maske ve eldiven dağıtmayı
        \item[D)] Uyarı levhaları asmayı
        \item[E)] Hiçbiri
    \end{itemize}
    \tcblower
    \textbf{Doğru Cevap: A} \\
    \textbf{Açıklama:} Eliminasyon, tehlikeyi kaynağında tamamen yok ederek riski sıfırlayan en etkili kontrol yöntemidir.
\end{testcard}

\begin{testcard}{Soru 8}
    Risk kontrol hiyerarşisinde 'İkame' (Substitution) neyi ifade eder?
    \begin{itemize}
        \item[A)] Çalışanı başka bir çalışanla değiştirmeyi
        \item[B)] Tehlikeli olan bir maddeyi daha az tehlikeli veya tehlikesiz olanla değiştirmeyi
        \item[C)] Koruyucu siperlik takmayı
        \item[D)] İş süresini kısaltmayı
        \item[E)] Hiçbiri
    \end{itemize}
    \tcblower
    \textbf{Doğru Cevap: B} \\
    \textbf{Açıklama:} İkame, tehlikeli olan nesnenin veya yöntemin yerine daha az tehlikeli bir alternatif getirilmesidir.
\end{testcard}

\begin{testcard}{Soru 9}
    Aşağıdakilerden hangisi risk kontrol hiyerarşisinde en son sırada yer alır ve en az koruma sağlar?
    \begin{itemize}
        \item[A)] Mühendislik önlemleri
        \item[B)] İdari önlemler
        \item[C)] Kişisel Koruyucu Donanım (KKD) kullanımı
        \item[D)] İkame
        \item[E)] Hiçbiri
    \end{itemize}
    \tcblower
    \textbf{Doğru Cevap: C} \\
    \textbf{Açıklama:} KKD kullanımı, risk kontrol hiyerarşisinde en son çaredir ve risk kaynağında çözülemediğinde kişiyi korumak için kullanılır.
\end{testcard}

\begin{testcard}{Soru 10}
    İşyerinde kişisel koruyucu donanımların (KKD) temin maliyetiyle ilgili hangisi yasal kuraldır?
    \begin{itemize}
        \item[A)] Maliyet tamamen işveren tarafından karşılanmalı, çalışana yansıtılmamalıdır
        \item[B)] İşçi ve işveren yarı yarıya öder
        \item[C)] Sendika aidatlarından karşılanır
        \item[D)] İşçi maaşından kesinti yapılarak tahsil edilir
        \item[E)] Hiçbiri
    \end{itemize}
    \tcblower
    \textbf{Doğru Cevap: A} \\
    \textbf{Açıklama:} 6331 sayılı İSG Kanunu uyarınca İSG tedbirlerinin ve KKD'lerin maliyeti hiçbir şekilde çalışanlara yansıtılamaz, işveren ücretsiz sağlamak zorundadır.
\end{testcard}

\begin{testcard}{Soru 11}
    Çalışanların temel İSG eğitimleri yönetmeliğine göre, 'Çok Tehlikeli' sınıftaki bir işyerinde eğitim süresi asgari kaç saattir?
    \begin{itemize}
        \item[A)] En az 8 saat
        \item[B)] En az 12 saat
        \item[C)] En az 16 saat
        \item[D)] En az 20 saat
        \item[E)] Hiçbiri
    \end{itemize}
    \tcblower
    \textbf{Doğru Cevap: C} \\
    \textbf{Açıklama:} Çok Tehlikeli sınıftaki işyerlerinde çalışanların alacağı temel İSG eğitimi en az 16 saat olmak zorundadır.
\end{testcard}

\begin{testcard}{Soru 12}
    Temel İSG eğitimleri yönetmeliğine göre, 'Tehlikeli' sınıftaki bir işyerinde eğitim süresi asgari kaç saattir?
    \begin{itemize}
        \item[A)] En az 8 saat
        \item[B)] En az 12 saat
        \item[C)] En az 16 saat
        \item[D)] En az 24 saat
        \item[E)] Hiçbiri
    \end{itemize}
    \tcblower
    \textbf{Doğru Cevap: B} \\
    \textbf{Açıklama:} Tehlikeli sınıftaki işyerlerinde çalışanların temel İSG eğitimi asgari 12 saat olarak belirlenmiştir.
\end{testcard}

\begin{testcard}{Soru 13}
    Temel İSG eğitimleri yönetmeliğine göre, 'Az Tehlikeli' sınıftaki bir işyerinde eğitim süresi asgari kaç saattir?
    \begin{itemize}
        \item[A)] En az 8 saat
        \item[B)] En az 10 saat
        \item[C)] En az 12 saat
        \item[D)] En az 16 saat
        \item[E)] Hiçbiri
    \end{itemize}
    \tcblower
    \textbf{Doğru Cevap: A} \\
    \textbf{Açıklama:} Az Tehlikeli sınıftaki işyerlerinde temel İSG eğitimi asgari 8 saattir.
\end{testcard}

\begin{testcard}{Soru 14}
    Çok Tehlikeli sınıftaki bir işyerinde temel İSG eğitimleri en geç hangi periyotla yenilenmelidir?
    \begin{itemize}
        \item[A)] Yılda en az bir defa
        \item[B)] 2 yılda en az bir defa
        \item[C)] 3 yılda en az bir defa
        \item[D)] 5 yılda en az bir defa
        \item[E)] Hiçbiri
    \end{itemize}
    \tcblower
    \textbf{Doğru Cevap: A} \\
    \textbf{Açıklama:} Çok Tehlikeli sınıftaki işyerlerinde temel İSG eğitimlerinin yılda en az bir defa yenilenmesi zorunludur.
\end{testcard}

\begin{testcard}{Soru 15}
    Tehlikeli sınıftaki bir işyerinde temel İSG eğitimleri en geç hangi periyotla yenilenmelidir?
    \begin{itemize}
        \item[A)] Yılda bir
        \item[B)] 2 yılda bir
        \item[C)] 3 yılda bir
        \item[D)] 4 yılda bir
        \item[E)] Hiçbiri
    \end{itemize}
    \tcblower
    \textbf{Doğru Cevap: B} \\
    \textbf{Açıklama:} Tehlikeli sınıftaki işyerlerinde temel İSG eğitimleri en geç 2 yılda bir tekrarlanmalıdır.
\end{testcard}

\begin{testcard}{Soru 16}
    Az Tehlikeli sınıftaki bir işyerinde temel İSG eğitimleri en geç hangi periyotla yenilenmelidir?
    \begin{itemize}
        \item[A)] Yılda bir
        \item[B)] 2 yılda bir
        \item[C)] 3 yılda bir
        \item[D)] 5 yılda bir
        \item[E)] Hiçbiri
    \end{itemize}
    \tcblower
    \textbf{Doğru Cevap: C} \\
    \textbf{Açıklama:} Az Tehlikeli sınıfta temel İSG eğitimlerinin en geç 3 yılda bir yenilenmesi yasal gerekliliktir.
\end{testcard}

\begin{testcard}{Soru 17}
    İş kazalarının nedenlerini açıklayan Heinrich oranına göre, kazaların yüzde kaçı 'tehlikeli davranışlardan' kaynaklanır?
    \begin{itemize}
        \item[A)] Yüzde 50
        \item[B)] Yüzde 70
        \item[C)] Yüzde 88
        \item[D)] Yüzde 98
        \item[E)] Hiçbiri
    \end{itemize}
    \tcblower
    \textbf{Doğru Cevap: C} \\
    \textbf{Açıklama:} Heinrich'in kaza teorisine göre iş kazalarının \%88'i tehlikeli hareketlerden (insan kusuru), \%10'u tehlikeli durumlardan ve \%2'si kaçınılmaz nedenlerden kaynaklanır.
\end{testcard}

\begin{testcard}{Soru 18}
    Aşağıdakilerden hangisi bir tehlikeli davranış (güvensiz hareket) örneğidir?
    \begin{itemize}
        \item[A)] Makinenin dönen dişli kısmında koruyucu kapak olmaması
        \item[B)] İşçinin yüksekte çalışırken baret ve emniyet kemeri takmaması
        \item[C)] İşyeri zemininde yağ döküntüsü bulunması
        \item[D)] Elektrik kablolarının yalıtımının sıyrılmış olması
        \item[E)] Hiçbiri
    \end{itemize}
    \tcblower
    \textbf{Doğru Cevap: B} \\
    \textbf{Açıklama:} Çalışanın kişisel güvenlik kurallarını ihlal etmesi veya KKD kullanmaması güvensiz/tehlikeli davranıştır. Diğer şıklar çevreye ait tehlikeli durumlardır.
\end{testcard}

\begin{testcard}{Soru 19}
    Aşağıdakilerden hangisi bir tehlikeli durum (güvensiz koşul) örneğidir?
    \begin{itemize}
        \item[A)] Forklifti aşırı hızlı kullanmak
        \item[B)] Merdiven kenarında korkuluk bulunmaması
        \item[C)] Baret takmadan şantiye sahasına girmek
        \item[D)] Çalışırken şakalaşmak
        \item[E)] Hiçbiri
    \end{itemize}
    \tcblower
    \textbf{Doğru Cevap: B} \\
    \textbf{Açıklama:} Fiziki ortamın veya ekipmanların korumasız, arızalı veya güvensiz olması tehlikeli durumdur. Korkuluksuz merdiven fiziksel bir tehlikeli durumdur.
\end{testcard}

\begin{testcard}{Soru 20}
    Türkiye'de müstakil olarak yürürlüğe giren ve tüm iş kollarını kapsayan temel İSG Kanunu hangisidir?
    \begin{itemize}
        \item[A)] 4857 Sayılı İş Kanunu
        \item[B)] 6331 Sayılı İş Sağlığı ve Güvenliği Kanunu
        \item[C)] 5510 Sayılı Sosyal Sigortalar Kanunu
        \item[D)] 1593 Sayılı Umumi Hıfzıssıhha Kanunu
        \item[E)] Hiçbiri
    \end{itemize}
    \tcblower
    \textbf{Doğru Cevap: B} \\
    \textbf{Açıklama:} 2012 yılında kabul edilen 6331 Sayılı Kanun, Türkiye'nin ilk müstakil ve kapsamlı İş Sağlığı ve Güvenliği Kanunu'dur.
\end{testcard}

\begin{testcard}{Soru 21}
    TS ISO 45001 standardına göre tehlike nasıl tanımlanır?
    \begin{itemize}
        \item[A)] Hasar ve yaralanmaya yol açabilecek kaynak, durum veya işlem
        \item[B)] Bir olayın gerçekleşme olasılığı ile şiddetinin bileşimi
        \item[C)] Sadece ölümcül iş kazaları
        \item[D)] İşyeri hekiminin yazdığı reçete sayısı
        \item[E)] Hiçbiri
    \end{itemize}
    \tcblower
    \textbf{Doğru Cevap: A} \\
    \textbf{Açıklama:} ISO 45001 standardı tehlikeyi; yaralanma veya sağlık bozulmasına neden olabilecek kaynak, durum veya eylem/işlem olarak tanımlar.
\end{testcard}

\begin{testcard}{Soru 22}
    Bir tehlike veya riskin algılanmasını etkileyen öznel faktörler arasında aşağıdakilerden hangisi yer almaz?
    \begin{itemize}
        \item[A)] Kişinin yaşı ve deneyimi
        \item[B)] İşyerinin sermaye büyüklüğü
        \item[C)] Korkutuculuk ve anlaşılabilirlik düzeyi
        \item[D)] Risk altındaki kişi sayısı
        \item[E)] Hiçbiri
    \end{itemize}
    \tcblower
    \textbf{Doğru Cevap: B} \\
    \textbf{Açıklama:} Tehlike algılaması kişisel tecrübe, yaş, cinsiyet ve riskin doğasından (korkutuculuk vb.) etkilenir; işyerinin sermaye büyüklüğü risk algısını belirleyen öznel bir faktör değildir.
\end{testcard}

\begin{testcard}{Soru 23}
    Dağcıların ve sirk cambazlarının risk algısının düşük olması, İSG ilkelerine göre ne tür bir tehlike doğurur?
    \begin{itemize}
        \item[A)] Önlemlerin göz ardı edilmesine ve aşırı özgüvene yol açabilir
        \item[B)] Kazaları tamamen sıfırlar
        \item[C)] Ekipman kullanım maliyetini düşürür
        \item[D)] Çalışma verimini artırır
        \item[E)] Hiçbiri
    \end{itemize}
    \tcblower
    \textbf{Doğru Cevap: A} \\
    \textbf{Açıklama:} Düşük risk algısı ve aşırı deneyim kaynaklı özgüven, güvenlik önlemlerinin kanıksanmasına ve ihmal edilmesine yol açabilir (Ekstra.pdf Sayfa 37).
\end{testcard}

\begin{testcard}{Soru 24}
    Aşağıdakilerden hangisi İSG'nin proaktif yaklaşım ilkelerinden biridir?
    \begin{itemize}
        \item[A)] Kaza nedenini cezalandırma ile çözmek
        \item[B)] Tehlikeyi kaynağında önleyerek koruma sağlamak
        \item[C)] Sadece kazadan sonra rapor hazırlamak
        \item[D)] Sigorta primlerini artırmak
        \item[E)] Hiçbiri
    \end{itemize}
    \tcblower
    \textbf{Doğru Cevap: B} \\
    \textbf{Açıklama:} Proaktif İSG anlayışının özü, tehlikeleri kaynağından yok etmek veya azaltarak kazaları en baştan önlemektir.
\end{testcard}

\begin{testcard}{Soru 25}
    Aşağıdakilerden hangisi bir İSG uzmanının görev yetki ve yükümlülükleri hakkında yanlıştır?
    \begin{itemize}
        \item[A)] İşverene rehberlik ve danışmanlık yapmak
        \item[B)] Risk değerlendirmesi çalışmalarına katılmak
        \item[C)] İşyerindeki tüm üretimi tek başına yönetmek ve planlamak
        \item[D)] Çalışma ortamının gözetimini yapmak
        \item[E)] Hiçbiri
    \end{itemize}
    \tcblower
    \textbf{Doğru Cevap: C} \\
    \textbf{Açıklama:} İSG Uzmanı rehberlik, danışmanlık ve gözetim yapar; üretimi doğrudan yönetmek veya planlamak işverenin veya üretim müdürlerinin görevidir.
\end{testcard}

\begin{testcard}{Soru 26}
    İşyerinde İSG tedbirlerine uymakla yükümlü olan taraflar kimlerdir?
    \begin{itemize}
        \item[A)] Sadece işverenler
        \item[B)] Sadece işçiler
        \item[C)] Hem işveren hem de çalışanlar ortaklaşa sorumludur
        \item[D)] Sadece iş müfettişleri
        \item[E)] Hiçbiri
    \end{itemize}
    \tcblower
    \textbf{Doğru Cevap: C} \\
    \textbf{Açıklama:} İşveren tedbirleri almak ve denetlemekle; çalışanlar ise bu kurallara uymak ve tehlikeleri bildirmekle yükümlüdür.
\end{testcard}

\begin{testcard}{Soru 27}
    Aşağıdakilerden hangisi fiziksel risk etmenleri arasında değerlendirilir?
    \begin{itemize}
        \item[A)] Ağır metaller ve solvent gazlar
        \item[B)] Gürültü, titreşim ve termal konfor şartları
        \item[C)] Bakteri ve virüsler
        \item[D)] Mobbing ve stres
        \item[E)] Hiçbiri
    \end{itemize}
    \tcblower
    \textbf{Doğru Cevap: B} \\
    \textbf{Açıklama:} Gürültü, titreşim, aydınlatma, radyasyon, termal konfor (sıcaklık, nem) fiziksel risk etmenleridir.
\end{testcard}

\begin{testcard}{Soru 28}
    Kimyasal risk etmenlerine bağlı hastalıklara yol açan maddelerden biri aşağıdakilerden hangisidir?
    \begin{itemize}
        \item[A)] İyonize radyasyon
        \item[B)] Titreşim
        \item[C)] Pestisitler (tarım ilaçları)
        \item[D)] Düşük aydınlatma
        \item[E)] Hiçbiri
    \end{itemize}
    \tcblower
    \textbf{Doğru Cevap: C} \\
    \textbf{Açıklama:} Pestisitler, solventler, asitler ve zehirli gazlar kimyasal risk faktörleridir (Ekstra.pdf Sayfa 48).
\end{testcard}

\begin{testcard}{Soru 29}
    İşyerinde asbest söküm işi yapan bir çalışanın karşı karşıya olduğu en büyük risk grubu hangisidir?
    \begin{itemize}
        \item[A)] Ergonomik riskler
        \item[B)] Biyolojik riskler
        \item[C)] Kimyasal riskler (asbest lifleri ve tozları)
        \item[D)] Psikososyal riskler
        \item[E)] Hiçbiri
    \end{itemize}
    \tcblower
    \textbf{Doğru Cevap: C} \\
    \textbf{Açıklama:} Asbest lifleri solunduğunda akciğer kanseri ve mezotelyomaya yol açan son derece tehlikeli kimyasal ve toz risk etmenidir.
\end{testcard}

\begin{testcard}{Soru 30}
    Risk değerlendirmesi çalışmalarında geniş çalışan katılımı sağlanması hakkında hangisi doğrudur?
    \begin{itemize}
        \item[A)] Kanunen zorunlu olup risk algısını ve sahiplenmeyi artırır
        \item[B)] İşlerin yavaşlamasına neden olduğu için gereksizdir
        \item[C)] Sadece az tehlikeli sınıfta uygulanır
        \item[D)] Sadece sendikalı işyerlerinde zorunludur
        \item[E)] Hiçbiri
    \end{itemize}
    \tcblower
    \textbf{Doğru Cevap: A} \\
    \textbf{Açıklama:} Çalışan katılımı proaktif yaklaşımın temelidir. Çalışanların görüşlerinin alınması yasal bir zorunluluktur.
\end{testcard}

\begin{testcard}{Soru 31}
    Aşağıdakilerden hangisi bir iş güvenliği uzmanının işyerinde hayati tehlike tespit ettiğinde yapması gereken yasal eylemdir?
    \begin{itemize}
        \item[A)] Durumu görmezden gelmek
        \item[B)] Tehlikeyi işverene yazılı olarak bildirmek, acil önlem alınmazsa Bakanlığa rapor etmek
        \item[C)] İşyerini doğrudan satın almak
        \item[D)] Kendi kendine makineyi sökmek
        \item[E)] Hiçbiri
    \end{itemize}
    \tcblower
    \textbf{Doğru Cevap: B} \\
    \textbf{Açıklama:} İSG uzmanı acil durumları tespit ettiğinde işverene yazar; önlem alınmazsa yetkili mercilere (ÇSGB) bildirmekle yükümlüdür.
\end{testcard}

\begin{testcard}{Soru 32}
    İşyerinde ergonomik risklerin kontrol altına alınması için yapılan çalışmaların odağında ne yer almalıdır?
    \begin{itemize}
        \item[A)] İşçiyi en ağır koşullarda çalışmaya zorlamak
        \item[B)] İşin ve çalışma ortamının çalışana uyumlu hale getirilmesi
        \item[C)] Sadece çalışma saatlerinin artırılması
        \item[D)] Ürünlerin nakliye maliyetini düşürmek
        \item[E)] Hiçbiri
    \end{itemize}
    \tcblower
    \textbf{Doğru Cevap: B} \\
    \textbf{Açıklama:} Ergonomi, işin insana uydurulmasını ve fiziksel/zihinsel zorlanmaların azaltılmasını hedefler.
\end{testcard}

\begin{testcard}{Soru 33}
    İSG'ye Giriş ders notu Sayfa 17-23 kapsamına göre, aşağıdakilerden hangisi İSG'nin tarihsel gelişimindeki sanayi devrimi sonrasındaki dönüm noktalarından biridir?
    \begin{itemize}
        \item[A)] İlk iş müfettişlerinin görevlendirilmesi ve çocuk çalıştırma yasakları
        \item[B)] Sadece bilgisayar yazılımlarının geliştirilmesi
        \item[C)] Tarım devriminin başlaması
        \item[D)] İlk tekerleğin icat edilmesi
        \item[E)] Hiçbiri
    \end{itemize}
    \tcblower
    \textbf{Doğru Cevap: A} \\
    \textbf{Açıklama:} Sanayi devrimiyle birlikte artan kazalar sonucu İngiltere ve Avrupa'da fabrikalar kanunu çıkarılmış ve çocuk çalıştırma sınırları getirilmiştir.
\end{testcard}

\begin{testcard}{Soru 34}
    Bir kimyasal maddenin malzeme güvenlik bilgi formu (MSDS) ne işe yarar?
    \begin{itemize}
        \item[A)] Kimyasalın sadece fiyatını belirtir
        \item[B)] Kimyasalın tehlikelerini, depolama ve ilk yardım kurallarını içeren teknik dokümandır
        \item[C)] Şirketin yıllık vergi levhasıdır
        \item[D)] Çalışanın iş sözleşmesidir
        \item[E)] Hiçbiri
    \end{itemize}
    \tcblower
    \textbf{Doğru Cevap: B} \\
    \textbf{Açıklama:} MSDS, kimyasalların güvenli kullanımı, riskleri, kişisel korunma ve acil müdahale adımlarını içeren belgedir.
\end{testcard}

\begin{testcard}{Soru 35}
    Aşağıdakilerden hangisi biyolojik risk etmenlerine karşı koruma tedbirleri arasında yer alır?
    \begin{itemize}
        \item[A)] Gürültü ölçümü yapmak
        \item[B)] Aşılamalar, hijyen kuralları ve havalandırma
        \item[C)] Çalışanı ergonomik koltukta oturtmak
        \item[D)] Topraklama hattı çekmek
        \item[E)] Hiçbiri
    \end{itemize}
    \tcblower
    \textbf{Doğru Cevap: B} \\
    \textbf{Açıklama:} Biyolojik etkenlere karşı koruma hijyen, aşılama ve uygun kişisel koruyucu donanım (maske vb.) ile sağlanır.
\end{testcard}

\begin{testcard}{Soru 36}
    İşyerinde termal konfor şartları denilince hangisi akla gelmez?
    \begin{itemize}
        \item[A)] Hava sıcaklığı ve nem oranı
        \item[B)] Hava akım hızı
        \item[C)] Radyant sıcaklık
        \item[D)] Elektromanyetik dalga frekansı
        \item[E)] Hiçbiri
    \end{itemize}
    \tcblower
    \textbf{Doğru Cevap: D} \\
    \textbf{Açıklama:} Sıcaklık, nem, hava akımı ve radyant sıcaklık termal konforu oluşturur; elektromanyetik dalgalar ise radyasyon başlığındadır.
\end{testcard}

\begin{testcard}{Soru 37}
    İşveren, çalışanların işe girişlerinde ve iş değişikliğinde hangisini yaptırmak zorundadır?
    \begin{itemize}
        \item[A)] Sağlık raporu aldırmak
        \item[B)] Sadece sabıka kaydı istemek
        \item[C)] Ehliyet sınavına sokmak
        \item[D)] Sendika üyeliğini zorunlu kılmak
        \item[E)] Hiçbiri
    \end{itemize}
    \tcblower
    \textbf{Doğru Cevap: A} \\
    \textbf{Açıklama:} Çalışanların işe girişlerinde, iş değişikliğinde ve periyodik aralıklarla sağlık muayenelerinin yapılması zorunludur.
\end{testcard}

\begin{testcard}{Soru 38}
    İSG Kanunu Madde 4 kapsamında işverenin çalışanları bilgilendirme yükümlülüğü neleri kapsar?
    \begin{itemize}
        \item[A)] Sadece rakip şirketlerin finansal raporlarını
        \item[B)] İşyerindeki tehlikeler, riskler ve alınması gereken yasal tedbirleri
        \item[C)] Şirketin gizli patent kodlarını
        \item[D)] Müşteri telefon listelerini
        \item[E)] Hiçbiri
    \end{itemize}
    \tcblower
    \textbf{Doğru Cevap: B} \\
    \textbf{Açıklama:} İşveren, çalışanları maruz kalabilecekleri tehlike ve riskler ile koruyucu önlemler konusunda bilgilendirmelidir.
\end{testcard}

\begin{testcard}{Soru 39}
    İş sağlığı ve güvenliği profesyoneli olarak adlandırılan işyeri hekimleri ve iş güvenliği uzmanlarının bağımsızlığı hakkında hangisi doğrudur?
    \begin{itemize}
        \item[A)] Kararlarında ve mesleki uygulamalarında bağımsızdırlar, görevleri nedeniyle hakları kısıtlanamaz
        \item[B)] Tamamen işverenin her dediğini yapmak zorundadırlar
        \item[C)] Bağımsızlıkları kanunen tanınmamıştır
        \item[D)] Sadece devlet memuru olmak zorundadırlar
        \item[E)] Hiçbiri
    \end{itemize}
    \tcblower
    \textbf{Doğru Cevap: A} \\
    \textbf{Açıklama:} Yasa gereği hekim ve uzmanlar rehberlik görevlerini bağımsız ve etik ilkeler doğrultusunda yerine getirirler.
\end{testcard}

\begin{testcard}{Soru 40}
    Çalışan temsilcisi kimdir?
    \begin{itemize}
        \item[A)] İşverenin en yakın akrabası
        \item[B)] İş sağlığı ve güvenliği konularında çalışanlarla işveren arasında iletişimi sağlayan seçilmiş veya atanmış çalışan
        \item[C)] Şirketin avukatı
        \item[D)] Sendika genel başkanı
        \item[E)] Hiçbiri
    \end{itemize}
    \tcblower
    \textbf{Doğru Cevap: B} \\
    \textbf{Açıklama:} Çalışan temsilcisi, çalışanların İSG süreçlerine katılımını sağlayan ve kurullarda onları temsil eden kişidir.
\end{testcard}

\begin{testcard}{Soru 41}
    İSG mevzuatına göre, tehlike sınıfı tebliğinde yer almayan bir faaliyet için sınıflandırma nasıl belirlenir?
    \begin{itemize}
        \item[A)] İşverenin kendi seçimine göre
        \item[B)] Benzer işleri yapan diğer işyerlerinin sınıfına ve Bakanlık görüşüne göre
        \item[C)] Sınıflandırma dışı tutulur
        \item[D)] En yüksek riskli sınıfa otomatik atanır
        \item[E)] Hiçbiri
    \end{itemize}
    \tcblower
    \textbf{Doğru Cevap: B} \\
    \textbf{Açıklama:} Tehlike sınıfı tebliğlerinde bulunmayan faaliyetler için Bakanlık görüşü ve benzer iş kollarının NACE kodları esas alınır.
\end{testcard}

\begin{testcard}{Soru 42}
    Temel İSG eğitimlerinde çalışanlara ilk yardım, yangın ve tahliye konularında eğitim verilmesi yasal olarak zorunlu mudur?
    \begin{itemize}
        \item[A)] Hayır, sadece isteğe bağlıdır
        \item[B)] Evet, temel İSG eğitim konuları arasında yer almaktadır
        \item[C)] Sadece itfaiyeciler için zorunludur
        \item[D)] Sadece çok tehlikeli sınıfta zorunludur
        \item[E)] Hiçbiri
    \end{itemize}
    \tcblower
    \textbf{Doğru Cevap: B} \\
    \textbf{Açıklama:} Yangın, tahliye, ilkyardım ve acil durum planlamaları temel eğitim müfredatının zorunlu parçalarıdır.
\end{testcard}

\begin{testcard}{Soru 43}
    6331 sayılı İSG Kanunu'na göre, aşağıdakilerden hangisi kapsam dışıdır?
    \begin{itemize}
        \item[A)] Fabrika işçileri
        \item[B)] Ev hizmetlerinde çalışanlar ve kendi namına çalışan tek şahıslar
        \item[C)] Kamu kurumundaki memurlar
        \item[D)] Maden ocağı çalışanları
        \item[E)] Hiçbiri
    \end{itemize}
    \tcblower
    \textbf{Doğru Cevap: B} \\
    \textbf{Açıklama:} Kanun; ev hizmetleri, askeri ve emniyet faaliyetleri ile afet müdahale ekiplerini kapsam dışı tutmuştur.
\end{testcard}

\begin{testcard}{Soru 44}
    İşyerinde ramak kala olayların kaydedilmesi ve raporlanmasının temel amacı nedir?
    \begin{itemize}
        \item[A)] Şirketin prestijini sarsmak
        \item[B)] Çalışanları işten atmak
        \item[C)] Olası kazaların önceden tahmin edilerek gerekli tedbirlerin alınması
        \item[D)] SGK cezalarından kaçınmak
        \item[E)] Hiçbiri
    \end{itemize}
    \tcblower
    \textbf{Doğru Cevap: C} \\
    \textbf{Açıklama:} Ramak kala olaylar kazaların habercisidir. Kaydedilip önlem alınması büyük kazaları engeller.
\end{testcard}

\begin{testcard}{Soru 45}
    İşverenin gözetim yükümlülüğü ne demektir?
    \begin{itemize}
        \item[A)] İşçileri sürekli kameralarla izlemek
        \item[B)] İşyerinde alınan İSG tedbirlerinin uygulanıp uygulanmadığını izlemek ve denetlemek
        \item[C)] İşçinin ev hayatını takip etmek
        \item[D)] Sadece işe geliş gidiş saatlerini kaydetmek
        \item[E)] Hiçbiri
    \end{itemize}
    \tcblower
    \textbf{Doğru Cevap: B} \\
    \textbf{Açıklama:} İşveren sadece önlem almakla kalmayıp, bu önlemlere uyulup uyulmadığını gözetlemek ve denetlemekle de sorumludur.
\end{testcard}

\begin{testcard}{Soru 46}
    Risk kontrol hiyerarşisinde 'idari önlemler' kapsamına aşağıdakilerden hangisi girer?
    \begin{itemize}
        \item[A)] Tehlikeli makineyi kaldırmak
        \item[B)] İş rotasyonu, çalışma sürelerinin düzenlenmesi ve eğitimler
        \item[C)] Çelik burunlu ayakkabı vermek
        \item[D)] Koruyucu bariyer koymak
        \item[E)] Hiçbiri
    \end{itemize}
    \tcblower
    \textbf{Doğru Cevap: B} \\
    \textbf{Açıklama:} İş organizasyonu, rotasyon, vardiya düzenlemesi ve İSG eğitimleri idari önlemler basamağındadır.
\end{testcard}

\begin{testcard}{Soru 47}
    İşyeri İSG organizasyonlarında OSGB ne anlama gelmektedir?
    \begin{itemize}
        \item[A)] Ortak Sağlık ve Güvenlik Birimi
        \item[B)] Özel Sektör Geliştirme Birliği
        \item[C)] Ortak Sanayi Gözetim Bürosu
        \item[D)] Ölçüm ve Sınıflandırma Genel Başkanlığı
        \item[E)] Hiçbiri
    \end{itemize}
    \tcblower
    \textbf{Doğru Cevap: A} \\
    \textbf{Açıklama:} OSGB, işyerlerine İSG hizmetlerini (hekim, uzman, personel) sunmak üzere Bakanlıkça yetkilendirilmiş Ortak Sağlık ve Güvenlik Birimleri'dir.
\end{testcard}

\begin{testcard}{Soru 48}
    Çalışanların İSG konusundaki temel yükümlülükleri arasında aşağıdakilerden hangisi yer almaz?
    \begin{itemize}
        \item[A)] Makine ve ekipmanları kurallara uygun kullanmak
        \item[B)] Kendilerine verilen kişisel koruyucu donanımları doğru kullanmak ve korumak
        \item[C)] İşyerindeki tehlikeleri ve eksiklikleri derhal işverene bildirmek
        \item[D)] Diğer işyerlerinin risk analizlerini hazırlamak
        \item[E)] Hiçbiri
    \end{itemize}
    \tcblower
    \textbf{Doğru Cevap: D} \\
    \textbf{Açıklama:} Çalışanlar kendi işyerlerindeki kurallara uymak ve bildirim yapmakla yükümlüdür; başka işyerlerinin analizini yapmak görevleri değildir.
\end{testcard}

\begin{testcard}{Soru 49}
    İş kazası teorisinde 'güvenli davranış kültürü' oluşturmanın önündeki en büyük engel hangisidir?
    \begin{itemize}
        \item[A)] Eğitim süresinin kısalığı
        \item[B)] Bana bir şey olmaz algısı ve güvensiz hareketlerin kanıksanması
        \item[C)] Koruyucu ekipmanların renkleri
        \item[D)] Fabrikadaki aydınlatma yetersizliği
        \item[E)] Hiçbiri
    \end{itemize}
    \tcblower
    \textbf{Doğru Cevap: B} \\
    \textbf{Açıklama:} Güvenlik kültürünün en büyük engeli bireysel kanıksama, 'bana bir şey olmaz' inancı ve kuralların ihlal edilmesinin alışkanlık haline gelmesidir.
\end{testcard}

\begin{testcard}{Soru 50}
    İSG'ye Giriş ders notu sayfa 29'da belirtilen risk tanımının en kritik boyutu hangisidir?
    \begin{itemize}
        \item[A)] Risk sadece maddi kayıptır
        \item[B)] Risk bir olasılık ve bu olasılığın yaratacağı şiddetin çarpımı / bileşkesidir
        \item[C)] Risk sadece işçinin yaşına bağlıdır
        \item[D)] Risk sadece doğal afetlerdir
        \item[E)] Hiçbiri
    \end{itemize}
    \tcblower
    \textbf{Doğru Cevap: B} \\
    \textbf{Açıklama:} Risk iki boyutludur: Olayın meydana gelme olasılığı (sıklığı) ile yol açacağı hasarın şiddeti (ciddiyeti).
\end{testcard}

\begin{testcard}{Soru 51}
    İşyerlerinde hazırlanan 'Acil Durum Planı' kavramının kapsamı, yasal amacı ve niteliği dikkate alındığında aşağıdakilerden hangisi yanlıştır?
    \begin{itemize}
        \item[A)] Acil durum planı; yangın, deprem, sızıntı gibi olaylarda kimin, ne zaman, nasıl davranacağını belirleyen dinamik ve uygulamalı bir plandır.
        \item[B)] Acil durum planı, sadece iş müfettişlerinin denetimlerinden geçmek veya ceza almamak amacıyla hazırlanıp çekmecede saklanan statik bir belgedir.
        \item[C)] Acil durum planı kapsamında söndürme, kurtarma, koruma ve ilk yardım gibi özel ekipler kurulmalı ve düzenli tatbikatlar yapılmalıdır.
        \item[D)] Acil durum planı, işyerindeki tüm çalışanların katılımını gerektiren ve sürekli güncel tutulması gereken bir süreçtir.
        \item[E)] Hiçbiri
    \end{itemize}
    \tcblower
    \textbf{Doğru Cevap: B} \\
    \textbf{Açıklama:} Acil durum planı kesinlikle sadece denetim için saklanan durağan bir kağıt parçası değildir; tatbikatlarla desteklenen, güncel tutulan ve hayat kurtarmayı amaçlayan canlı bir plandır.
\end{testcard}

\begin{testcard}{Soru 52}
    Bir işyerinin NACE kodunun belirlenmesi ve buna bağlı olarak tehlike sınıfının (Az Tehlikeli, Tehlikeli, Çok Tehlikeli) tespiti, iş sağlığı ve güvenliği uygulamalarındaki hangi yasal süre veya yükümlülük üzerinde doğrudan belirleyici bir etkiye sahip değildir?
    \begin{itemize}
        \item[A)] Çalışanların temel İSG eğitimlerinin saat süreleri ve yenilenme periyotları
        \item[B)] Risk değerlendirmesinin ve acil durum planlarının yenilenme periyotları
        \item[C)] Çalışanların kıdem tazminatlarının ödeme günü ve yıllık izin haklarının hesaplanması
        \item[D)] İşyeri hekimi ve iş güvenliği uzmanının çalışan başına aylık asgari çalışma süreleri
        \item[E)] Hiçbiri
    \end{itemize}
    \tcblower
    \textbf{Doğru Cevap: C} \\
    \textbf{Açıklama:} Tehlike sınıfı/NACE kodu; İSG eğitim sürelerini, risk analizi yenileme sürelerini (2-4-6 yıl), uzman ve hekim çalışma sürelerini belirler; ancak kıdem tazminatı veya yıllık izin haklarının hesaplanması iş hukuku (4857 sayılı kanun) kapsamında çalışanın kıdem yılına göre belirlenir, tehlike sınıfıyla ilgisi yoktur.
\end{testcard}

\subsection{Ünite 2: Risk Değerlendirme Metodolojileri}
\begin{testcard}{Soru 1}
    FMEA (Hata Türleri ve Etkileri Analizi) metodolojisinde Risk Öncelik Sayısı (RÖS) hesaplanırken hangi parametreler çarpılır?
    \begin{itemize}
        \item[A)] Olasılık × Şiddet × Frekans
        \item[B)] Olasılık × Şiddet × Saptanabilirlik
        \item[C)] İhtimal × Frekans × Şiddet
        \item[D)] Olasılık × Şiddet × Etki
        \item[E)] Hiçbiri
    \end{itemize}
    \tcblower
    \textbf{Doğru Cevap: B} \\
    \textbf{Açıklama:} FMEA metodolojisinde RÖS = Olasılık (Occurrence) x Şiddet (Severity) x Saptanabilirlik (Detection) formülü kullanılır.
\end{testcard}

\begin{testcard}{Soru 2}
    FMEA çalışmasında, bir hatayı saptamak ve fark etmek çok kolaysa saptanabilirlik (D) puanı nasıl belirlenir?
    \begin{itemize}
        \item[A)] Düşük puan verilir (örneğin 1 veya 2)
        \item[B)] En yüksek puan verilir (10)
        \item[C)] Puanlama dışı tutulur
        \item[D)] Olasılık puanıyla eşitlenir
        \item[E)] Hiçbiri
    \end{itemize}
    \tcblower
    \textbf{Doğru Cevap: A} \\
    \textbf{Açıklama:} FMEA'da hatayı saptamak kolaysa risk düşüktür, bu yüzden saptanabilirlik (Detection) derecesi düşük puanlanır.
\end{testcard}

\begin{testcard}{Soru 3}
    Fine-Kinney risk değerlendirme metodunun klasik risk matrisinden (5x5) en temel farkı nedir?
    \begin{itemize}
        \item[A)] Sadece kimya tesislerinde kullanılması
        \item[B)] Risk hesabına 'Frekans' (tehlikeye maruz kalma sıklığı) parametresini eklemesi
        \item[C)] Nitel bir yöntem olması
        \item[D)] Sadece ölümlü kazaları hesaplaması
        \item[E)] Hiçbiri
    \end{itemize}
    \tcblower
    \textbf{Doğru Cevap: B} \\
    \textbf{Açıklama:} Fine-Kinney metodunda Risk = Olasılık x Frekans x Şiddet formülü kullanılır ve maruziyet sıklığı (Frekans) hesaba dahil edilir.
\end{testcard}

\begin{testcard}{Soru 4}
    HAZOP (Tehlike ve İşletilebilirlik Analizi) metodolojisinde kullanılan 'Fazla', 'Eksik', 'Ters' gibi kelimelere ne ad verilir?
    \begin{itemize}
        \item[A)] Sapma katsayıları
        \item[B)] Kılavuz kelimeler (Guide Words)
        \item[C)] Risk puanları
        \item[D)] Şiddet seviyeleri
        \item[E)] Hiçbiri
    \end{itemize}
    \tcblower
    \textbf{Doğru Cevap: B} \\
    \textbf{Açıklama:} HAZOP analizinde boru hatlarındaki akış, basınç, sıcaklık gibi sapmaları belirlemek için kılavuz kelimeler kullanılır.
\end{testcard}

\begin{testcard}{Soru 5}
    Hata Ağacı Analizi (FTA) metodolojisiyle ilgili aşağıdakilerden hangisi yanlıştır?
    \begin{itemize}
        \item[A)] Yukarıdan aşağıya (Top-down) çalışır
        \item[B)] Mantıksal VE/VEYA kapılarını kullanır
        \item[C)] İleriye doğru çalışan tümevarımcı (Bottom-up) bir yöntemdir
        \item[D)] İstenmeyen bir tepe olaydan kök nedenlere iner
        \item[E)] Hiçbiri
    \end{itemize}
    \tcblower
    \textbf{Doğru Cevap: C} \\
    \textbf{Açıklama:} FTA tümdengelimci (yukarıdan aşağıya) çalışır. Aşağıdan yukarıya / ileriye doğru çalışan tümevarımcı yöntem Olay Ağacı Analizidir (ETA).
\end{testcard}

\begin{testcard}{Soru 6}
    Olay Ağacı Analizi (ETA) hakkında hangisi doğrudur?
    \begin{itemize}
        \item[A)] Başlatıcı olayla başlayıp ileriye doğru güvenlik bariyerlerinin başarısını sorgulayan tümevarımcı bir metottur
        \item[B)] Sadece VE mantıksal kapısını kullanır
        \item[C)] Geriye doğru kök neden arar
        \item[D)] Tasarım aşamasında asla kullanılmaz
        \item[E)] Hiçbiri
    \end{itemize}
    \tcblower
    \textbf{Doğru Cevap: A} \\
    \textbf{Açıklama:} ETA, bir başlangıç olayından yola çıkarak güvenlik bariyerlerinin çalışıp çalışmamasına göre sistemin ulaşacağı sonuçları tümevarımla belirler.
\end{testcard}

\begin{testcard}{Soru 7}
    Ergonomi biliminin iş sağlığı ve güvenliğindeki temel amacı nedir?
    \begin{itemize}
        \item[A)] İşçiyi yüksek hızda çalışmaya zorlamak
        \item[B)] İşyerinin estetik görünümünü düzenlemek
        \item[C)] İş ve çalışma ortamının insanın fiziksel ve zihinsel kapasitesine uygun hale getirilmesi
        \item[D)] Sadece oturarak çalışma düzeni oluşturmak
        \item[E)] Hiçbiri
    \end{itemize}
    \tcblower
    \textbf{Doğru Cevap: C} \\
    \textbf{Açıklama:} Ergonominin amacı çalışma ortamını, araçları ve işi insanın kapasitesine ve sınırlarına uydurmaktır.
\end{testcard}

\begin{testcard}{Soru 8}
    Aşağıdakilerden hangisi fiziksel ergonominin inceleme konuları arasında yer alır?
    \begin{itemize}
        \item[A)] Çalışma duruşları, tekrarlı hareketler, malzeme taşıma ve yerleşim
        \item[B)] Karar verme ve zihinsel iş yükü
        \item[C)] Şirket içi iletişim ve takım çalışması
        \item[D)] Şirketin yıllık bütçe planlaması
        \item[E)] Hiçbiri
    \end{itemize}
    \tcblower
    \textbf{Doğru Cevap: A} \\
    \textbf{Açıklama:} Duruşlar, el hareketleri, yük kaldırma ve antropometri fiziksel ergonominin temel inceleme alanlarıdır.
\end{testcard}

\begin{testcard}{Soru 9}
    Zihinsel iş yükü, stres, insan-bilgisayar etkileşimi ve karar verme süreçleri ergonominin hangi alt dalına girer?
    \begin{itemize}
        \item[A)] Fiziksel Ergonomi
        \item[B)] Bilişsel (Kognitif) Ergonomi
        \item[C)] Örgütsel Ergonomi
        \item[D)] Çevresel Ergonomi
        \item[E)] Hiçbiri
    \end{itemize}
    \tcblower
    \textbf{Doğru Cevap: B} \\
    \textbf{Açıklama:} Bilişsel ergonomi, algılama, hafıza, muhakeme ve motor tepkiler gibi zihinsel süreçlerle ilgilenir.
\end{testcard}

\begin{testcard}{Soru 10}
    Hata Ağacı Analizinde (FTA) kullanılan 'VE' kapısı (AND gate) neyi temsil eder?
    \begin{itemize}
        \item[A)] Giriş olaylarından sadece birinin gerçekleşmesi durumunda çıkış olayının gerçekleşeceğini
        \item[B)] Çıkış olayının gerçekleşmesi için giriş olaylarının tamamının aynı anda gerçekleşmesi gerektiğini
        \item[C)] Kaza olasılığının sıfır olduğunu
        \item[D)] Sürecin durdurulduğunu
        \item[E)] Hiçbiri
    \end{itemize}
    \tcblower
    \textbf{Doğru Cevap: B} \\
    \textbf{Açıklama:} VE kapısı, bağlı olan tüm alt olayların/nedenlerin aynı anda gerçekleşmesi halinde tepe olayın oluşacağını belirtir.
\end{testcard}

\begin{testcard}{Soru 11}
    Hata Ağacı Analizinde (FTA) kullanılan 'VEYA' kapısı (OR gate) neyi temsil eder?
    \begin{itemize}
        \item[A)] Çıkış olayının gerçekleşmesi için giriş olaylarından en az birinin gerçekleşmesinin yeterli olduğunu
        \item[B)] Tüm olayların aynı anda olmasının zorunlu olduğunu
        \item[C)] Olayların birbirini engellediğini
        \item[D)] Hiçbir olayın gerçekleşmediğini
        \item[E)] Hiçbiri
    \end{itemize}
    \tcblower
    \textbf{Doğru Cevap: A} \\
    \textbf{Açıklama:} VEYA kapısı, alt olaylardan herhangi birinin gerçekleşmesinin tepe olayı tetiklemek için yeterli olduğunu gösterir.
\end{testcard}

\begin{testcard}{Soru 12}
    Kontrol Listesi (Check-List) analiz metodu hakkında aşağıdakilerden hangisi doğrudur?
    \begin{itemize}
        \item[A)] Son derece karmaşık matematiksel olasılıklar içerir
        \item[B)] İşyerindeki eksikliklerin önceden hazırlanmış sorular/maddeler eşliğinde evet/hayır şeklinde denetlenmesidir
        \item[C)] Sadece nükleer santrallerde kullanılır
        \item[D)] Bir beyin fırtınası yöntemidir
        \item[E)] Hiçbiri
    \end{itemize}
    \tcblower
    \textbf{Doğru Cevap: B} \\
    \textbf{Açıklama:} Check-list yöntemi, önceden hazırlanmış kontrol listeleriyle mevcut durumun standartlara uygunluğunu hızlıca tespit etmeye yarar.
\end{testcard}

\begin{testcard}{Soru 13}
    İşyerinde el aletlerinin (örneğin kerpeten, tornavida) ergonomik olarak tasarlanmasında hangisi göz önünde bulundurulmalıdır?
    \begin{itemize}
        \item[A)] Sadece en ucuz plastik malzemenin kullanılması
        \item[B)] El anatomisi, kavrama kuvveti dağılımı ve kaymaz sap yüzeyi
        \item[C)] Aletin olabildiğince ağır yapılması
        \item[D)] Aletin renginin mavi olması
        \item[E)] Hiçbiri
    \end{itemize}
    \tcblower
    \textbf{Doğru Cevap: B} \\
    \textbf{Açıklama:} Ergonomik el aletleri elin anatomik yapısına uymalı, el bileğinde aşırı bükülmelere yol açmamalıdır.
\end{testcard}

\begin{testcard}{Soru 14}
    Aşağıdakilerden hangisi antropometrinin tanımıdır?
    \begin{itemize}
        \item[A)] İnsan vücut ölçülerinin sistematik olarak ölçülüp sınıflandırılması
        \item[B)] Çalışanların haftalık çalışma saatlerinin takibi
        \item[C)] İş kazası istatistiklerinin toplanması
        \item[D)] Kimyasal maddelerin analiz edilmesi
        \item[E)] Hiçbiri
    \end{itemize}
    \tcblower
    \textbf{Doğru Cevap: A} \\
    \textbf{Açıklama:} Antropometri, insan vücudunun boyutlarını, erişim mesafelerini ve fiziksel sınırlarını ölçen bilim dalıdır.
\end{testcard}

\begin{testcard}{Soru 15}
    Ergonomik koltuk tasarımında aşağıdaki özelliklerden hangisi aranmaz?
    \begin{itemize}
        \item[A)] Yükseklik ve sırt eğiminin ayarlanabilir olması
        \item[B)] Bel desteğinin bulunması
        \item[C)] Tamamen sert metalden yapılmış ve ayarlanamaz olması
        \item[D)] Oturma yüzeyinin hava alabilen malzemeden olması
        \item[E)] Hiçbiri
    \end{itemize}
    \tcblower
    \textbf{Doğru Cevap: C} \\
    \textbf{Açıklama:} Ergonomik koltuklar kullanıcının boyuna ve oturma pozisyonuna göre ayarlanabilir esneklikte olmalıdır.
\end{testcard}

\begin{testcard}{Soru 16}
    HAZOP analizinde 'Sıcaklık' parametresi için 'Yüksek' kılavuz kelimesi birleştirildiğinde oluşan sapma neyi ifade eder?
    \begin{itemize}
        \item[A)] Süreçte aşırı ısınma riskini ve borularda basınç artışını
        \item[B)] Akışın tamamen durmasını
        \item[C)] Borularda donma olayını
        \item[D)] Sıvının renginin değişmesini
        \item[E)] Hiçbiri
    \end{itemize}
    \tcblower
    \textbf{Doğru Cevap: A} \\
    \textbf{Açıklama:} Parametre + kılavuz kelime sapmayı verir. Sıcaklık + Yüksek = Aşırı ısınma/reaksiyon tehlikesi.
\end{testcard}

\begin{testcard}{Soru 17}
    FMEA yönteminde eğer bir arıza türünün müşteriye ulaşmadan fabrikada yakalanma ihtimali çok zorsa (yani saptanması neredeyse imkansızsa) 'Saptanabilirlik (D)' puanı kaç olmalıdır?
    \begin{itemize}
        \item[A)] En düşük puan (1)
        \item[B)] Orta seviye (5)
        \item[C)] En yüksek puan (10)
        \item[D)] Sıfır puan
        \item[E)] Hiçbiri
    \end{itemize}
    \tcblower
    \textbf{Doğru Cevap: C} \\
    \textbf{Açıklama:} Saptanabilirlik zorlaştıkça risk artar. Fark edilmesi en zor hatalar için D katsayısı 10 puan olarak belirlenir.
\end{testcard}

\begin{testcard}{Soru 18}
    Bir işyerinde tekrarlı hareketlerin (örneğin sürekli paketleme veya montaj hattı işi) yol açabileceği en yaygın ergonomik sağlık sorunu hangisidir?
    \begin{itemize}
        \item[A)] İşitme kaybı
        \item[B)] Kas-iskelet sistemi hastalıkları (örneğin karpal tünel sendromu)
        \item[C)] Akciğer kanseri
        \item[D)] Bulaşıcı hastalıklar
        \item[E)] Hiçbiri
    \end{itemize}
    \tcblower
    \textbf{Doğru Cevap: B} \\
    \textbf{Açıklama:} Tekrarlı hareketler ve aşırı zorlanmalar kas, tendon ve sinir sıkışmalarına yol açarak kas-iskelet rahatsızlıkları oluşturur.
\end{testcard}

\begin{testcard}{Soru 19}
    Fine-Kinney metodunda risk skoru '200' ile '400' arasında çıkan bir durum için hangi seviyede eylem planlanmalıdır?
    \begin{itemize}
        \item[A)] Kabul edilebilir risk, önlem gerekmez
        \item[B)] Önemli risk, kısa vadede iyileştirme yapılmalıdır
        \item[C)] İhmal edilebilir risk
        \item[D)] İşyeri kapatılmalıdır
        \item[E)] Hiçbiri
    \end{itemize}
    \tcblower
    \textbf{Doğru Cevap: B} \\
    \textbf{Açıklama:} Fine-Kinney risk baremine göre 200-400 aralığı 'Önemli Risk' sınıfındadır ve kısa-orta vadede düzeltici eylem gerektirir.
\end{testcard}

\begin{testcard}{Soru 20}
    Aşağıdakilerden hangisi risk değerlendirme ekibinde yer alması yasal olarak zorunlu kişilerden biri değildir?
    \begin{itemize}
        \item[A)] İşveren veya işveren vekili
        \item[B)] İş güvenliği uzmanı ve işyeri hekimi
        \item[C)] Çalışan temsilcileri ve destek elemanları
        \item[D)] Şirketin tüm müşterileri ve tedarikçileri
        \item[E)] Hiçbiri
    \end{itemize}
    \tcblower
    \textbf{Doğru Cevap: D} \\
    \textbf{Açıklama:} Müşteriler veya tedarikçiler risk değerlendirme ekibinde yer almaz. Ekip işyeri yöneticileri, İSG profesyonelleri ve çalışan temsilcilerinden oluşur.
\end{testcard}

\begin{testcard}{Soru 21}
    İşyerinde yapılacak bir risk değerlendirmesinde, tehlikelerin belirlenmesi aşamasında hangisinden yararlanılmaz?
    \begin{itemize}
        \item[A)] Ramak kala raporları ve kaza istatistikleri
        \item[B)] Çalışanların şikayetleri ve sağlık gözetim sonuçları
        \item[C)] Kimyasal MSDS formları ve ortam ölçüm raporları
        \item[D)] Şirketin hisse senedi değer grafiklerinden
        \item[E)] Hiçbiri
    \end{itemize}
    \tcblower
    \textbf{Doğru Cevap: D} \\
    \textbf{Açıklama:} Hisse senedi verileri İSG tehlikelerini belirlemede kullanılacak teknik bir veri değildir.
\end{testcard}

\begin{testcard}{Soru 22}
    FMEA analizinin çeşitleri arasında aşağıdakilerden hangisi yer almaz?
    \begin{itemize}
        \item[A)] Tasarım FMEA (DFMEA)
        \item[B)] Proses FMEA (PFMEA)
        \item[C)] Finansal Yatırım FMEA (FFMEA)
        \item[D)] Servis FMEA (SFMEA)
        \item[E)] Hiçbiri
    \end{itemize}
    \tcblower
    \textbf{Doğru Cevap: C} \\
    \textbf{Açıklama:} FMEA tasarım, proses, servis ve sistem alanlarında uygulanır; finansal yatırım FMEA standart bir İSG aracı değildir.
\end{testcard}

\begin{testcard}{Soru 23}
    HAZOP analizlerinin genellikle çok disiplinli bir uzman grubu tarafından beyin fırtınası şeklinde yapılmasının temel amacı nedir?
    \begin{itemize}
        \item[A)] Toplantıyı uzatmak
        \item[B)] Farklı mühendislik ve işletme disiplinlerinin deneyimlerini birleştirerek karmaşık proses risklerini gözden kaçırmamak
        \item[C)] Ceza tutanaklarını hazırlamak
        \item[D)] Çalışanların maaşlarını belirlemek
        \item[E)] Hiçbiri
    \end{itemize}
    \tcblower
    \textbf{Doğru Cevap: B} \\
    \textbf{Açıklama:} Proses süreçleri elektrik, kimya, makine gibi birçok uzmanlık gerektirdiğinden HAZOP takım çalışmasıyla yürütülür.
\end{testcard}

\begin{testcard}{Soru 24}
    Fine-Kinney metodunda olasılık (O) değeri '10' neyi ifade eder?
    \begin{itemize}
        \item[A)] Kazanın neredeyse kesin olarak gerçekleşmesini (beklenen bir durum)
        \item[B)] Kazanın gerçekleşmesinin imkansız olduğunu
        \item[C)] Yılda bir kez olabileceğini
        \item[D)] Şiddetin en hafif seviyede olduğunu
        \item[E)] Hiçbiri
    \end{itemize}
    \tcblower
    \textbf{Doğru Cevap: A} \\
    \textbf{Açıklama:} Fine-Kinney olasılık skalasında 10 en yüksek değerdir ve olayın gerçekleşmesinin neredeyse kesin olduğunu ifade eder.
\end{testcard}

\begin{testcard}{Soru 25}
    Fine-Kinney metodunda şiddet (S) değeri '100' neyi ifade eder?
    \begin{itemize}
        \item[A)] Çoklu ölümlere yol açabilecek felaket durumunu
        \item[B)] İlkyardım gerektiren hafif sıyrığı
        \item[C)] Maddi hasarın sıfır olduğunu
        \item[D)] Sadece gürültü maruziyetini
        \item[E)] Hiçbiri
    \end{itemize}
    \tcblower
    \textbf{Doğru Cevap: A} \\
    \textbf{Açıklama:} Fine-Kinney şiddet tablosunda 100 felaket (birden fazla ölüm) durumuna karşılık gelen en yüksek şiddet derecesidir.
\end{testcard}

\begin{testcard}{Soru 26}
    FMEA çalışmasında RÖS (Risk Öncelik Sayısı) değeri kaçın üzerine çıktığında düzeltici faaliyetlerin yapılması zorunlu kabul edilir?
    \begin{itemize}
        \item[A)] RÖS > 10
        \item[B)] RÖS > 50
        \item[C)] RÖS > 100 (Genel kabul gören sınır)
        \item[D)] RÖS > 500
        \item[E)] Hiçbiri
    \end{itemize}
    \tcblower
    \textbf{Doğru Cevap: C} \\
    \textbf{Açıklama:} Klasik FMEA kabullerinde RÖS değeri 100'ü aştığında risk kabul edilemez/kritik görülür ve düzeltici önlem alınması zorunludur.
\end{testcard}

\begin{testcard}{Soru 27}
    Bir risk değerlendirmesinde risklerin analiz edilmesi ve derecelendirilmesinin temel amacı nedir?
    \begin{itemize}
        \item[A)] Önlemlerin aciliyet sırasını belirlemek (önceliklendirme)
        \item[B)] Müfettişlere sahte belgeler sunmak
        \item[C)] Şirketin yıllık bütçesini tüketmek
        \item[D)] Sorumluluğu tamamen çalışanlara yıkmak
        \item[E)] Hiçbiri
    \end{itemize}
    \tcblower
    \textbf{Doğru Cevap: A} \\
    \textbf{Açıklama:} Derecelendirme sayesinde en yüksek riskli alanlardan başlanarak bütçe ve kaynaklar doğru yönetilir.
\end{testcard}

\begin{testcard}{Soru 28}
    Ofis çalışmalarında bilgisayar ekranının üst kenarı çalışanın göz hizasıyla nasıl ayarlanmalıdır?
    \begin{itemize}
        \item[A)] Göz hizasından çok yukarıda olmalı
        \item[B)] Göz hizasıyla aynı yükseklikte veya hafifçe aşağıda olmalı
        \item[C)] Ekran yere dik olarak tamamen yatırılmalı
        \item[D)] Göz hizasından 50 cm aşağıda olmalı
        \item[E)] Hiçbiri
    \end{itemize}
    \tcblower
    \textbf{Doğru Cevap: B} \\
    \textbf{Açıklama:} Boyun zorlanmalarını önlemek için ekranın üst çizgisi göz seviyesinde veya hafifçe altında olmalıdır.
\end{testcard}

\begin{testcard}{Soru 29}
    Sürekli ayakta durarak çalışan bir işçiye ergonomik destek olarak aşağıdakilerden hangisi sunulabilir?
    \begin{itemize}
        \item[A)] Sert beton zemin
        \item[B)] Yorgunluk önleyici paspas ve uygun destek tabanlı ayakkabı
        \item[C)] Vardiyayı 16 saate çıkarmak
        \item[D)] Ağır yük kaldırma görevi eklemek
        \item[E)] Hiçbiri
    \end{itemize}
    \tcblower
    \textbf{Doğru Cevap: B} \\
    \textbf{Açıklama:} Yorgunluk önleyici paspaslar ayaktaki statik yükü azaltır ve dolaşım sistemini korur.
\end{testcard}

\begin{testcard}{Soru 30}
    Hata Ağacı Analizindeki (FTA) temel olay (primary event) neyi temsil eder?
    \begin{itemize}
        \item[A)] Kendi başına daha alt nedenlere bölünemeyen en temel kök arızayı
        \item[B)] Kazanın sonrasındaki felaket senaryosunu
        \item[C)] Alınan ilk yardımı
        \item[D)] Yangın söndürme tüpünü
        \item[E)] Hiçbiri
    \end{itemize}
    \tcblower
    \textbf{Doğru Cevap: A} \\
    \textbf{Açıklama:} FTA ağacının en altında daire ile gösterilen temel olaylar, analizin durdurulduğu en kök hata veya arıza durumlarıdır.
\end{testcard}

\begin{testcard}{Soru 31}
    Aşağıdakilerden hangisi proaktif risk değerlendirme metotlarından biri değildir?
    \begin{itemize}
        \item[A)] FMEA
        \item[B)] HAZOP
        \item[C)] Kaza sonrası adli inceleme ve raporlama
        \item[D)] Check-List
        \item[E)] Hiçbiri
    \end{itemize}
    \tcblower
    \textbf{Doğru Cevap: C} \\
    \textbf{Açıklama:} Kaza sonrası inceleme reaktif bir faaliyettir. Diğerleri ise kaza oluşmadan yapılan proaktif analizlerdir.
\end{testcard}

\begin{testcard}{Soru 32}
    FMEA'da RÖS puanını düşürmek için öncelikle hangi parametreye müdahale edilmelidir?
    \begin{itemize}
        \item[A)] Şiddet veya olasılığı düşürecek mühendislik kontrolleri ile saptanabilirliği artıracak önlemler
        \item[B)] Sadece çalışan sayısını azaltmak
        \item[C)] Şirketin adını değiştirmek
        \item[D)] İşyeri hekimini değiştirmek
        \item[E)] Hiçbiri
    \end{itemize}
    \tcblower
    \textbf{Doğru Cevap: A} \\
    \textbf{Açıklama:} Tasarım değişiklikleri veya koruyucular ile şiddet ve olasılık düşürülür; sensör ve denetimlerle saptanabilirlik kolaylaştırılır.
\end{testcard}

\begin{testcard}{Soru 33}
    İşyerinde elle yük taşıma işlerinde bel yaralanmalarını önlemek için asgari kural hangisidir?
    \begin{itemize}
        \item[A)] Yükü bel bükülerek ve dizler düz tutularak kaldırmak
        \item[B)] Yükü dizler bükülerek (çökülerek) ve gövdeye yakın tutarak kaldırmak
        \item[C)] Yükü tek elle ve koşarak taşımak
        \item[D)] Yükü başın üstünde taşımak
        \item[E)] Hiçbiri
    \end{itemize}
    \tcblower
    \textbf{Doğru Cevap: B} \\
    \textbf{Açıklama:} Elle kaldırmada yük vücuda yakın olmalı ve bacak kaslarından güç alınarak sırt düz tutulmalıdır.
\end{testcard}

\begin{testcard}{Soru 34}
    Bilişsel ergonomi çalışmalarının odağında aşağıdakilerden hangisi yer alır?
    \begin{itemize}
        \item[A)] Kontrol panellerinin ve göstergelerin kolay okunabilir ve anlaşılır tasarlanması
        \item[B)] Sandalye minderinin yumuşaklığı
        \item[C)] Masanın yerden yüksekliği
        \item[D)] Fabrikanın dış cephe rengi
        \item[E)] Hiçbiri
    \end{itemize}
    \tcblower
    \textbf{Doğru Cevap: A} \\
    \textbf{Açıklama:} Göstergelerin ve ekran arayüzlerinin anlaşılır olması zihinsel yükü azaltarak insan hatalarını önler.
\end{testcard}

\begin{testcard}{Soru 35}
    İSG risk değerlendirmelerinde riskin tamamen yok edilmesinin (eliminasyon) mümkün olmadığı durumlarda ne yapılmalıdır?
    \begin{itemize}
        \item[A)] Çalışmayı tamamen durdurup fabrikayı kapatmak
        \item[B)] Riski kabul edilebilir seviyeye indirecek kontrol hiyerarşisi adımlarını sırasıyla uygulamak
        \item[C)] Hiçbir önlem almadan çalışmaya devam etmek
        \item[D)] Sorumluluğu devlet sigortasına devretmek
        \item[E)] Hiçbiri
    \end{itemize}
    \tcblower
    \textbf{Doğru Cevap: B} \\
    \textbf{Açıklama:} Tehlike sıfırlanamıyorsa ikame, mühendislik kontrolleri, idari önlemler ve en son KKD ile kabul edilebilir sınıra çekilmelidir.
\end{testcard}

\begin{testcard}{Soru 36}
    Risk Değerlendirmesi Yönetmeliği'ne göre, hazırlanan risk değerlendirmesi dokümanında hangisinin bulunması zorunlu değildir?
    \begin{itemize}
        \item[A)] Belirlenen tehlikeler ve risk analiz sonuçları
        \item[B)] Karar verilen kontrol tedbirleri ve uygulama adımları
        \item[C)] Risk değerlendirmesini yapan ekibin imzaları
        \item[D)] Çalışanların kişisel banka hesap şifreleri
        \item[E)] Hiçbiri
    \end{itemize}
    \tcblower
    \textbf{Doğru Cevap: D} \\
    \textbf{Açıklama:} Kişisel finansal bilgiler İSG risk değerlendirme dokümanında yer almaz.
\end{testcard}

\begin{testcard}{Soru 37}
    ETA (Olay Ağacı Analizi) dallanmalarında genellikle hangi iki durum sorgulanır?
    \begin{itemize}
        \item[A)] Başarı / Başarısızlık (Evet / Hayır)
        \item[B)] Sıcak / Soğuk
        \item[C)] Pahalı / Ucuz
        \item[D)] Hızlı / Yavaş
        \item[E)] Hiçbiri
    \end{itemize}
    \tcblower
    \textbf{Doğru Cevap: A} \\
    \textbf{Açıklama:} ETA dallarında güvenlik bariyerinin çalışıp çalışmadığı (başarı/başarısızlık) sorgulanarak senaryolar üretilir.
\end{testcard}

\begin{testcard}{Soru 38}
    İşyeri aydınlatmasının ergonomik standartlara uygun olmamasının (yetersiz veya kamaşma yaratan ışık) çalışana etkisi hangisidir?
    \begin{itemize}
        \item[A)] Kas erimesi
        \item[B)] Göz yorgunluğu, baş ağrısı ve dikkat dağılmasına bağlı iş kazası riskinde artış
        \item[C)] İşitme kaybı
        \item[D)] Alerjik reaksiyonlar
        \item[E)] Hiçbiri
    \end{itemize}
    \tcblower
    \textbf{Doğru Cevap: B} \\
    \textbf{Açıklama:} Kötü aydınlatma göz sağlığını bozar ve tehlikelerin fark edilmesini zorlaştırarak kazaları tetikler.
\end{testcard}

\begin{testcard}{Soru 39}
    Ofis ergonomisinde, otururken uyluk ile gövde arasındaki açının yaklaşık kaç derece olması tavsiye edilir?
    \begin{itemize}
        \item[A)] 45 derece
        \item[B)] 90-100 derece
        \item[C)] 180 derece
        \item[D)] 15 derece
        \item[E)] Hiçbiri
    \end{itemize}
    \tcblower
    \textbf{Doğru Cevap: B} \\
    \textbf{Açıklama:} Doğal ve rahat bir oturuş için dizler ve kalça açısı 90-100 derece civarında olmalıdır.
\end{testcard}

\begin{testcard}{Soru 40}
    FMEA analizinde yer alan 'S' katsayısının (Şiddet) puanlamasında, çalışanın ölümüne yol açabilecek kritik hata türlerine kaç puan verilir?
    \begin{itemize}
        \item[A)] 1 Puan
        \item[B)] 3 Puan
        \item[C)] 5 Puan
        \item[D)] 9 veya 10 Puan
        \item[E)] Hiçbiri
    \end{itemize}
    \tcblower
    \textbf{Doğru Cevap: D} \\
    \textbf{Açıklama:} Ölümcül veya ağır yaralanmalı sonuçları olan hata türlerinin şiddet derecesi FMEA tablosunda en yüksek (9 veya 10) puanlanır.
\end{testcard}

\begin{testcard}{Soru 41}
    Aşağıdakilerden hangisi Fine-Kinney metodunda maruz kalma sıklığı (frekans) için kullanılan en yüksek puanı (10) tanımlar?
    \begin{itemize}
        \item[A)] Sürekli veya günde birçok kez maruz kalma
        \item[B)] Yılda bir kez karşılaşma
        \item[C)] Sadece işe ilk girişte karşılaşma
        \item[D)] Ayda bir karşılaşma
        \item[E)] Hiçbiri
    \end{itemize}
    \tcblower
    \textbf{Doğru Cevap: A} \\
    \textbf{Açıklama:} Frekans skalasında 10 puan, tehlikeye kesintisiz ya da gün içinde çok sık maruz kalındığını belirtir.
\end{testcard}

\begin{testcard}{Soru 42}
    Bir risk değerlendirmesi çalışmasının başarısını gösteren en önemli kriter hangisidir?
    \begin{itemize}
        \item[A)] Sadece kağıt üzerinde arşivlenmesi
        \item[B)] Aksiyon planlarının hayata geçirilmesi ve kaza oranlarında düşüş sağlanması
        \item[C)] Çok kalın bir dosya haline getirilmesi
        \item[D)] İşyeri sahibinin odasında saklanması
        \item[E)] Hiçbiri
    \end{itemize}
    \tcblower
    \textbf{Doğru Cevap: B} \\
    \textbf{Açıklama:} Düzeltici tedbirler sahada uygulanmadığı sürece risk analizleri sadece bürokratik bir belgeden ibaret kalır.
\end{testcard}

\begin{testcard}{Soru 43}
    Ergonomik tasarımlarda 'antropometrik veri aralığı' belirlenirken genellikle nüfusun yüzde kaçı hedef alınır?
    \begin{itemize}
        \item[A)] Yüzde 100'ü
        \item[B)] Yüzde 5 ile Yüzde 95 arasındaki dilim (\%90'lık kesim)
        \item[C)] Sadece en uzun boylular
        \item[D)] Sadece en kısa boylular
        \item[E)] Hiçbiri
    \end{itemize}
    \tcblower
    \textbf{Doğru Cevap: B} \\
    \textbf{Açıklama:} Tasarımların çoğunluğa uyması için uç değerler (\%5 ve \%95 dışındakiler) elenerek ortadaki \%90'lık nüfus dilimi baz alınır.
\end{testcard}

\begin{testcard}{Soru 44}
    Hata Ağacı Analizindeki (FTA) mantıksal kapılar hangi matematiksel disipline dayanır?
    \begin{itemize}
        \item[A)] Trigonometri
        \item[B)] Boole Cebri (Mantık Matematiği)
        \item[C)] Geometri
        \item[D)] Matris Determinantı
        \item[E)] Hiçbiri
    \end{itemize}
    \tcblower
    \textbf{Doğru Cevap: B} \\
    \textbf{Açıklama:} FTA'daki VE/VEYA kapıları ikili mantık sistemine (0 ve 1) ve Boole cebri kurallarına dayanır.
\end{testcard}

\begin{testcard}{Soru 45}
    FMEA analizi yaparken 'hata türü' (failure mode) neyi ifade eder?
    \begin{itemize}
        \item[A)] Kaza sonrası ödenen tazminat miktarını
        \item[B)] Bir ürün veya prosesin işlevini yerine getirememe şeklini (nasıl bozulduğunu)
        \item[C)] İSG uzmanının sertifika numarasını
        \item[D)] Çalışanın eğitim durumunu
        \item[E)] Hiçbiri
    \end{itemize}
    \tcblower
    \textbf{Doğru Cevap: B} \\
    \textbf{Açıklama:} Hata türü, bir parçanın veya sürecin standart dışı çalışma veya arıza yapma biçimidir.
\end{testcard}

\begin{testcard}{Soru 46}
    Örgütsel (Organizasyonel) ergonomi aşağıdakilerden hangisini inceler?
    \begin{itemize}
        \item[A)] Vardiya düzenleri, çalışma süreleri, ekip çalışması ve kalite yönetimi
        \item[B)] Masa yüksekliğini
        \item[C)] Monitör parlaklığını
        \item[D)] Tornavida sapı şeklini
        \item[E)] Hiçbiri
    \end{itemize}
    \tcblower
    \textbf{Doğru Cevap: A} \\
    \textbf{Açıklama:} Örgütsel ergonomi, İSG sistemlerini, organizasyonel yapıları ve süreç zamanlamalarını optimize eder.
\end{testcard}

\begin{testcard}{Soru 47}
    HAZOP analizinde kullanılan 'Eksik / Yok' (No / Less) kılavuz kelimesi hangisine yol açabilir?
    \begin{itemize}
        \item[A)] Boru hattında akışın kesilmesi veya vananın tıkanması
        \item[B)] Reaktörün aşırı dolması
        \item[C)] Basıncın iki katına çıkması
        \item[D)] Sıvının ısınması
        \item[E)] Hiçbiri
    \end{itemize}
    \tcblower
    \textbf{Doğru Cevap: A} \\
    \textbf{Açıklama:} Akış parametresine 'Eksik/Yok' uygulandığında akışın durması veya yetersiz besleme senaryoları incelenir.
\end{testcard}

\begin{testcard}{Soru 48}
    Çalışma ortamındaki gürültünün çalışanın konsantrasyonu ve zihinsel performansı üzerindeki olumsuz etkisi ergonominin hangi boyutuyla ilgilidir?
    \begin{itemize}
        \item[A)] Fiziksel Ergonomi
        \item[B)] Bilişsel ve Çevresel Ergonomi
        \item[C)] Antropometri
        \item[D)] Anatomik Ergonomi
        \item[E)] Hiçbiri
    \end{itemize}
    \tcblower
    \textbf{Doğru Cevap: B} \\
    \textbf{Açıklama:} Gürültünün dikkat dağıtması ve zihinsel yük yaratması bilişsel ve çevresel ergonominin konusudur.
\end{testcard}

\begin{testcard}{Soru 49}
    Risk Değerlendirmesi Yönetmeliği uyarınca, yeni bir makine alınması veya üretim yönteminin değiştirilmesi durumunda risk değerlendirmesi nasıl güncellenmelidir?
    \begin{itemize}
        \item[A)] Yasal sürenin (2/4/6 yıl) dolması beklenir
        \item[B)] Değişiklikten etkilenen alanlar için kısmen veya tamamen derhal yenilenmelidir
        \item[C)] Güncellemeye gerek yoktur
        \item[D)] Sadece iş müfettişi geldiğinde güncellenir
        \item[E)] Hiçbiri
    \end{itemize}
    \tcblower
    \textbf{Doğru Cevap: B} \\
    \textbf{Açıklama:} İşyerinde yeni bir tehlike ortaya çıktığında veya teknoloji değiştiğinde risk analizi derhal güncellenir.
\end{testcard}

\begin{testcard}{Soru 50}
    Bir risk analizinde olasılık değerinin düşürülmesi ne demektir?
    \begin{itemize}
        \item[A)] Kazanın zarar şiddetini azaltmak
        \item[B)] Kazanın meydana gelme sıklığını/ihtimalini azaltacak önlemler almak
        \item[C)] Rapor sayısını azaltmak
        \item[D)] Çalışan sayısını artırmak
        \item[E)] Hiçbiri
    \end{itemize}
    \tcblower
    \textbf{Doğru Cevap: B} \\
    \textbf{Açıklama:} Olasılığın düşürülmesi, o kazayı tetikleyen kök nedenlerin kontrol altına alınarak gerçekleşme şansının azaltılmasıdır.
\end{testcard}

\begin{testcard}{Soru 51}
    FMEA'da düzeltici önlemler alındıktan sonra ne yapılır?
    \begin{itemize}
        \item[A)] Analiz sonlandırılır
        \item[B)] Yeni parametrelerle RÖS puanı tekrar hesaplanır ve riskin kabul edilebilir sınıra inip inmediği kontrol edilir
        \item[C)] Doküman çöpe atılır
        \item[D)] Eski RÖS puanı aynen saklanır
        \item[E)] Hiçbiri
    \end{itemize}
    \tcblower
    \textbf{Doğru Cevap: B} \\
    \textbf{Açıklama:} FMEA dinamik bir süreçtir. Tedbir sonrası 'Kalıntı Risk' hesaplanarak RÖS tekrar ölçülür.
\end{testcard}

\begin{testcard}{Soru 52}
    Çalışma koltuğunun kolçaklarının bulunmasının ergonomik faydası nedir?
    \begin{itemize}
        \item[A)] Çalışanın uyumasını kolaylaştırmak
        \item[B)] Omuz ve boyun kaslarındaki statik yükü azaltmak
        \item[C)] Koltuğu ağırlaştırmak
        \item[D)] Çalışanın ayağa kalkmasını engellemek
        \item[E)] Hiçbiri
    \end{itemize}
    \tcblower
    \textbf{Doğru Cevap: B} \\
    \textbf{Açıklama:} Kolçaklar dirsekleri destekleyerek omuz kuşağı ve boyun kaslarının gerilmesini önler.
\end{testcard}

\begin{testcard}{Soru 53}
    Hata Ağacı Analizinde (FTA) 'kapı' (gate) kavramı neyi ifade eder?
    \begin{itemize}
        \item[A)] Fabrikanın giriş kapısını
        \item[B)] Alt olayların bir üst olayı nasıl tetiklediğini belirleyen mantıksal ilişkileri
        \item[C)] Yangın çıkış kapılarını
        \item[D)] Güvenlik kulübesini
        \item[E)] Hiçbiri
    \end{itemize}
    \tcblower
    \textbf{Doğru Cevap: B} \\
    \textbf{Açıklama:} Kapılar (VE, VEYA vb.) mantıksal geçişleri ve olayların birleşme kurallarını temsil eder.
\end{testcard}

\begin{testcard}{Soru 54}
    İşyerinde ergonomi ilkelerine uyulmasının işletmeye sağladığı dolaylı fayda hangisidir?
    \begin{itemize}
        \item[A)] Hammadde fiyatlarının düşmesi
        \item[B)] İş kazası ve meslek hastalığı oranlarının azalması, verimlilik ve motivasyonun artması
        \item[C)] Vergilerin azalması
        \item[D)] Elektrik faturasının düşmesi
        \item[E)] Hiçbiri
    \end{itemize}
    \tcblower
    \textbf{Doğru Cevap: B} \\
    \textbf{Açıklama:} Ergonomi çalışan sağlığını korurken hata oranlarını düşürür, kayıp günleri azaltır ve verimliliği artırır.
\end{testcard}

\begin{testcard}{Soru 55}
    İSG risk değerlendirmelerinde 'Kabul Edilebilir Risk' düzeyi neyi ifade eder?
    \begin{itemize}
        \item[A)] Yasal yükümlülüklere uygun ve işletmenin tolere edebileceği düzeye indirilmiş risk seviyesi
        \item[B)] Sıfır riskli durum
        \item[C)] Hiçbir önlem alınmamış durum
        \item[D)] Kazanın kesinlikle olacağı durum
        \item[E)] Hiçbiri
    \end{itemize}
    \tcblower
    \textbf{Doğru Cevap: A} \\
    \textbf{Açıklama:} Gerçek hayatta risk tamamen sıfırlanamaz. Kabul edilebilir risk, yasal sınırlara uygun ve kontrol altındaki tolere edilebilir seviyedir.
\end{testcard}

\begin{testcard}{Soru 56}
    Sağlık ve Güvenlik İşaretleri Yönetmeliği'ne göre, kırmızı renkli daire biçiminde, beyaz zemin üzerine siyah piktogramlı ve kırmızı diyagonal (eğik) çizgili bir levha neyi ifade eder ve hangi işaret grubuna girer?
    \begin{itemize}
        \item[A)] Yangınla mücadele işareti - Uyarı
        \item[B)] Yasaklayıcı işaret - Yasaklama
        \item[C)] Acil çıkış işareti - Bilgilendirme
        \item[D)] Emredici işaret - Zorunluluk
        \item[E)] Hiçbiri
    \end{itemize}
    \tcblower
    \textbf{Doğru Cevap: B} \\
    \textbf{Açıklama:} Kırmızı renkli, yuvarlak şekilli, siyah piktogramlı ve kırmızı diyagonal çizgili işaretler yasaklayıcı işaretlerdir. (Örn: Sigara içilmez, Giriş yasaktır).
\end{testcard}

\begin{testcard}{Soru 57}
    İşyerinde sarı üçgen zemin üzerine siyah piktogram ve siyah çerçeve ile belirtilen bir güvenlik işareti neyi temsil eder?
    \begin{itemize}
        \item[A)] Zorunlu eylem veya kişisel koruyucu donanım takma gerekliliği
        \item[B)] Yangın söndürme cihazlarının yerleri
        \item[C)] Belirli bir tehlike veya riske karşı uyarı (Warning/Caution)
        \item[D)] Acil kaçış yolları ve ilk yardım noktaları
        \item[E)] Hiçbiri
    \end{itemize}
    \tcblower
    \textbf{Doğru Cevap: C} \\
    \textbf{Açıklama:} Sarı renkli ve üçgen şekilli, siyah piktogramlı işaretler uyarı işaretleridir. (Örn: Yüksek gerilim, Asılı yük, İş makinesi).
\end{testcard}

\begin{testcard}{Soru 58}
    Daire (yuvarlak) şekilli, mavi zemin üzerine beyaz piktogram içeren güvenlik ve sağlık işaretlerinin temel amacı nedir?
    \begin{itemize}
        \item[A)] Bir tehlikeyi bildirmek ve uyarmak
        \item[B)] Yangın ekipmanlarının yerini göstermek
        \item[C)] Belirli bir davranışı veya kişisel koruyucu donanım kullanımını zorunlu kılmak (Emredici/Mandatory)
        \item[D)] Yasaklanmış olan eylemleri göstermek
        \item[E)] Hiçbiri
    \end{itemize}
    \tcblower
    \textbf{Doğru Cevap: C} \\
    \textbf{Açıklama:} Mavi zemin üzerine beyaz piktogramlı yuvarlak işaretler emredici/zorunluluk bildiren işaretlerdir. (Örn: Baret tak, Kulaklık kullan).
\end{testcard}

\begin{testcard}{Soru 59}
    Kişisel Koruyucu Donanımların (KKD) işyerinde kullanımı ve seçimiyle ilgili yasal ve pratik esaslar dikkate alındığında aşağıdakilerden hangisi yanlıştır?
    \begin{itemize}
        \item[A)] KKD'ler kendisi ek bir risk oluşturmadan ilgili riski önlemeye uygun olmalıdır.
        \item[B)] KKD'ler ergonomik gereksinimlere uygun olmalı ve kullanan çalışanın sağlık durumuna uymalıdır.
        \item[C)] Riskleri önlemek için öncelikli olarak KKD kullanımına başvurulmalı, toplu korunma önlemleri en son çare olarak düşünülmelidir.
        \item[D)] KKD'lerin temin edilmesi ve bakımıyla ilgili tüm maliyetler işveren tarafından karşılanmalı, çalışana yansıtılmamalıdır.
        \item[E)] Hiçbiri
    \end{itemize}
    \tcblower
    \textbf{Doğru Cevap: C} \\
    \textbf{Açıklama:} Risk kontrol hiyerarşisinde kişisel koruyucu donanım (KKD) kullanımı en son çaredir. Öncelik her zaman tehlikeyi kaynakta yok etmekte veya toplu koruma önlemlerindedir.
\end{testcard}

\begin{testcard}{Soru 60}
    Bir kişisel koruyucu donanımın (KKD) tasarımı ve taşıması gereken özelliklerle ilgili olarak, "Kendisi ek bir risk oluşturmamalı, çalışma şartlarına ve çalışana tam uyum sağlamalıdır" ilkesi neyi ifade eder?
    \begin{itemize}
        \item[A)] Donanımın sadece görsel şıklığını
        \item[B)] KKD'nin koruyucu etkisinin yanında çalışan üzerinde ek bir yük, tehlike veya engelleme yaratmaması gerekliliğini
        \item[C)] Donanımın piyasadaki en ucuz ürün olması gerektiğini
        \item[D)] Donanımın sadece çok tehlikeli sınıfta kullanılabilir olduğunu
        \item[E)] Hiçbiri
    \end{itemize}
    \tcblower
    \textbf{Doğru Cevap: B} \\
    \textbf{Açıklama:} KKD'ler çalışanı korurken hareket kısıtlılığı, görüş engeli veya cildinde tahriş gibi yeni tehlikeler yaratmamalı; işe ve kişiye uygun olarak tasarlanmalıdır.
\end{testcard}

\begin{testcard}{Soru 61}
    Fine-Kinney veya L Tipi Matris gibi risk analizi yöntemlerinde, risk puanının '15, 16, 20, 25' gibi en yüksek aralıklarda çıkıp 'Kabul Edilemez Risk' bölgesinde yer alması durumunda yapılması gereken yasal ve operasyonel işlem hangisidir?
    \begin{itemize}
        \item[A)] Mevcut önlemler sürdürülmeli ve süreç periyodik olarak izlenmelidir.
        \item[B)] Risk kabul edilebilir seviyeye çekilene kadar çalışma derhal durdurulmalı, gerekirse risk düşürülene kadar işe başlanmamalıdır.
        \item[C)] Ekipmanların temizliği yapılarak çalışmaya aynı şekilde devam edilir.
        \item[D)] Risk sadece yazılı olarak kaydedilir ve herhangi bir fiili müdahale yapılmaz.
        \item[E)] Hiçbiri
    \end{itemize}
    \tcblower
    \textbf{Doğru Cevap: B} \\
    \textbf{Açıklama:} Ders notu 9, Sayfa 139'daki verilere göre 15, 16, 20, 25 puan alan riskler 'Kabul Edilemez Risk' sınıfındadır. Bu durumda iş durdurulur ve risk kabul edilebilir seviyeye çekilmeden çalışmaya başlanamaz.
\end{testcard}

\begin{testcard}{Soru 62}
    Risk değerlendirme tablosunda hesaplanan risk değerinin '8, 9, 10, 12' aralığında çıkarak 'Dikkate Değer Risk' olarak tanımlanması durumunda izlenmesi gereken eylem planı aşağıdakilerden hangisi doğrudur?
    \begin{itemize}
        \item[A)] Acil tedbir gerektirmeyebilir ancak izlenmesi ve kontrolü sürdürülmelidir.
        \item[B)] İş derhal durdurulur ve fabrika tahliye edilir.
        \item[C)] Bu risklere mümkün olduğunca çabuk müdahale edilmeli, acil tedbirler veya ilave tedbirler planlanarak kontrol altına alınmalıdır.
        \item[D)] Bu puanlar tamamen risksiz bölgeleri temsil ettiği için hiçbir eyleme gerek yoktur.
        \item[E)] Hiçbiri
    \end{itemize}
    \tcblower
    \textbf{Doğru Cevap: C} \\
    \textbf{Açıklama:} Sayfa 139'a göre 8, 9, 10, 12 puanlı 'Dikkate Değer Riskler' için mümkün olduğunca çabuk müdahale edilmeli ve ilave tedbirler/acil eylemler alınmalıdır.
\end{testcard}

\begin{testcard}{Soru 63}
    Risk puanı '1, 2, 3' veya '4, 5, 6' olarak hesaplanan bir tehlike için risk tablosundaki karşılık gelen durum ve eylem nedir?
    \begin{itemize}
        \item[A)] Kabul edilemez risk olup iş derhal durdurulmalıdır.
        \item[B)] Kabul edilebilir risk olup acil tedbir gerektirmeyebilir, ancak mevcut kontroller sürdürülmeli ve iyileştirilerek izlenmeye devam edilmelidir.
        \item[C)] Bu tehlike tamamen listeden çıkarılmalı ve bir daha değerlendirilmemelidir.
        \item[D)] Çalışanlar işten uzaklaştırılarak makine yenilenmelidir.
        \item[E)] Hiçbiri
    \end{itemize}
    \tcblower
    \textbf{Doğru Cevap: B} \\
    \textbf{Açıklama:} Sayfa 139'a göre 1-6 arası riskler 'Kabul Edilebilir Risk' sınıfındadır. Acil önlem gerektirmeyebilir ama mevcut tedbirler korunarak izlemeye ve iyileştirmeye devam edilir.
\end{testcard}

\begin{testcard}{Soru 64}
    Risk kontrol önlemlerinin hiyerarşik sıralamasında yer alan 'İkame' (Substitution) yönteminin temel mantığı ve bu yönteme uygun örnek aşağıdakilerin hangisinde doğru verilmiştir?
    \begin{itemize}
        \item[A)] Tehlikeli makineyi tamamen kaldırmak - Makineyi hurdaya ayırmak
        \item[B)] Tehlikeli olan bir kaynağı veya yöntemi daha az tehlikeli olanla değiştirmek - Solvent bazlı tiner boya yerine su bazlı boya kullanılması
        \item[C)] Çalışana koruyucu eldiven ve gözlük dağıtmak - KKD temini
        \item[D)] Riskli bölgeye uyarı levhası asmak - İdari önlem
        \item[E)] Hiçbiri
    \end{itemize}
    \tcblower
    \textbf{Doğru Cevap: B} \\
    \textbf{Açıklama:} İkame, tehlikeli olanın yerine tehlikesiz veya daha az tehlikeli olanın getirilmesidir. Solvent bazlı (zehirli/parlayıcı) boya yerine su bazlı boya kullanılması klasik bir ikame örneğidir.
\end{testcard}

\subsection{Ünite 3: İSG Kurulları ve İş Kazaları}
\begin{testcard}{Soru 1}
    6331 sayılı İSG Kanunu'na göre, bir işyerinde yasal olarak İSG kurulu kurulması için asgari çalışan sayısı kaç olmalıdır?
    \begin{itemize}
        \item[A)] En az 30 çalışan
        \item[B)] En az 50 çalışan
        \item[C)] En az 100 çalışan
        \item[D)] En az 250 çalışan
        \item[E)] Hiçbiri
    \end{itemize}
    \tcblower
    \textbf{Doğru Cevap: B} \\
    \textbf{Açıklama:} 6331 sayılı İSG Kanunu Madde 22 uyarınca fiilen 50 ve daha fazla çalışanın bulunduğu işyerlerinde kurul kurulmalıdır.
\end{testcard}

\begin{testcard}{Soru 2}
    İSG kurulunun kurulabilmesi için yapılan işin süre yönünden yasal özelliği ne olmalıdır?
    \begin{itemize}
        \item[A)] En az 1 ay süren geçici işler
        \item[B)] 6 aydan fazla süren sürekli işlerin yapıldığı yerler
        \item[C)] Sadece mevsimlik işler
        \item[D)] Süre kısıtlaması yoktur
        \item[E)] Hiçbiri
    \end{itemize}
    \tcblower
    \textbf{Doğru Cevap: B} \\
    \textbf{Açıklama:} Kurul kurulması için işyerinde yapılan işin 6 aydan fazla süren sürekli bir iş olması gerekir.
\end{testcard}

\begin{testcard}{Soru 3}
    İSG kurullarında kurulun başkanı yasal olarak kimdir?
    \begin{itemize}
        \item[A)] İş güvenliği uzmanı
        \item[B)] İşyeri hekimi
        \item[C)] İşveren veya işveren vekili
        \item[D)] Çalışan temsilcisi
        \item[E)] Hiçbiri
    \end{itemize}
    \tcblower
    \textbf{Doğru Cevap: C} \\
    \textbf{Açıklama:} İSG Kurulları yönetmeliğine göre kurulun başkanı işveren veya işveren vekilidir.
\end{testcard}

\begin{testcard}{Soru 4}
    İSG kurullarında kurulun sekretaryasını yürütmekle görevli üye yasal olarak kimdir?
    \begin{itemize}
        \item[A)] İşyeri hekimi
        \item[B)] İş güvenliği uzmanı
        \item[C)] Çalışan temsilcisi
        \item[D)] İdari işler müdürü
        \item[E)] Hiçbiri
    \end{itemize}
    \tcblower
    \textbf{Doğru Cevap: B} \\
    \textbf{Açıklama:} Kurulun sekreteri iş güvenliği uzmanıdır. Eğer uzman tam zamanlı değilse, sekretaryayı insan kaynakları veya idari işler üstlenir.
\end{testcard}

\begin{testcard}{Soru 5}
    İSG kurullarında kararlar oy çokluğuyla alınır. Oyların eşit çıkması durumunda nihai karar nasıl belirlenir?
    \begin{itemize}
        \item[A)] Kurul başkanının (işveren/vekilinin) oy verdiği yön çoğunluk sayılır
        \item[B)] Toplantı iptal edilir
        \item[C)] Uzmanın oyu çift sayılır
        \item[D)] Kura çekimi yapılır
        \item[E)] Hiçbiri
    \end{itemize}
    \tcblower
    \textbf{Doğru Cevap: A} \\
    \textbf{Açıklama:} Oyların eşitliği halinde kurul başkanının oyu yönünde karar kesinleşmiş sayılır.
\end{testcard}

\begin{testcard}{Soru 6}
    İSG kurulunun yasal toplantı periyodu normal şartlarda en az ne kadardır?
    \begin{itemize}
        \item[A)] Haftada bir
        \item[B)] Ayda en az bir kez
        \item[C)] 3 ayda bir
        \item[D)] 6 ayda bir
        \item[E)] Hiçbiri
    \end{itemize}
    \tcblower
    \textbf{Doğru Cevap: B} \\
    \textbf{Açıklama:} İSG kurulları olağan olarak ayda en az bir kez toplanmak zorundadır.
\end{testcard}

\begin{testcard}{Soru 7}
    Tehlikeli sınıfta yer alan bir işyerinde kurul kararıyla toplantı periyodu en geç kaç ayda bir olacak şekilde uzatılabilir?
    \begin{itemize}
        \item[A)] 2 ayda bir
        \item[B)] 3 ayda bir
        \item[C)] 6 ayda bir
        \item[D)] Uzatılamaz, her ay zorunludur
        \item[E)] Hiçbiri
    \end{itemize}
    \tcblower
    \textbf{Doğru Cevap: A} \\
    \textbf{Açıklama:} Tehlikeli sınıftaki işyerlerinde kurul periyodu kurulun kararıyla en fazla 2 ayda bire kadar uzatılabilir.
\end{testcard}

\begin{testcard}{Soru 8}
    Az Tehlikeli sınıfta yer alan bir işyerinde kurul kararıyla toplantı periyodu en geç kaç ayda bir olacak şekilde uzatılabilir?
    \begin{itemize}
        \item[A)] 2 ay
        \item[B)] 3 ay
        \item[C)] 4 ay
        \item[D)] 6 ay
        \item[E)] Hiçbiri
    \end{itemize}
    \tcblower
    \textbf{Doğru Cevap: B} \\
    \textbf{Açıklama:} Az Tehlikeli sınıftaki işyerlerinde kurul toplantı sıklığı en fazla 3 ayda bire kadar uzatılabilmektedir.
\end{testcard}

\begin{testcard}{Soru 9}
    İşyerinde yaşanan bir iş kazasını işverenin SGK'ya (Sosyal Güvenlik Kurumu) bildirmesi için yasal süre nedir?
    \begin{itemize}
        \item[A)] Kazadan sonraki 3 iş günü içinde
        \item[B)] Kazanın olduğu gün mesai bitimine kadar
        \item[C)] Kazadan sonraki 24 saat içinde
        \item[D)] Haftanın son cuma günü
        \item[E)] Hiçbiri
    \end{itemize}
    \tcblower
    \textbf{Doğru Cevap: A} \\
    \textbf{Açıklama:} İşveren iş kazalarını, kazanın gerçekleştiği günü takip eden 3 iş günü içinde SGK'ya bildirmekle yükümlüdür.
\end{testcard}

\begin{testcard}{Soru 10}
    İşyeri hekimi veya sağlık hizmeti sunucuları, teşhis ettikleri meslek hastalıklarını en geç kaç gün içinde SGK'ya bildirmelidir?
    \begin{itemize}
        \item[A)] 3 iş günü içinde
        \item[B)] 5 iş günü içinde
        \item[C)] 10 gün içinde
        \item[D)] Aynı gün
        \item[E)] Hiçbiri
    \end{itemize}
    \tcblower
    \textbf{Doğru Cevap: C} \\
    \textbf{Açıklama:} Sağlık hizmeti sunucuları kendilerine intikal eden veya teşhis koydukları meslek hastalıklarını en geç 10 gün içinde SGK'ya bildirir.
\end{testcard}

\begin{testcard}{Soru 11}
    5510 Sayılı Kanun Madde 13'e göre, aşağıdakilerden hangisi bir olayın iş kazası sayılması için gerekli şartlardan biri değildir?
    \begin{itemize}
        \item[A)] Kazaya uğrayanın yasal olarak sigortalı sayılması
        \item[B)] Olayın işyerinde veya işverenin yürütmekte olduğu iş nedeniyle olması
        \item[C)] Sigortalının bedenen veya ruhen engelli hale gelmesi
        \item[D)] Çalışanın en az 10 yıllık kıdeminin olması
        \item[E)] Hiçbiri
    \end{itemize}
    \tcblower
    \textbf{Doğru Cevap: D} \\
    \textbf{Açıklama:} Kıdem süresinin iş kazası tanımıyla ilgisi yoktur. İşe girildiği ilk gün hatta ilk saniye yaşanan kaza da iş kazasıdır.
\end{testcard}

\begin{testcard}{Soru 12}
    Emziren kadın çalışanların çocuklarına süt vermek için ayrılan zamanlarda geçirdikleri kazalar yasal olarak hangi sınıfa girer?
    \begin{itemize}
        \item[A)] Ev kazası
        \item[B)] İş kazası (5510 md. 13)
        \item[C)] Trafik kazası
        \item[D)] Doğal afet hasarı
        \item[E)] Hiçbiri
    \end{itemize}
    \tcblower
    \textbf{Doğru Cevap: B} \\
    \textbf{Açıklama:} Kadın çalışanın çocuk emzirme izni saatlerinde başına gelen kazalar yasal olarak iş kazası kabul edilir.
\end{testcard}

\begin{testcard}{Soru 13}
    İşverence sağlanan bir servis aracıyla çalışanların topluca işe gidip gelişi sırasında meydana gelen trafik kazası yasal olarak ne kabul edilir?
    \begin{itemize}
        \item[A)] Genel trafik kazası (İSG kapsamı dışı)
        \item[B)] İş kazası
        \item[C)] Özel seyahat kazası
        \item[D)] Doğal kaza
        \item[E)] Hiçbiri
    \end{itemize}
    \tcblower
    \textbf{Doğru Cevap: B} \\
    \textbf{Açıklama:} İşverenin sağladığı taşıtla işin yapıldığı yere gidiş geliş esnasında yaşanan kazalar iş kazasıdır.
\end{testcard}

\begin{testcard}{Soru 14}
    NACE kodları sistemi ne işe yarar?
    \begin{itemize}
        \item[A)] Şirketlerin vergi borçlarını hesaplamaya
        \item[B)] Ekonomik faaliyet kollarına göre işyerlerinin tehlike sınıfını belirlemeye
        \item[C)] Çalışanların performans puanlarını ölçmeye
        \item[D)] Yıllık izin günlerini kaydetmeye
        \item[E)] Hiçbiri
    \end{itemize}
    \tcblower
    \textbf{Doğru Cevap: B} \\
    \textbf{Açıklama:} NACE kodu, işyerinde yürütülen ana faaliyetin kodudur ve tehlike sınıfının (az tehlikeli/tehlikeli/çok tehlikeli) belirlenmesinde esastır.
\end{testcard}

\begin{testcard}{Soru 15}
    Tehlike sınıfı tebliğine göre, inşaat işleri (bina yapımı, yıkımı vb.) hangi tehlike sınıfında yer alır?
    \begin{itemize}
        \item[A)] Az Tehlikeli
        \item[B)] Tehlikeli
        \item[C)] Çok Tehlikeli
        \item[D)] Tehlikesiz muaf sınıf
        \item[E)] Hiçbiri
    \end{itemize}
    \tcblower
    \textbf{Doğru Cevap: C} \\
    \textbf{Açıklama:} Yapı şantiyeleri, madenler, tersaneler yüksek risk barındırdığı için Çok Tehlikeli sınıftadır.
\end{testcard}

\begin{testcard}{Soru 16}
    Güvenlik ve Sağlık İşaretleri Yönetmeliğine göre, 'Kırmızı' renkli daire biçimindeki işaretlerin temel anlamı nedir?
    \begin{itemize}
        \item[A)] Uyarı (dikkat tehlike)
        \item[B)] Yasaklayıcı işaret veya yangın ekipmanı
        \item[C)] Zorunluluk (KKD takma)
        \item[D)] Acil çıkış ve ilk yardım
        \item[E)] Hiçbiri
    \end{itemize}
    \tcblower
    \textbf{Doğru Cevap: B} \\
    \textbf{Açıklama:} Kırmızı renk; yasaklama, tehlike alarmı ve yangınla mücadele ekipmanlarının işaretlenmesinde kullanılır.
\end{testcard}

\begin{testcard}{Soru 17}
    Güvenlik ve Sağlık İşaretleri Yönetmeliğine göre, 'Mavi' renkli daire şeklindeki işaretlerin anlamı nedir?
    \begin{itemize}
        \item[A)] Yasaklama
        \item[B)] Emredici / Zorunluluk işareti (KKD giyme zorunluluğu vb.)
        \item[C)] Yangın ekipmanı
        \item[D)] Acil durum kaçış yolu
        \item[E)] Hiçbiri
    \end{itemize}
    \tcblower
    \textbf{Doğru Cevap: B} \\
    \textbf{Açıklama:} Mavi zemin üzerine beyaz piktogramlı yuvarlak işaretler kişisel koruyucu donanım takma veya belirli bir davranışı yapma zorunluluğunu belirtir.
\end{testcard}

\begin{testcard}{Soru 18}
    Güvenlik ve Sağlık İşaretleri Yönetmeliğine göre, 'Sarı' renkli üçgen işaretlerin anlamı nedir?
    \begin{itemize}
        \item[A)] Yangınla mücadele
        \item[B)] Uyarı işareti (dikkat tehlike olasılığı)
        \item[C)] Acil çıkış kapısı
        \item[D)] Zorunlu baret kullanımı
        \item[E)] Hiçbiri
    \end{itemize}
    \tcblower
    \textbf{Doğru Cevap: B} \\
    \textbf{Açıklama:} Sarı zemin üzerine siyah piktogramlı üçgen işaretler çalışanları çevrelerindeki tehlikelere karşı uyarma amaçlıdır.
\end{testcard}

\begin{testcard}{Soru 19}
    Güvenlik ve Sağlık İşaretleri Yönetmeliğine göre, 'Yeşil' renkli dikdörtgen veya kare işaretlerin anlamı nedir?
    \begin{itemize}
        \item[A)] Yasaklama
        \item[B)] Uyarı
        \item[C)] Acil çıkış, ilk yardım veya güvenli durum işaretleri
        \item[D)] Yangın hortumu yeri
        \item[E)] Hiçbiri
    \end{itemize}
    \tcblower
    \textbf{Doğru Cevap: C} \\
    \textbf{Açıklama:} Yeşil renkli tabelalar acil kaçış yolları, acil duşlar, ilk yardım odaları ve güvenli alanları gösterir.
\end{testcard}

\begin{testcard}{Soru 20}
    İşyerinde kurulacak İSG kurulunun kararlarının bağlayıcılığı hakkında hangisi doğrudur?
    \begin{itemize}
        \item[A)] Kurul kararları sadece tavsiye niteliğindedir, işveren uymak zorunda değildir
        \item[B)] Kurulun aldığı kararlar işveren ve çalışanlar için bağlayıcıdır, uygulanması zorunludur
        \item[C)] Sadece çalışanları bağlar, yönetimi bağlamaz
        \item[D)] Yasal hiçbir geçerliliği yoktur
        \item[E)] Hiçbiri
    \end{itemize}
    \tcblower
    \textbf{Doğru Cevap: B} \\
    \textbf{Açıklama:} İSG Kurulu kararları işyerinde resmi karar hükmündedir ve işveren bu kararları uygulamakla yükümlüdür.
\end{testcard}

\begin{testcard}{Soru 21}
    Aşağıdakilerden hangisi İSG Kurulu üyeleri arasında bulunması zorunlu olan kişilerden biri değildir?
    \begin{itemize}
        \item[A)] İşveren veya işveren vekili
        \item[B)] İş güvenliği uzmanı ve işyeri hekimi
        \item[C)] Çalışan temsilcisi ve sendika temsilcisi
        \item[D)] İşyerinin en büyük müşterisi
        \item[E)] Hiçbiri
    \end{itemize}
    \tcblower
    \textbf{Doğru Cevap: D} \\
    \textbf{Açıklama:} Müşteriler İSG kurulu üyesi olamazlar.
\end{testcard}

\begin{testcard}{Soru 22}
    NACE kodu ve tehlike sınıfı tebliğine göre, 'Ofis ve yazılım geliştirme hizmetleri' hangi tehlike sınıfındadır?
    \begin{itemize}
        \item[A)] Az Tehlikeli
        \item[B)] Tehlikeli
        \item[C)] Çok Tehlikeli
        \item[D)] Sorumluluk dışı muaf
        \item[E)] Hiçbiri
    \end{itemize}
    \tcblower
    \textbf{Doğru Cevap: A} \\
    \textbf{Açıklama:} Ofis ortamları, bankacılık ve yazılım işleri düşük riskli kabul edildiğinden Az Tehlikeli sınıftadır.
\end{testcard}

\begin{testcard}{Soru 23}
    NACE kodu ve tehlike sınıfı tebliğine göre, 'Mobilya ve ahşap imalat atölyeleri' hangi tehlike sınıfına girmektedir?
    \begin{itemize}
        \item[A)] Az Tehlikeli
        \item[B)] Tehlikeli
        \item[C)] Çok Tehlikeli
        \item[D)] Özel muafiyetli sınıf
        \item[E)] Hiçbiri
    \end{itemize}
    \tcblower
    \textbf{Doğru Cevap: B} \\
    \textbf{Açıklama:} Ahşap kesim, mobilya üretimi ve boyama işleri barındırdığı riskler nedeniyle Tehlikeli sınıf kapsamındadır.
\end{testcard}

\begin{testcard}{Soru 24}
    Bir iş kazası sonrasında İSG kurulunun yapması gereken ilk eylemlerden biri hangisidir?
    \begin{itemize}
        \item[A)] Kazayı araştırmak, kök neden analizi yapmak ve tekrarlanmaması için önlemler belirlemek
        \item[B)] Kazayı basından gizlemek
        \item[C)] Kazalı çalışanın sözleşmesini derhal feshetmek
        \item[D)] Fabrikayı satmak
        \item[E)] Hiçbiri
    \end{itemize}
    \tcblower
    \textbf{Doğru Cevap: A} \\
    \textbf{Açıklama:} Kurul, iş kazası sonrasında olayı inceleyerek rapor hazırlar ve düzeltici faaliyetleri karara bağlar.
\end{testcard}

\begin{testcard}{Soru 25}
    Aslıhan Kayık Altınalp Hoca'nın ders notuna göre, iş kazası sonrasında SGK'ya yapılacak bildirimde süre hesaplanırken hangi günler hesaba katılmaz?
    \begin{itemize}
        \item[A)] Resmi tatiller ve hafta sonları (iş günü hesabı yapıldığından)
        \item[B)] Pazartesi günleri
        \item[C)] Yarım gün mesailer
        \item[D)] Hiçbir gün hariç tutulmaz, takvim günü sayılır
        \item[E)] Hiçbiri
    \end{itemize}
    \tcblower
    \textbf{Doğru Cevap: A} \\
    \textbf{Açıklama:} Yasal bildirim süresi 3 'iş günü'dür. Resmi tatil ve pazar günleri iş günü sayılmadığı için süre hesabına dahil edilmez.
\end{testcard}

\begin{testcard}{Soru 26}
    İşyerinde birden fazla iş güvenliği uzmanı görev yapıyorsa, İSG kuruluna hangi uzman katılır?
    \begin{itemize}
        \item[A)] Kura ile belirlenen uzman
        \item[B)] Tam zamanlı çalışan veya işverenin görevlendirdiği tam zamanlı koordinatör uzman
        \item[C)] En genç olan uzman
        \item[D)] Hepsi birden sekreter olmak zorundadır
        \item[E)] Hiçbiri
    \end{itemize}
    \tcblower
    \textbf{Doğru Cevap: B} \\
    \textbf{Açıklama:} Koordinasyonu sağlamak adına işverenin görevlendirdiği koordinatör İSG uzmanı kurulda yer alır.
\end{testcard}

\begin{testcard}{Soru 27}
    Kurul üyelerinin İSG eğitimleri hakkında hangisi doğrudur?
    \begin{itemize}
        \item[A)] Kurul üyelerine kurulun görevleri, yasal hak ve sorumlulukları konusunda özel eğitim verilmelidir
        \item[B)] Eğitim alınması isteğe bağlıdır
        \item[C)] Sadece çalışan temsilcisi eğitim alır
        \item[D)] Eğitimler sadece uzaktan video izleme ile geçiştirilebilir
        \item[E)] Hiçbiri
    \end{itemize}
    \tcblower
    \textbf{Doğru Cevap: A} \\
    \textbf{Açıklama:} Kurul üyelerinin etkin çalışabilmesi için kurul mevzuatı ve işleyişi hakkında özel eğitim almaları yasal zorunluluktur.
\end{testcard}

\begin{testcard}{Soru 28}
    İSG kurullarının sekretaryası görevinde, toplantı tutanakları ve kararların yazıldığı deftere ne ad verilir?
    \begin{itemize}
        \item[A)] Günlük Mesai Defteri
        \item[B)] İSG Kurul Karar Defteri (onaylı defter)
        \item[C)] Muhasebe Yevmiye Defteri
        \item[D)] Müşteri Kayıt Defteri
        \item[E)] Hiçbiri
    \end{itemize}
    \tcblower
    \textbf{Doğru Cevap: B} \\
    \textbf{Açıklama:} Kurul kararları resmi olarak İSG Kurul Karar Defterine yazılır ve üyelerce imzalanır.
\end{testcard}

\begin{testcard}{Soru 29}
    İş kazası oranlarının yaş gruplarına göre dağılımını inceleyen SGK verilerine göre, iş kazaları en fazla hangi yaş grubunda görülmektedir?
    \begin{itemize}
        \item[A)] 18 yaş altı
        \item[B)] 25-29 yaş grubu (ve genel olarak 25-34 genç çalışanlar)
        \item[C)] 50-60 yaş arası
        \item[D)] 65 yaş üstü
        \item[E)] Hiçbiri
    \end{itemize}
    \tcblower
    \textbf{Doğru Cevap: B} \\
    \textbf{Açıklama:} İstatistiklere göre tecrübesizlik ve dinamizmin getirdiği risk almayla birleşen 25-29 yaş aralığı kazalara en çok maruz kalan gruptur (Ekstra.pdf Sayfa 14).
\end{testcard}

\begin{testcard}{Soru 30}
    Aşağıdakilerden hangisi iş kazasını tetikleyen insan faktörleri arasında yer alan 'fiziksel yetersizlik' kapsamında değerlendirilir?
    \begin{itemize}
        \item[A)] Kronik uykusuzluk ve aşırı yorgunluk
        \item[B)] İşyerindeki kötü havalandırma
        \item[C)] Makinede koruyucu olmaması
        \item[D)] Yanlış aydınlatma
        \item[E)] Hiçbiri
    \end{itemize}
    \tcblower
    \textbf{Doğru Cevap: A} \\
    \textbf{Açıklama:} Uykusuzluk, yorgunluk, görme veya işitme bozuklukları insana bağlı fiziksel yetersizlik etmenleridir.
\end{testcard}

\begin{testcard}{Soru 31}
    Aşağıdaki durumlardan hangisinde İSG kurulunun olağanüstü (acil) toplanması gerekir?
    \begin{itemize}
        \item[A)] İşyerinde ölümlü veya ağır yaralanmalı bir iş kazası meydana geldiğinde veya yakın bir tehlike ihbarında
        \item[B)] İşyerine yeni bir bilgisayar alındığında
        \item[C)] Çalışanlardan biri işe geç kaldığında
        \item[D)] Yıllık bilanço açıklandığında
        \item[E)] Hiçbiri
    \end{itemize}
    \tcblower
    \textbf{Doğru Cevap: A} \\
    \textbf{Açıklama:} Ağır kazalar, ramak kala durumlar veya ciddi tehlike tespitlerinde kurul acilen olağanüstü toplantıya çağrılır.
\end{testcard}

\begin{testcard}{Soru 32}
    İSG kurul toplantılarının gündemi kim tarafından belirlenir ve üyelere duyurulur?
    \begin{itemize}
        \item[A)] Sadece çalışan temsilcisi
        \item[B)] Kurul başkanı ve sekreteri tarafından ortaklaşa belirlenir ve toplantı öncesi üyelere bildirilir
        \item[C)] Sivil savunma uzmanı
        \item[D)] Şirket muhasebecisi
        \item[E)] Hiçbiri
    \end{itemize}
    \tcblower
    \textbf{Doğru Cevap: B} \\
    \textbf{Açıklama:} Gündem, kurul başkanı (işveren) ve sekreteri (uzman) tarafından hazırlanarak üyelere önceden iletilir.
\end{testcard}

\begin{testcard}{Soru 33}
    İşyerinde çalışan temsilcisi sayısı neye göre belirlenir?
    \begin{itemize}
        \item[A)] İşyerinin cirosuna göre
        \item[B)] Çalışan sayısına göre (belirli baremlerle)
        \item[C)] İSG uzmanının kıdemine göre
        \item[D)] İşyeri hekiminin kararına göre
        \item[E)] Hiçbiri
    \end{itemize}
    \tcblower
    \textbf{Doğru Cevap: B} \\
    \textbf{Açıklama:} Çalışan temsilcisi sayısı yasa gereği çalışan sayısına paralel olarak artar (örneğin 2-50 çalışan için 1, 51-100 için 2, vb.).
\end{testcard}

\begin{testcard}{Soru 34}
    İSG Kurulu kararlarının çalışanlara duyurulması yasal olarak nasıl yapılmalıdır?
    \begin{itemize}
        \item[A)] Duyuru panoları, e-posta veya ortak bilgilendirme kanalları ile herkesin görebileceği şekilde ilan edilmelidir
        \item[B)] Gizli tutulmalı, kimseye söylenmemelidir
        \item[C)] Sadece sendika odasında saklanır
        \item[D)] Sadece kulaktan kulağa söylenir
        \item[E)] Hiçbiri
    \end{itemize}
    \tcblower
    \textbf{Doğru Cevap: A} \\
    \textbf{Açıklama:} Alınan kararlar tüm çalışanları etkilediği için panolarda veya dijital ekranlarda ilan edilmelidir.
\end{testcard}

\begin{testcard}{Soru 35}
    İşyerinde yürütülen asıl işin tehlike sınıfının tespiti konusunda ortaya çıkan uyuşmazlıklarda nihai kararı kim verir?
    \begin{itemize}
        \item[A)] İşveren
        \item[B)] Çalışma ve Sosyal Güvenlik Bakanlığı (yetkili mahkemeler yoluyla)
        \item[C)] Belediye başkanlığı
        \item[D)] Ticaret odası
        \item[E)] Hiçbiri
    \end{itemize}
    \tcblower
    \textbf{Doğru Cevap: B} \\
    \textbf{Açıklama:} Tehlike sınıfı uyuşmazlıklarında yetkili merci Çalışma ve Sosyal Güvenlik Bakanlığı'dır.
\end{testcard}

\begin{testcard}{Soru 36}
    NACE kodu ve tehlike sınıfı tebliğine göre, 'Kimyasal boya ve patlayıcı madde üretimi' hangi tehlike sınıfında yer almaktadır?
    \begin{itemize}
        \item[A)] Az Tehlikeli
        \item[B)] Tehlikeli
        \item[C)] Çok Tehlikeli
        \item[D)] Muaf sınıf
        \item[E)] Hiçbiri
    \end{itemize}
    \tcblower
    \textbf{Doğru Cevap: C} \\
    \textbf{Açıklama:} Kimyasal patlayıcı üretimi yüksek kaza ve yangın/patlama riski nedeniyle Çok Tehlikeli sınıftadır.
\end{testcard}

\begin{testcard}{Soru 37}
    Çalışanın görevli olarak işverenin sağladığı araçla başka bir şubeye gönderilmesi esnasında yolda geçirdiği kaza iş kazası mıdır?
    \begin{itemize}
        \item[A)] Hayır, yolda olduğu için genel kazadır
        \item[B)] Evet, işveren tarafından görevlendirildiği ve işverenin sağladığı taşıtla seyahat ettiği için iş kazasıdır
        \item[C)] Sadece ölümlü olursa iş kazası sayılır
        \item[D)] Sadece mesai saatleri dışındaysa iş kazasıdır
        \item[E)] Hiçbiri
    \end{itemize}
    \tcblower
    \textbf{Doğru Cevap: B} \\
    \textbf{Açıklama:} İşverenin verdiği görev nedeniyle seyahat ederken yaşanan kazalar doğrudan iş kazası kapsamında değerlendirilir.
\end{testcard}

\begin{testcard}{Soru 38}
    İşyerinde çalışanların İSG kuruluna başvurarak tehlikeli durumları bildirme hakkı hakkında hangisi doğrudur?
    \begin{itemize}
        \item[A)] Bildirim yapan çalışanlar cezalandırılır
        \item[B)] Yasal haklarıdır, tehlike tespitlerini kurula yazılı veya sözlü iletebilirler
        \item[C)] Sadece mesai bitiminden sonra bildirim yapabilirler
        \item[D)] Bildirim yapmak yasaktır
        \item[E)] Hiçbiri
    \end{itemize}
    \tcblower
    \textbf{Doğru Cevap: B} \\
    \textbf{Açıklama:} Çalışanların tehlikeleri tespit edip kurula veya yönetime iletmeleri hem yasal hakları hem de İSG katılımcılığı için gerekliliktir.
\end{testcard}

\begin{testcard}{Soru 39}
    İSG kurulunun kurulması için gereken süreklilik şartı (6 aydan fazla süren işler) neyi ifade eder?
    \begin{itemize}
        \item[A)] İş sözleşmesinin süresiz olmasını
        \item[B)] İşyerinde yapılan faaliyetin mevsimlik veya geçici olmayıp, en az 6 ay kesintisiz devam eden işlerden olmasını
        \item[C)] İşçinin deneme süresini
        \item[D)] Şirketin kuruluş tarihini
        \item[E)] Hiçbiri
    \end{itemize}
    \tcblower
    \textbf{Doğru Cevap: B} \\
    \textbf{Açıklama:} Kısa süreli inşaat onarımları veya geçici etkinlikler gibi 6 aydan kısa süren işler kurul kurma zorunluluğundan muaftır.
\end{testcard}

\begin{testcard}{Soru 40}
    Ortam ölçümleri kapsamında işyerinde yapılması gereken toz ve gaz ölçümleri neye göre denetlenir?
    \begin{itemize}
        \item[A)] İşverenin keyfi sınırlarına göre
        \item[B)] İlgili yönetmeliklerde belirtilen yasal maruziyet sınır değerlerine (ESD/OEL) göre
        \item[C)] Müşteri taleplerine göre
        \item[D)] Ölçüm yapılmadan tahminle karar verilir
        \item[E)] Hiçbiri
    \end{itemize}
    \tcblower
    \textbf{Doğru Cevap: B} \\
    \textbf{Açıklama:} Kimyasal maddelerle çalışmalarda ve tozla mücadele yönetmeliklerinde belirtilen eşik sınır değerler (TWA/STEL) aşılmamalıdır.
\end{testcard}

\begin{testcard}{Soru 41}
    İSG kurullarında yedek üyelerin durumu hakkında hangisi doğrudur?
    \begin{itemize}
        \item[A)] Yedek üyelerin kurulda hiçbir hakkı yoktur
        \item[B)] Asıl üyelerin katılamadığı durumlarda onların yerine kurula katılırlar ve oy hakları olur
        \item[C)] Sadece sekretarya yaparlar
        \item[D)] Şirket dışından seçilirler
        \item[E)] Hiçbiri
    \end{itemize}
    \tcblower
    \textbf{Doğru Cevap: B} \\
    \textbf{Açıklama:} Her asıl üyenin bir yedeği seçilir/atanır ve asıl üye bulunmadığında yedek üye onun yetkileriyle toplantıya katılır.
\end{testcard}

\begin{testcard}{Soru 42}
    Yasal olarak kurul kurması gereken bir işveren, alt işveren (taşeron) çalıştırıyorsa kurul ilişkisi nasıl kurulur?
    \begin{itemize}
        \item[A)] Asıl işveren alt işverenin kuruluna müdahale edemez
        \item[B)] Asıl işveren ve alt işveren kurulları arasında koordinasyon sağlanmalı, gerekirse ortak kurul veya temsilci bulundurulmalıdır
        \item[C)] Alt işveren kurul kurmaktan tamamen muaftır
        \item[D)] Sorumluluk sadece alt işverene aittir
        \item[E)] Hiçbiri
    \end{itemize}
    \tcblower
    \textbf{Doğru Cevap: B} \\
    \textbf{Açıklama:} Asıl işveren-alt işveren ilişkisinde koordinasyon asıl işveren tarafından kurulur ve İSG ortak yürütülür.
\end{testcard}

\begin{testcard}{Soru 43}
    İşyerinde tehlike sınıfına göre 'Çok Tehlikeli' olan bir işletmede İSG kurulunun en geç kaç ayda bir toplanması yasal olarak ertelenemez?
    \begin{itemize}
        \item[A)] 1 ay (her ay toplanmak zorundadır, ertelenemez)
        \item[B)] 2 ay
        \item[C)] 3 ay
        \item[D)] 6 ay
        \item[E)] Hiçbiri
    \end{itemize}
    \tcblower
    \textbf{Doğru Cevap: A} \\
    \textbf{Açıklama:} Çok Tehlikeli sınıfta kurul periyodu uzatılamaz; her ay olağan toplantı yapılması zorunludur.
\end{testcard}

\begin{testcard}{Soru 44}
    Sağlık ve güvenlik işaretlerinde mavi renk üzerine beyaz baret sembolü neyi ifade eder?
    \begin{itemize}
        \item[A)] Baret takmanın yasak olduğunu
        \item[B)] Baret takmanın zorunlu olduğunu
        \item[C)] Baretlerin satılık olduğunu
        \item[D)] Yakında şantiye olduğunu
        \item[E)] Hiçbiri
    \end{itemize}
    \tcblower
    \textbf{Doğru Cevap: B} \\
    \textbf{Açıklama:} Mavi zemin emredici/zorunluluk belirtir. Baret resmi baret takılmasını emreder.
\end{testcard}

\begin{testcard}{Soru 45}
    Sağlık ve güvenlik işaretlerinde sarı üçgen içerisindeki kuru kafa sembolü ne anlama gelir?
    \begin{itemize}
        \item[A)] Güvenli alan
        \item[B)] Toksik / zehirli madde tehlikesi uyarısı
        \item[C)] Yangın ekipmanı
        \item[D)] Zorunlu ilk yardım malzemesi
        \item[E)] Hiçbiri
    \end{itemize}
    \tcblower
    \textbf{Doğru Cevap: B} \\
    \textbf{Açıklama:} Sarı üçgen uyarıdır. Kuru kafa ise zehirli/toksik madde tehlikesine karşı uyarır.
\end{testcard}

\begin{testcard}{Soru 46}
    Sağlık ve güvenlik işaretlerinde kırmızı daire içerisinden geçen eğik kırmızı çizgi altındaki sigara sembolü neyi ifade eder?
    \begin{itemize}
        \item[A)] Sigara içmenin serbest olduğunu
        \item[B)] Sigara içilmesinin yasak olduğunu
        \item[C)] Yangın çıkışını
        \item[D)] Sigara içme alanı olduğunu
        \item[E)] Hiçbiri
    \end{itemize}
    \tcblower
    \textbf{Doğru Cevap: B} \\
    \textbf{Açıklama:} Kırmızı çember ve çapraz çizgi yasaklama belirtir. Bu işaret sigara içilmesini yasaklar.
\end{testcard}

\begin{testcard}{Soru 47}
    Sağlık ve güvenlik işaretlerinde yeşil zemin üzerindeki beyaz koşan insan ve kapı sembolü ne anlama gelir?
    \begin{itemize}
        \item[A)] Yasak bölge
        \item[B)] Acil çıkış / Kaçış yolu
        \item[C)] Asansör yeri
        \item[D)] Spor salonu girişi
        \item[E)] Hiçbiri
    \end{itemize}
    \tcblower
    \textbf{Doğru Cevap: B} \\
    \textbf{Açıklama:} Yeşil renk acil kaçış ve kurtarma işaretidir. Kapıya koşan insan acil çıkış yönünü gösterir.
\end{testcard}

\begin{testcard}{Soru 48}
    İşyerinde ortam ölçümleri (gürültü, toz vb.) kimler tarafından yapılmalıdır?
    \begin{itemize}
        \item[A)] Herhangi bir çalışan
        \item[B)] Bakanlıkça yetkilendirilmiş akredite laboratuvarlar
        \item[C)] İşyeri muhasebecisi
        \item[D)] Sadece iş müfettişleri
        \item[E)] Hiçbiri
    \end{itemize}
    \tcblower
    \textbf{Doğru Cevap: B} \\
    \textbf{Açıklama:} Ölçümlerin hukuki geçerliliği olması için İSGÜM yetkili akredite laboratuvarlarca kalibreli cihazlarla yapılması şarttır.
\end{testcard}

\begin{testcard}{Soru 49}
    İSG Kurulları yönetmeliğine göre, kurul üyelerinin oylamada çekimser kalma hakkı var mıdır?
    \begin{itemize}
        \item[A)] Evet, herkes çekimser kalabilir
        \item[B)] Hayır, kurul üyeleri oylamada çekimser oy kullanamazlar (kabul veya ret oyu vermelidirler)
        \item[C)] Sadece başkan çekimser kalabilir
        \item[D)] Sadece hekim çekimser kalabilir
        \item[E)] Hiçbiri
    \end{itemize}
    \tcblower
    \textbf{Doğru Cevap: B} \\
    \textbf{Açıklama:} Kurulda net karar alınabilmesi için üyelerin çekimser oy kullanması yasaktır; kabul ya da ret yönünde oy kullanmalıdırlar.
\end{testcard}

\begin{testcard}{Soru 50}
    İSG Kurul kararları ne zaman yürürlüğe girer?
    \begin{itemize}
        \item[A)] 1 yıl sonra
        \item[B)] Karar defteri imzalanıp taraflara duyurulduğu andan itibaren
        \item[C)] Sadece Bakanlık onayladığında
        \item[D)] Yıl sonunda
        \item[E)] Hiçbiri
    \end{itemize}
    \tcblower
    \textbf{Doğru Cevap: B} \\
    \textbf{Açıklama:} Kararlar imzalanıp ilan edildiği an yürürlüğe girer ve aksiyon süreleri başlar.
\end{testcard}

\begin{testcard}{Soru 51}
    Heinrich'in Domino Etkisi Kuramına göre, kazaları önlemede en zayıf halka olan ve aradan çekildiğinde kazanın meydana gelmesini engelleyen kilit zincir (halka) aşağıdakilerden hangisidir?
    \begin{itemize}
        \item[A)] Kalıtımsal özellikler ve sosyal çevre
        \item[B)] Kişisel kusurlar / İnsan hataları
        \item[C)] Güvensiz Hareketler / Güvensiz Koşullar (Tehlikeli Durum-Davranış)
        \item[D)] Yaralanma ve hasar
        \item[E)] Hiçbiri
    \end{itemize}
    \tcblower
    \textbf{Doğru Cevap: C} \\
    \textbf{Açıklama:} Domino Etkisi Kuramına göre beş domino taşından (çevre, insan hatası, tehlikeli durum/davranış, kaza, yaralanma) üçüncüsü olan 'Güvensiz Hareketler ve Güvensiz Koşullar' zincirden çekilirse kaza meydana gelmez ve zincir kırılır.
\end{testcard}

\begin{testcard}{Soru 52}
    Heinrich'in 1-29-300 Kaza Piramidi modeline göre, iş sağlığı ve güvenliği çalışmalarında ağır yaralanmalı veya ölümcül büyük kazaları (1) engelleyebilmek için asıl odaklanılması gereken alan aşağıdakilerden hangisidir?
    \begin{itemize}
        \item[A)] Sadece ölümlü kazaların hukuki sonuçlarını araştırmak
        \item[B)] 300 adet yaralanmasız (ramak kala/olay) ve 29 adet hafif yaralanmalı olayın kök nedenlerini ortadan kaldırmak
        \item[C)] Kişisel koruyucu donanımların renklerini değiştirmek
        \item[D)] Sadece 1 adet büyük kazanın olduğu güne odaklanmak
        \item[E)] Hiçbiri
    \end{itemize}
    \tcblower
    \textbf{Doğru Cevap: B} \\
    \textbf{Açıklama:} Kaza piramidine göre her 1 ağır kazanın altında 29 hafif yaralanmalı kaza ve 300 yaralanmasız ramak kala olay yatar. Büyük kazaları önlemenin yolu piramidin tabanındaki ramak kala olayları ve küçük kazaları önlemektir.
\end{testcard}

\begin{testcard}{Soru 53}
    5510 Sayılı Kanun Madde 13 kapsamında iş kazası tanımları dikkate alındığında, bir çalışanın öğle yemeği molasında işyeri bahçesinde voleybol oynarken düşüp ayağını kırması olayı hakkında yapılan değerlendirmelerden hangisi doğrudur?
    \begin{itemize}
        \item[A)] İş saatleri dışında olduğu için iş kazası sayılmaz.
        \item[B)] Çalışan kendi rızasıyla oynadığı ve asıl işini yapmadığı için iş kazası sayılmaz.
        \item[C)] Olay işyeri sınırları (avlu/bahçe) içinde meydana geldiği için asıl işini yapıp yapmadığına bakılmaksızın yasal olarak iş kazasıdır.
        \item[D)] Spor faaliyeti sadece Çok Tehlikeli işyerlerinde iş kazası sayılır.
        \item[E)] Hiçbiri
    \end{itemize}
    \tcblower
    \textbf{Doğru Cevap: C} \\
    \textbf{Açıklama:} 5510 sayılı Kanun'a göre, sigortalının işyerinde bulunduğu sırada meydana gelen her kaza, asıl işini yapıp yapmadığına bakılmaksızın yasal olarak iş kazası kabul edilir. İşyeri sınırları içindeki avlu, bahçe ve dinlenme alanları da işyeri sayılır.
\end{testcard}

\begin{testcard}{Soru 54}
    Emziren bir kadın çalışanın, mevzuatın kendisine tanıdığı günlük süt izni saatinde çocuğuna süt vermek amacıyla işyeri dışındaki evine giderken yolda trafik kazası geçirmesi durumunun yasal statüsü aşağıdakilerden hangisidir?
    \begin{itemize}
        \item[A)] Çalışan işyeri dışında ve kendi özel işiyle ilgilendiği için iş kazası sayılmaz.
        \item[B)] 5510 sayılı Kanun uyarınca, emziren kadın sigortalının çocuğuna süt vermek için ayrılan zamanlarda meydana gelen kazalar iş kazası kabul edilir.
        \item[C)] Sadece işverenin tahsis ettiği araçla gidiyorsa iş kazasıdır.
        \item[D)] Yasal mevzuatta bu duruma dair herhangi bir düzenleme bulunmamaktadır.
        \item[E)] Hiçbiri
    \end{itemize}
    \tcblower
    \textbf{Doğru Cevap: B} \\
    \textbf{Açıklama:} 5510 Sayılı Sosyal Sigortalar ve Genel Sağlık Sigortası Kanunu Madde 13/d bendine göre emziren kadın sigortalının çocuğuna süt vermek için ayrılan zamanlarda (süt izninde) geçirdiği kazalar iş kazası sayılır.
\end{testcard}

\begin{testcard}{Soru 55}
    Aşağıdaki durumların hangisinde işe gidiş-dönüş esnasında gerçekleşen bir kaza 5510 sayılı Kanun kapsamında 'İş Kazası' olarak nitelendirilir?
    \begin{itemize}
        \item[A)] Çalışanın kendi özel otomobiliyle sabah işe giderken kaza yapması
        \item[B)] Çalışanın iş çıkışı toplu taşıma (metro/otobüs) kullanırken kaza geçirmesi
        \item[C)] Çalışanın işveren tarafından sağlanan servis aracıyla işin yapıldığı yere toplu olarak götürülüp getirilmesi sırasında servis kazası geçirmesi
        \item[D)] Çalışanın işe gitmek üzere evinden çıkıp otobüs durağına yürürken düşmesi
        \item[E)] Hiçbiri
    \end{itemize}
    \tcblower
    \textbf{Doğru Cevap: C} \\
    \textbf{Açıklama:} 5510 sayılı Kanun'un 13. maddesine göre, sigortalıların işverence sağlanan bir taşıtla işin yapıldığı yere toplu olarak götürülüp getirilmeleri sırasında meydana gelen kazalar iş kazası sayılır. Bireysel ulaşımlar kapsam dışıdır.
\end{testcard}

\subsection{Ünite 4: Çalışmaktan Kaçınma ve Formüller}
\begin{testcard}{Soru 1}
    6331 sayılı İSG Kanunu Madde 13'e göre ciddi ve yakın bir tehlike ile karşılaşan çalışan ne talep edebilir?
    \begin{itemize}
        \item[A)] Ücret artışı
        \item[B)] İSG kuruluna başvurarak durumun tespitini ve gerekli tedbirlerin alınmasını
        \item[C)] İşyeri hekiminin değiştirilmesini
        \item[D)] Şirketin başka bir şehre taşınmasını
        \item[E)] Hiçbiri
    \end{itemize}
    \tcblower
    \textbf{Doğru Cevap: B} \\
    \textbf{Açıklama:} Çalışan, ciddi ve yakın tehlike durumunda kurula (veya kurul yoksa işverene) başvurarak tedbir talep etme hakkına sahiptir.
\end{testcard}

\begin{testcard}{Soru 2}
    Çalışmaktan kaçınma hakkını kullanan çalışanların bu süre zarfındaki hakları hakkında hangisi doğrudur?
    \begin{itemize}
        \item[A)] Ücretleri ve diğer özlük hakları aynen devam eder, sözleşmeleri feshedilemez
        \item[B)] Çalışmadıkları için ücretleri kesilir
        \item[C)] Tazminatsız olarak işten çıkarılırlar
        \item[D)] Yıllık izinlerinden düşülür
        \item[E)] Hiçbiri
    \end{itemize}
    \tcblower
    \textbf{Doğru Cevap: A} \\
    \textbf{Açıklama:} Yasa gereği çalışmaktan kaçınılan süre boyunca çalışanın hakları saklıdır, ücret ve sosyal haklarında kesinti yapılamaz.
\end{testcard}

\begin{testcard}{Soru 3}
    Kaza Sıklık Hızı (KSH) hesaplanırken pay kısmında yer alan kazalar hangi katsayı ile çarpılarak standardize edilir?
    \begin{itemize}
        \item[A)] 1.000
        \item[B)] 10.000
        \item[C)] 100.000
        \item[D)] 1.000.000
        \item[E)] Hiçbiri
    \end{itemize}
    \tcblower
    \textbf{Doğru Cevap: D} \\
    \textbf{Açıklama:} Kaza Sıklık Hızı uluslararası standartlarda 1.000.000 adam-saat katsayısı kullanılarak hesaplanır.
\end{testcard}

\begin{testcard}{Soru 4}
    KSH formülünde kullanılan 1.000.000 katsayısı neyi temsil eder?
    \begin{itemize}
        \item[A)] Fabrikanın cirosunu
        \item[B)] 100 çalışanın ömür boyu çalışma saatini
        \item[C)] Yaklaşık 500 çalışanın 1 yılda yaptığı toplam çalışma adam-saat süresini (referans sanayi ölçeği)
        \item[D)] İş kazası başına kesilen para cezasını
        \item[E)] Hiçbiri
    \end{itemize}
    \tcblower
    \textbf{Doğru Cevap: C} \\
    \textbf{Açıklama:} 1.000.000 katsayısı, 500 tam zamanlı çalışanın haftada 40 saatten 50 haftalık yıllık çalışma süresine denk gelen standart bir kabuldür.
\end{testcard}

\begin{testcard}{Soru 5}
    Kaza Ağırlık Hızı (KAH) hesaplanırken katsayı çarpanı standart olarak kaç kabul edilir?
    \begin{itemize}
        \item[A)] 100
        \item[B)] 1.000 (veya bazı ülkelerde 1.000.000)
        \item[C)] 5.000
        \item[D)] 10.000
        \item[E)] Hiçbiri
    \end{itemize}
    \tcblower
    \textbf{Doğru Cevap: B} \\
    \textbf{Açıklama:} KAH hesaplanırken kayıp gün oranı genellikle 1.000 adam-saat çalışma birimi üzerinden standardize edilir.
\end{testcard}

\begin{testcard}{Soru 6}
    Kaza Ağırlık Hızı (KAH) hesaplamasında, kazanın ölümlü olması veya çalışanın sürekli tam iş göremez hale gelmesi durumunda kayıp gün hanesine standart olarak kaç gün eklenir?
    \begin{itemize}
        \item[A)] 1.000 gün
        \item[B)] 3.000 gün
        \item[C)] 6.000 gün
        \item[D)] 7.500 gün
        \item[E)] Hiçbiri
    \end{itemize}
    \tcblower
    \textbf{Doğru Cevap: C} \\
    \textbf{Açıklama:} Uluslararası standartlarda ölümlü kazalar ve kalıcı maluliyetler için kayıp gün sayısı 6.000 gün kabul edilerek hesaba eklenir.
\end{testcard}

\begin{testcard}{Soru 7}
    Ciddi ve yakın tehlike nedeniyle çalışmaktan kaçınma hakkını kullanan bir çalışana karşı işveren ne yapamaz?
    \begin{itemize}
        \item[A)] İş sözleşmesini haklı nedenle feshedemez veya disiplin cezası uygulayamaz
        \item[B)] Eğitim veremez
        \item[C)] Sağlık raporu isteyemez
        \item[D)] İşyeri kurallarını anlatamaz
        \item[E)] Hiçbiri
    \end{itemize}
    \tcblower
    \textbf{Doğru Cevap: A} \\
    \textbf{Açıklama:} Kaçınma hakkını kullanan çalışanın sözleşmesini feshetmek veya ona yaptırım uygulamak yasal olarak suçtur.
\end{testcard}

\begin{testcard}{Soru 8}
    Kaza Sıklık Hızı (KSH) hesaplamasında payda kısmında yer alan 'Toplam Fiili Çalışma Süresi' neleri kapsamaz?
    \begin{itemize}
        \item[A)] İşçilerin fabrikada fiilen çalışarak geçirdiği süreleri
        \item[B)] Fazla mesai sürelerini
        \item[C)] Hafta sonu tatili, yıllık izin ve resmi tatilde geçen süreleri
        \item[D)] Hazırlık ve bakım çalışmalarında geçen süreleri
        \item[E)] Hiçbiri
    \end{itemize}
    \tcblower
    \textbf{Doğru Cevap: C} \\
    \textbf{Açıklama:} Formülde sadece 'fiilen çalışılan' saatler (mesailer dahil) hesaba katılır; izinli olunan veya çalışılmayan süreler düşülür.
\end{testcard}

\begin{testcard}{Soru 9}
    İşyerinde çalışmaktan kaçınma hakkını kullanan çalışana rağmen işverenin acil tedbir almaması durumunda çalışan ne yapabilir?
    \begin{itemize}
        \item[A)] İş sözleşmesini haklı nedenle derhal feshedebilir
        \item[B)] Gizlice makineleri bozabilir
        \item[C)] Grev başlatabilir
        \item[D)] Hiçbir şey yapamaz, çalışmak zorundadır
        \item[E)] Hiçbiri
    \end{itemize}
    \tcblower
    \textbf{Doğru Cevap: A} \\
    \textbf{Açıklama:} Gerekli tedbirler alınmazsa çalışan iş sözleşmesini haklı nedenle (tazminatlı) derhal feshetme hakkına kavuşur.
\end{testcard}

\begin{testcard}{Soru 10}
    Kaza Ağırlık Hızı (KAH) hesaplamasındaki pay kısmındaki 'Kayıp Gün Sayısı' neleri ifade eder?
    \begin{itemize}
        \item[A)] İşçilerin tatil günlerini
        \item[B)] Yaşanan kazalar nedeniyle hekim tarafından verilen geçici iş göremezlik (istirahat) gün sayılarını
        \item[C)] Şirketin kapalı olduğu günleri
        \item[D)] Malzeme tedarik sürelerini
        \item[E)] Hiçbiri
    \end{itemize}
    \tcblower
    \textbf{Doğru Cevap: B} \\
    \textbf{Açıklama:} Kayıp gün sayısı kazalı çalışanların tedavi ve istirahat nedeniyle işe gelemedikleri yasal raporlu günlerin toplamıdır.
\end{testcard}

\begin{testcard}{Soru 11}
    Kaza Sıklık Hızı (KSH) değeri 15 çıkan bir işletme neyi ifade eder?
    \begin{itemize}
        \item[A)] Çalışılan her 1.000.000 adam-saatte 15 iş kazası yaşandığını
        \item[B)] 15 gün iş kaybı olduğunu
        \item[C)] Şirkette 15 çalışan olduğunu
        \item[D)] Yılda 15 saat çalışıldığını
        \item[E)] Hiçbiri
    \end{itemize}
    \tcblower
    \textbf{Doğru Cevap: A} \\
    \textbf{Açıklama:} KSH katsayısı standardı temsil eder. 15 değeri, 1 milyon çalışma saatinde ortalama 15 kazanın yaşandığı anlamına gelir.
\end{testcard}

\begin{testcard}{Soru 12}
    Çalışmaktan kaçınma hakkı kapsamında kurul kararı beklenmeksizin işin derhal durdurulabileceği durum hangisidir?
    \begin{itemize}
        \item[A)] Küçük bir aydınlatma arızası olduğunda
        \item[B)] Tehlikenin çok ciddi, yakın ve kaçınılmaz olması durumunda (örneğin yangın, çökme riski)
        \item[C)] Yemekhane menüsü beğenilmediğinde
        \item[D)] İş elbisesinin rengi sevilmediğinde
        \item[E)] Hiçbiri
    \end{itemize}
    \tcblower
    \textbf{Doğru Cevap: B} \\
    \textbf{Açıklama:} Tehlikenin kaçınılmaz ve kurul toplantısı beklenemeyecek kadar acil olduğu hallerde çalışan işi derhal bırakıp güvenli yere çekilebilir.
\end{testcard}

\begin{testcard}{Soru 13}
    Kaza Ağırlık Hızı (KAH) değeri 2.5 çıkan bir fabrika ne anlama gelmektedir?
    \begin{itemize}
        \item[A)] Çalışılan her 1.000 saatte 2.5 gün iş gücü kaybı yaşandığını
        \item[B)] Fabrikada 2.5 kaza olduğunu
        \item[C)] İş kazalarının 2.5 saat sürdüğünü
        \item[D)] 2.5 işçinin yaralandığını
        \item[E)] Hiçbiri
    \end{itemize}
    \tcblower
    \textbf{Doğru Cevap: A} \\
    \textbf{Açıklama:} KAH paydasındaki 1000 katsayısı nedeniyle sonuç, 1000 çalışma saatindeki kayıp gün oranını gösterir.
\end{testcard}

\begin{testcard}{Soru 14}
    6331 sayılı Kanun Madde 13 uyarınca kurulun bulunmadığı küçük işyerlerinde çalışmaktan kaçınma hakkı başvurusu nereye yapılır?
    \begin{itemize}
        \item[A)] Doğrudan işverene veya işveren vekiline
        \item[B)] Belediye başkanlığına
        \item[C)] Ticaret odasına
        \item[D)] Karakola
        \item[E)] Hiçbiri
    \end{itemize}
    \tcblower
    \textbf{Doğru Cevap: A} \\
    \textbf{Açıklama:} Kurul kurulması zorunlu olmayan (50'den az çalışanlı) işyerlerinde başvuru doğrudan işverene iletilir.
\end{testcard}

\begin{testcard}{Soru 15}
    İSG istatistiklerinde kullanılan KSH ve KAH değerlerinin hesaplanmasındaki temel amaç nedir?
    \begin{itemize}
        \item[A)] Şirketlerin vergi matrahlarını düşürmek
        \item[B)] İşyerinin İSG performansını zaman içinde izlemek ve diğer işletmelerle karşılaştırabilmek
        \item[C)] Çalışanlara ceza kesmek
        \item[D)] Sigorta şirketlerine daha fazla prim ödemek
        \item[E)] Hiçbiri
    \end{itemize}
    \tcblower
    \textbf{Doğru Cevap: B} \\
    \textbf{Açıklama:} Standart katsayılarla yapılan hesaplamalar, farklı boyutlardaki fabrikaların kaza oranlarını adil şekilde karşılaştırmayı sağlar.
\end{testcard}

\begin{testcard}{Soru 16}
    Kaza Ağırlık Hızında (KAH) ölümlü kazalar için eklenen 6000 kayıp gün sayısı hangi esasa dayanır?
    \begin{itemize}
        \item[A)] İşçinin ortalama kalan çalışma ömrünün (yıl bazında) yasal gün karşılığı kabulüdür
        \item[B)] Rastgele seçilmiş bir sayıdır
        \item[C)] İş Kanunu ceza gün sınırıdır
        \item[D)] Yıllık izin günlerinin toplamıdır
        \item[E)] Hiçbiri
    \end{itemize}
    \tcblower
    \textbf{Doğru Cevap: A} \\
    \textbf{Açıklama:} Uluslararası Çalışma Örgütü (ILO) kabullerinde, bir işçinin ölümlü kazayla kaybolan iş gücü potansiyeli 20 yıllık çalışma ömrüne (6000 gün) eşitlenmiştir.
\end{testcard}

\begin{testcard}{Soru 17}
    İSG Kanunu Madde 13'e göre, çalışmaktan kaçınma hakkını kullanan bir işçi, işverenin gerekli tedbirleri almaması halinde iş sözleşmesini nasıl sonlandırabilir? (Örnek No: 1)
    \begin{itemize}
        \item[A)] Haklı nedenle derhal fesih hakkını kullanarak (tazminat hakkı saklı kalacak şekilde)
        \item[B)] İstifa ederek tazminat hakkından vazgeçerek
        \item[C)] Hiçbir şekilde sonlandıramaz
        \item[D)] Sadece arabulucu kararıyla
        \item[E)] Hiçbiri
    \end{itemize}
    \tcblower
    \textbf{Doğru Cevap: A} \\
    \textbf{Açıklama:} İşverenin tedbir almama ısrarı durumunda çalışan sözleşmesini haklı nedenle tazminatlı olarak feshedebilir.
\end{testcard}

\begin{testcard}{Soru 18}
    İSG Kanunu Madde 13'e göre, çalışmaktan kaçınma hakkını kullanan bir işçi, işverenin gerekli tedbirleri almaması halinde iş sözleşmesini nasıl sonlandırabilir? (Örnek No: 2)
    \begin{itemize}
        \item[A)] Haklı nedenle derhal fesih hakkını kullanarak (tazminat hakkı saklı kalacak şekilde)
        \item[B)] İstifa ederek tazminat hakkından vazgeçerek
        \item[C)] Hiçbir şekilde sonlandıramaz
        \item[D)] Sadece arabulucu kararıyla
        \item[E)] Hiçbiri
    \end{itemize}
    \tcblower
    \textbf{Doğru Cevap: A} \\
    \textbf{Açıklama:} İşverenin tedbir almama ısrarı durumunda çalışan sözleşmesini haklı nedenle tazminatlı olarak feshedebilir.
\end{testcard}

\begin{testcard}{Soru 19}
    İSG Kanunu Madde 13'e göre, çalışmaktan kaçınma hakkını kullanan bir işçi, işverenin gerekli tedbirleri almaması halinde iş sözleşmesini nasıl sonlandırabilir? (Örnek No: 3)
    \begin{itemize}
        \item[A)] Haklı nedenle derhal fesih hakkını kullanarak (tazminat hakkı saklı kalacak şekilde)
        \item[B)] İstifa ederek tazminat hakkından vazgeçerek
        \item[C)] Hiçbir şekilde sonlandıramaz
        \item[D)] Sadece arabulucu kararıyla
        \item[E)] Hiçbiri
    \end{itemize}
    \tcblower
    \textbf{Doğru Cevap: A} \\
    \textbf{Açıklama:} İşverenin tedbir almama ısrarı durumunda çalışan sözleşmesini haklı nedenle tazminatlı olarak feshedebilir.
\end{testcard}

\begin{testcard}{Soru 20}
    İSG Kanunu Madde 13'e göre, çalışmaktan kaçınma hakkını kullanan bir işçi, işverenin gerekli tedbirleri almaması halinde iş sözleşmesini nasıl sonlandırabilir? (Örnek No: 4)
    \begin{itemize}
        \item[A)] Haklı nedenle derhal fesih hakkını kullanarak (tazminat hakkı saklı kalacak şekilde)
        \item[B)] İstifa ederek tazminat hakkından vazgeçerek
        \item[C)] Hiçbir şekilde sonlandıramaz
        \item[D)] Sadece arabulucu kararıyla
        \item[E)] Hiçbiri
    \end{itemize}
    \tcblower
    \textbf{Doğru Cevap: A} \\
    \textbf{Açıklama:} İşverenin tedbir almama ısrarı durumunda çalışan sözleşmesini haklı nedenle tazminatlı olarak feshedebilir.
\end{testcard}

\begin{testcard}{Soru 21}
    İSG Kanunu Madde 13'e göre, çalışmaktan kaçınma hakkını kullanan bir işçi, işverenin gerekli tedbirleri almaması halinde iş sözleşmesini nasıl sonlandırabilir? (Örnek No: 5)
    \begin{itemize}
        \item[A)] Haklı nedenle derhal fesih hakkını kullanarak (tazminat hakkı saklı kalacak şekilde)
        \item[B)] İstifa ederek tazminat hakkından vazgeçerek
        \item[C)] Hiçbir şekilde sonlandıramaz
        \item[D)] Sadece arabulucu kararıyla
        \item[E)] Hiçbiri
    \end{itemize}
    \tcblower
    \textbf{Doğru Cevap: A} \\
    \textbf{Açıklama:} İşverenin tedbir almama ısrarı durumunda çalışan sözleşmesini haklı nedenle tazminatlı olarak feshedebilir.
\end{testcard}

\begin{testcard}{Soru 22}
    İSG Kanunu Madde 13'e göre, çalışmaktan kaçınma hakkını kullanan bir işçi, işverenin gerekli tedbirleri almaması halinde iş sözleşmesini nasıl sonlandırabilir? (Örnek No: 6)
    \begin{itemize}
        \item[A)] Haklı nedenle derhal fesih hakkını kullanarak (tazminat hakkı saklı kalacak şekilde)
        \item[B)] İstifa ederek tazminat hakkından vazgeçerek
        \item[C)] Hiçbir şekilde sonlandıramaz
        \item[D)] Sadece arabulucu kararıyla
        \item[E)] Hiçbiri
    \end{itemize}
    \tcblower
    \textbf{Doğru Cevap: A} \\
    \textbf{Açıklama:} İşverenin tedbir almama ısrarı durumunda çalışan sözleşmesini haklı nedenle tazminatlı olarak feshedebilir.
\end{testcard}

\begin{testcard}{Soru 23}
    İSG Kanunu Madde 13'e göre, çalışmaktan kaçınma hakkını kullanan bir işçi, işverenin gerekli tedbirleri almaması halinde iş sözleşmesini nasıl sonlandırabilir? (Örnek No: 7)
    \begin{itemize}
        \item[A)] Haklı nedenle derhal fesih hakkını kullanarak (tazminat hakkı saklı kalacak şekilde)
        \item[B)] İstifa ederek tazminat hakkından vazgeçerek
        \item[C)] Hiçbir şekilde sonlandıramaz
        \item[D)] Sadece arabulucu kararıyla
        \item[E)] Hiçbiri
    \end{itemize}
    \tcblower
    \textbf{Doğru Cevap: A} \\
    \textbf{Açıklama:} İşverenin tedbir almama ısrarı durumunda çalışan sözleşmesini haklı nedenle tazminatlı olarak feshedebilir.
\end{testcard}

\begin{testcard}{Soru 24}
    İSG Kanunu Madde 13'e göre, çalışmaktan kaçınma hakkını kullanan bir işçi, işverenin gerekli tedbirleri almaması halinde iş sözleşmesini nasıl sonlandırabilir? (Örnek No: 8)
    \begin{itemize}
        \item[A)] Haklı nedenle derhal fesih hakkını kullanarak (tazminat hakkı saklı kalacak şekilde)
        \item[B)] İstifa ederek tazminat hakkından vazgeçerek
        \item[C)] Hiçbir şekilde sonlandıramaz
        \item[D)] Sadece arabulucu kararıyla
        \item[E)] Hiçbiri
    \end{itemize}
    \tcblower
    \textbf{Doğru Cevap: A} \\
    \textbf{Açıklama:} İşverenin tedbir almama ısrarı durumunda çalışan sözleşmesini haklı nedenle tazminatlı olarak feshedebilir.
\end{testcard}

\begin{testcard}{Soru 25}
    İSG Kanunu Madde 13'e göre, çalışmaktan kaçınma hakkını kullanan bir işçi, işverenin gerekli tedbirleri almaması halinde iş sözleşmesini nasıl sonlandırabilir? (Örnek No: 9)
    \begin{itemize}
        \item[A)] Haklı nedenle derhal fesih hakkını kullanarak (tazminat hakkı saklı kalacak şekilde)
        \item[B)] İstifa ederek tazminat hakkından vazgeçerek
        \item[C)] Hiçbir şekilde sonlandıramaz
        \item[D)] Sadece arabulucu kararıyla
        \item[E)] Hiçbiri
    \end{itemize}
    \tcblower
    \textbf{Doğru Cevap: A} \\
    \textbf{Açıklama:} İşverenin tedbir almama ısrarı durumunda çalışan sözleşmesini haklı nedenle tazminatlı olarak feshedebilir.
\end{testcard}

\begin{testcard}{Soru 26}
    İSG Kanunu Madde 13'e göre, çalışmaktan kaçınma hakkını kullanan bir işçi, işverenin gerekli tedbirleri almaması halinde iş sözleşmesini nasıl sonlandırabilir? (Örnek No: 10)
    \begin{itemize}
        \item[A)] Haklı nedenle derhal fesih hakkını kullanarak (tazminat hakkı saklı kalacak şekilde)
        \item[B)] İstifa ederek tazminat hakkından vazgeçerek
        \item[C)] Hiçbir şekilde sonlandıramaz
        \item[D)] Sadece arabulucu kararıyla
        \item[E)] Hiçbiri
    \end{itemize}
    \tcblower
    \textbf{Doğru Cevap: A} \\
    \textbf{Açıklama:} İşverenin tedbir almama ısrarı durumunda çalışan sözleşmesini haklı nedenle tazminatlı olarak feshedebilir.
\end{testcard}

\begin{testcard}{Soru 27}
    İSG Kanunu Madde 13'e göre, çalışmaktan kaçınma hakkını kullanan bir işçi, işverenin gerekli tedbirleri almaması halinde iş sözleşmesini nasıl sonlandırabilir? (Örnek No: 11)
    \begin{itemize}
        \item[A)] Haklı nedenle derhal fesih hakkını kullanarak (tazminat hakkı saklı kalacak şekilde)
        \item[B)] İstifa ederek tazminat hakkından vazgeçerek
        \item[C)] Hiçbir şekilde sonlandıramaz
        \item[D)] Sadece arabulucu kararıyla
        \item[E)] Hiçbiri
    \end{itemize}
    \tcblower
    \textbf{Doğru Cevap: A} \\
    \textbf{Açıklama:} İşverenin tedbir almama ısrarı durumunda çalışan sözleşmesini haklı nedenle tazminatlı olarak feshedebilir.
\end{testcard}

\begin{testcard}{Soru 28}
    İSG Kanunu Madde 13'e göre, çalışmaktan kaçınma hakkını kullanan bir işçi, işverenin gerekli tedbirleri almaması halinde iş sözleşmesini nasıl sonlandırabilir? (Örnek No: 12)
    \begin{itemize}
        \item[A)] Haklı nedenle derhal fesih hakkını kullanarak (tazminat hakkı saklı kalacak şekilde)
        \item[B)] İstifa ederek tazminat hakkından vazgeçerek
        \item[C)] Hiçbir şekilde sonlandıramaz
        \item[D)] Sadece arabulucu kararıyla
        \item[E)] Hiçbiri
    \end{itemize}
    \tcblower
    \textbf{Doğru Cevap: A} \\
    \textbf{Açıklama:} İşverenin tedbir almama ısrarı durumunda çalışan sözleşmesini haklı nedenle tazminatlı olarak feshedebilir.
\end{testcard}

\begin{testcard}{Soru 29}
    İSG Kanunu Madde 13'e göre, çalışmaktan kaçınma hakkını kullanan bir işçi, işverenin gerekli tedbirleri almaması halinde iş sözleşmesini nasıl sonlandırabilir? (Örnek No: 13)
    \begin{itemize}
        \item[A)] Haklı nedenle derhal fesih hakkını kullanarak (tazminat hakkı saklı kalacak şekilde)
        \item[B)] İstifa ederek tazminat hakkından vazgeçerek
        \item[C)] Hiçbir şekilde sonlandıramaz
        \item[D)] Sadece arabulucu kararıyla
        \item[E)] Hiçbiri
    \end{itemize}
    \tcblower
    \textbf{Doğru Cevap: A} \\
    \textbf{Açıklama:} İşverenin tedbir almama ısrarı durumunda çalışan sözleşmesini haklı nedenle tazminatlı olarak feshedebilir.
\end{testcard}

\begin{testcard}{Soru 30}
    İSG Kanunu Madde 13'e göre, çalışmaktan kaçınma hakkını kullanan bir işçi, işverenin gerekli tedbirleri almaması halinde iş sözleşmesini nasıl sonlandırabilir? (Örnek No: 14)
    \begin{itemize}
        \item[A)] Haklı nedenle derhal fesih hakkını kullanarak (tazminat hakkı saklı kalacak şekilde)
        \item[B)] İstifa ederek tazminat hakkından vazgeçerek
        \item[C)] Hiçbir şekilde sonlandıramaz
        \item[D)] Sadece arabulucu kararıyla
        \item[E)] Hiçbiri
    \end{itemize}
    \tcblower
    \textbf{Doğru Cevap: A} \\
    \textbf{Açıklama:} İşverenin tedbir almama ısrarı durumunda çalışan sözleşmesini haklı nedenle tazminatlı olarak feshedebilir.
\end{testcard}

\begin{testcard}{Soru 31}
    İSG Kanunu Madde 13'e göre, çalışmaktan kaçınma hakkını kullanan bir işçi, işverenin gerekli tedbirleri almaması halinde iş sözleşmesini nasıl sonlandırabilir? (Örnek No: 15)
    \begin{itemize}
        \item[A)] Haklı nedenle derhal fesih hakkını kullanarak (tazminat hakkı saklı kalacak şekilde)
        \item[B)] İstifa ederek tazminat hakkından vazgeçerek
        \item[C)] Hiçbir şekilde sonlandıramaz
        \item[D)] Sadece arabulucu kararıyla
        \item[E)] Hiçbiri
    \end{itemize}
    \tcblower
    \textbf{Doğru Cevap: A} \\
    \textbf{Açıklama:} İşverenin tedbir almama ısrarı durumunda çalışan sözleşmesini haklı nedenle tazminatlı olarak feshedebilir.
\end{testcard}

\begin{testcard}{Soru 32}
    İSG Kanunu Madde 13'e göre, çalışmaktan kaçınma hakkını kullanan bir işçi, işverenin gerekli tedbirleri almaması halinde iş sözleşmesini nasıl sonlandırabilir? (Örnek No: 16)
    \begin{itemize}
        \item[A)] Haklı nedenle derhal fesih hakkını kullanarak (tazminat hakkı saklı kalacak şekilde)
        \item[B)] İstifa ederek tazminat hakkından vazgeçerek
        \item[C)] Hiçbir şekilde sonlandıramaz
        \item[D)] Sadece arabulucu kararıyla
        \item[E)] Hiçbiri
    \end{itemize}
    \tcblower
    \textbf{Doğru Cevap: A} \\
    \textbf{Açıklama:} İşverenin tedbir almama ısrarı durumunda çalışan sözleşmesini haklı nedenle tazminatlı olarak feshedebilir.
\end{testcard}

\begin{testcard}{Soru 33}
    İSG Kanunu Madde 13'e göre, çalışmaktan kaçınma hakkını kullanan bir işçi, işverenin gerekli tedbirleri almaması halinde iş sözleşmesini nasıl sonlandırabilir? (Örnek No: 17)
    \begin{itemize}
        \item[A)] Haklı nedenle derhal fesih hakkını kullanarak (tazminat hakkı saklı kalacak şekilde)
        \item[B)] İstifa ederek tazminat hakkından vazgeçerek
        \item[C)] Hiçbir şekilde sonlandıramaz
        \item[D)] Sadece arabulucu kararıyla
        \item[E)] Hiçbiri
    \end{itemize}
    \tcblower
    \textbf{Doğru Cevap: A} \\
    \textbf{Açıklama:} İşverenin tedbir almama ısrarı durumunda çalışan sözleşmesini haklı nedenle tazminatlı olarak feshedebilir.
\end{testcard}

\begin{testcard}{Soru 34}
    İSG Kanunu Madde 13'e göre, çalışmaktan kaçınma hakkını kullanan bir işçi, işverenin gerekli tedbirleri almaması halinde iş sözleşmesini nasıl sonlandırabilir? (Örnek No: 18)
    \begin{itemize}
        \item[A)] Haklı nedenle derhal fesih hakkını kullanarak (tazminat hakkı saklı kalacak şekilde)
        \item[B)] İstifa ederek tazminat hakkından vazgeçerek
        \item[C)] Hiçbir şekilde sonlandıramaz
        \item[D)] Sadece arabulucu kararıyla
        \item[E)] Hiçbiri
    \end{itemize}
    \tcblower
    \textbf{Doğru Cevap: A} \\
    \textbf{Açıklama:} İşverenin tedbir almama ısrarı durumunda çalışan sözleşmesini haklı nedenle tazminatlı olarak feshedebilir.
\end{testcard}

\begin{testcard}{Soru 35}
    İSG Kanunu Madde 13'e göre, çalışmaktan kaçınma hakkını kullanan bir işçi, işverenin gerekli tedbirleri almaması halinde iş sözleşmesini nasıl sonlandırabilir? (Örnek No: 19)
    \begin{itemize}
        \item[A)] Haklı nedenle derhal fesih hakkını kullanarak (tazminat hakkı saklı kalacak şekilde)
        \item[B)] İstifa ederek tazminat hakkından vazgeçerek
        \item[C)] Hiçbir şekilde sonlandıramaz
        \item[D)] Sadece arabulucu kararıyla
        \item[E)] Hiçbiri
    \end{itemize}
    \tcblower
    \textbf{Doğru Cevap: A} \\
    \textbf{Açıklama:} İşverenin tedbir almama ısrarı durumunda çalışan sözleşmesini haklı nedenle tazminatlı olarak feshedebilir.
\end{testcard}

\begin{testcard}{Soru 36}
    İSG Kanunu Madde 13'e göre, çalışmaktan kaçınma hakkını kullanan bir işçi, işverenin gerekli tedbirleri almaması halinde iş sözleşmesini nasıl sonlandırabilir? (Örnek No: 20)
    \begin{itemize}
        \item[A)] Haklı nedenle derhal fesih hakkını kullanarak (tazminat hakkı saklı kalacak şekilde)
        \item[B)] İstifa ederek tazminat hakkından vazgeçerek
        \item[C)] Hiçbir şekilde sonlandıramaz
        \item[D)] Sadece arabulucu kararıyla
        \item[E)] Hiçbiri
    \end{itemize}
    \tcblower
    \textbf{Doğru Cevap: A} \\
    \textbf{Açıklama:} İşverenin tedbir almama ısrarı durumunda çalışan sözleşmesini haklı nedenle tazminatlı olarak feshedebilir.
\end{testcard}

\begin{testcard}{Soru 37}
    İSG Kanunu Madde 13'e göre, çalışmaktan kaçınma hakkını kullanan bir işçi, işverenin gerekli tedbirleri almaması halinde iş sözleşmesini nasıl sonlandırabilir? (Örnek No: 21)
    \begin{itemize}
        \item[A)] Haklı nedenle derhal fesih hakkını kullanarak (tazminat hakkı saklı kalacak şekilde)
        \item[B)] İstifa ederek tazminat hakkından vazgeçerek
        \item[C)] Hiçbir şekilde sonlandıramaz
        \item[D)] Sadece arabulucu kararıyla
        \item[E)] Hiçbiri
    \end{itemize}
    \tcblower
    \textbf{Doğru Cevap: A} \\
    \textbf{Açıklama:} İşverenin tedbir almama ısrarı durumunda çalışan sözleşmesini haklı nedenle tazminatlı olarak feshedebilir.
\end{testcard}

\begin{testcard}{Soru 38}
    İSG Kanunu Madde 13'e göre, çalışmaktan kaçınma hakkını kullanan bir işçi, işverenin gerekli tedbirleri almaması halinde iş sözleşmesini nasıl sonlandırabilir? (Örnek No: 22)
    \begin{itemize}
        \item[A)] Haklı nedenle derhal fesih hakkını kullanarak (tazminat hakkı saklı kalacak şekilde)
        \item[B)] İstifa ederek tazminat hakkından vazgeçerek
        \item[C)] Hiçbir şekilde sonlandıramaz
        \item[D)] Sadece arabulucu kararıyla
        \item[E)] Hiçbiri
    \end{itemize}
    \tcblower
    \textbf{Doğru Cevap: A} \\
    \textbf{Açıklama:} İşverenin tedbir almama ısrarı durumunda çalışan sözleşmesini haklı nedenle tazminatlı olarak feshedebilir.
\end{testcard}

\begin{testcard}{Soru 39}
    İSG Kanunu Madde 13'e göre, çalışmaktan kaçınma hakkını kullanan bir işçi, işverenin gerekli tedbirleri almaması halinde iş sözleşmesini nasıl sonlandırabilir? (Örnek No: 23)
    \begin{itemize}
        \item[A)] Haklı nedenle derhal fesih hakkını kullanarak (tazminat hakkı saklı kalacak şekilde)
        \item[B)] İstifa ederek tazminat hakkından vazgeçerek
        \item[C)] Hiçbir şekilde sonlandıramaz
        \item[D)] Sadece arabulucu kararıyla
        \item[E)] Hiçbiri
    \end{itemize}
    \tcblower
    \textbf{Doğru Cevap: A} \\
    \textbf{Açıklama:} İşverenin tedbir almama ısrarı durumunda çalışan sözleşmesini haklı nedenle tazminatlı olarak feshedebilir.
\end{testcard}

\begin{testcard}{Soru 40}
    İSG Kanunu Madde 13'e göre, çalışmaktan kaçınma hakkını kullanan bir işçi, işverenin gerekli tedbirleri almaması halinde iş sözleşmesini nasıl sonlandırabilir? (Örnek No: 24)
    \begin{itemize}
        \item[A)] Haklı nedenle derhal fesih hakkını kullanarak (tazminat hakkı saklı kalacak şekilde)
        \item[B)] İstifa ederek tazminat hakkından vazgeçerek
        \item[C)] Hiçbir şekilde sonlandıramaz
        \item[D)] Sadece arabulucu kararıyla
        \item[E)] Hiçbiri
    \end{itemize}
    \tcblower
    \textbf{Doğru Cevap: A} \\
    \textbf{Açıklama:} İşverenin tedbir almama ısrarı durumunda çalışan sözleşmesini haklı nedenle tazminatlı olarak feshedebilir.
\end{testcard}

\begin{testcard}{Soru 41}
    İSG Kanunu Madde 13'e göre, çalışmaktan kaçınma hakkını kullanan bir işçi, işverenin gerekli tedbirleri almaması halinde iş sözleşmesini nasıl sonlandırabilir? (Örnek No: 25)
    \begin{itemize}
        \item[A)] Haklı nedenle derhal fesih hakkını kullanarak (tazminat hakkı saklı kalacak şekilde)
        \item[B)] İstifa ederek tazminat hakkından vazgeçerek
        \item[C)] Hiçbir şekilde sonlandıramaz
        \item[D)] Sadece arabulucu kararıyla
        \item[E)] Hiçbiri
    \end{itemize}
    \tcblower
    \textbf{Doğru Cevap: A} \\
    \textbf{Açıklama:} İşverenin tedbir almama ısrarı durumunda çalışan sözleşmesini haklı nedenle tazminatlı olarak feshedebilir.
\end{testcard}

\begin{testcard}{Soru 42}
    İSG Kanunu Madde 13'e göre, çalışmaktan kaçınma hakkını kullanan bir işçi, işverenin gerekli tedbirleri almaması halinde iş sözleşmesini nasıl sonlandırabilir? (Örnek No: 26)
    \begin{itemize}
        \item[A)] Haklı nedenle derhal fesih hakkını kullanarak (tazminat hakkı saklı kalacak şekilde)
        \item[B)] İstifa ederek tazminat hakkından vazgeçerek
        \item[C)] Hiçbir şekilde sonlandıramaz
        \item[D)] Sadece arabulucu kararıyla
        \item[E)] Hiçbiri
    \end{itemize}
    \tcblower
    \textbf{Doğru Cevap: A} \\
    \textbf{Açıklama:} İşverenin tedbir almama ısrarı durumunda çalışan sözleşmesini haklı nedenle tazminatlı olarak feshedebilir.
\end{testcard}

\begin{testcard}{Soru 43}
    İSG Kanunu Madde 13'e göre, çalışmaktan kaçınma hakkını kullanan bir işçi, işverenin gerekli tedbirleri almaması halinde iş sözleşmesini nasıl sonlandırabilir? (Örnek No: 27)
    \begin{itemize}
        \item[A)] Haklı nedenle derhal fesih hakkını kullanarak (tazminat hakkı saklı kalacak şekilde)
        \item[B)] İstifa ederek tazminat hakkından vazgeçerek
        \item[C)] Hiçbir şekilde sonlandıramaz
        \item[D)] Sadece arabulucu kararıyla
        \item[E)] Hiçbiri
    \end{itemize}
    \tcblower
    \textbf{Doğru Cevap: A} \\
    \textbf{Açıklama:} İşverenin tedbir almama ısrarı durumunda çalışan sözleşmesini haklı nedenle tazminatlı olarak feshedebilir.
\end{testcard}

\begin{testcard}{Soru 44}
    İSG Kanunu Madde 13'e göre, çalışmaktan kaçınma hakkını kullanan bir işçi, işverenin gerekli tedbirleri almaması halinde iş sözleşmesini nasıl sonlandırabilir? (Örnek No: 28)
    \begin{itemize}
        \item[A)] Haklı nedenle derhal fesih hakkını kullanarak (tazminat hakkı saklı kalacak şekilde)
        \item[B)] İstifa ederek tazminat hakkından vazgeçerek
        \item[C)] Hiçbir şekilde sonlandıramaz
        \item[D)] Sadece arabulucu kararıyla
        \item[E)] Hiçbiri
    \end{itemize}
    \tcblower
    \textbf{Doğru Cevap: A} \\
    \textbf{Açıklama:} İşverenin tedbir almama ısrarı durumunda çalışan sözleşmesini haklı nedenle tazminatlı olarak feshedebilir.
\end{testcard}

\begin{testcard}{Soru 45}
    İSG Kanunu Madde 13'e göre, çalışmaktan kaçınma hakkını kullanan bir işçi, işverenin gerekli tedbirleri almaması halinde iş sözleşmesini nasıl sonlandırabilir? (Örnek No: 29)
    \begin{itemize}
        \item[A)] Haklı nedenle derhal fesih hakkını kullanarak (tazminat hakkı saklı kalacak şekilde)
        \item[B)] İstifa ederek tazminat hakkından vazgeçerek
        \item[C)] Hiçbir şekilde sonlandıramaz
        \item[D)] Sadece arabulucu kararıyla
        \item[E)] Hiçbiri
    \end{itemize}
    \tcblower
    \textbf{Doğru Cevap: A} \\
    \textbf{Açıklama:} İşverenin tedbir almama ısrarı durumunda çalışan sözleşmesini haklı nedenle tazminatlı olarak feshedebilir.
\end{testcard}

\begin{testcard}{Soru 46}
    İSG Kanunu Madde 13'e göre, çalışmaktan kaçınma hakkını kullanan bir işçi, işverenin gerekli tedbirleri almaması halinde iş sözleşmesini nasıl sonlandırabilir? (Örnek No: 30)
    \begin{itemize}
        \item[A)] Haklı nedenle derhal fesih hakkını kullanarak (tazminat hakkı saklı kalacak şekilde)
        \item[B)] İstifa ederek tazminat hakkından vazgeçerek
        \item[C)] Hiçbir şekilde sonlandıramaz
        \item[D)] Sadece arabulucu kararıyla
        \item[E)] Hiçbiri
    \end{itemize}
    \tcblower
    \textbf{Doğru Cevap: A} \\
    \textbf{Açıklama:} İşverenin tedbir almama ısrarı durumunda çalışan sözleşmesini haklı nedenle tazminatlı olarak feshedebilir.
\end{testcard}

\begin{testcard}{Soru 47}
    İSG Kanunu Madde 13'e göre, çalışmaktan kaçınma hakkını kullanan bir işçi, işverenin gerekli tedbirleri almaması halinde iş sözleşmesini nasıl sonlandırabilir? (Örnek No: 31)
    \begin{itemize}
        \item[A)] Haklı nedenle derhal fesih hakkını kullanarak (tazminat hakkı saklı kalacak şekilde)
        \item[B)] İstifa ederek tazminat hakkından vazgeçerek
        \item[C)] Hiçbir şekilde sonlandıramaz
        \item[D)] Sadece arabulucu kararıyla
        \item[E)] Hiçbiri
    \end{itemize}
    \tcblower
    \textbf{Doğru Cevap: A} \\
    \textbf{Açıklama:} İşverenin tedbir almama ısrarı durumunda çalışan sözleşmesini haklı nedenle tazminatlı olarak feshedebilir.
\end{testcard}

\begin{testcard}{Soru 48}
    İSG Kanunu Madde 13'e göre, çalışmaktan kaçınma hakkını kullanan bir işçi, işverenin gerekli tedbirleri almaması halinde iş sözleşmesini nasıl sonlandırabilir? (Örnek No: 32)
    \begin{itemize}
        \item[A)] Haklı nedenle derhal fesih hakkını kullanarak (tazminat hakkı saklı kalacak şekilde)
        \item[B)] İstifa ederek tazminat hakkından vazgeçerek
        \item[C)] Hiçbir şekilde sonlandıramaz
        \item[D)] Sadece arabulucu kararıyla
        \item[E)] Hiçbiri
    \end{itemize}
    \tcblower
    \textbf{Doğru Cevap: A} \\
    \textbf{Açıklama:} İşverenin tedbir almama ısrarı durumunda çalışan sözleşmesini haklı nedenle tazminatlı olarak feshedebilir.
\end{testcard}

\begin{testcard}{Soru 49}
    İSG Kanunu Madde 13'e göre, çalışmaktan kaçınma hakkını kullanan bir işçi, işverenin gerekli tedbirleri almaması halinde iş sözleşmesini nasıl sonlandırabilir? (Örnek No: 33)
    \begin{itemize}
        \item[A)] Haklı nedenle derhal fesih hakkını kullanarak (tazminat hakkı saklı kalacak şekilde)
        \item[B)] İstifa ederek tazminat hakkından vazgeçerek
        \item[C)] Hiçbir şekilde sonlandıramaz
        \item[D)] Sadece arabulucu kararıyla
        \item[E)] Hiçbiri
    \end{itemize}
    \tcblower
    \textbf{Doğru Cevap: A} \\
    \textbf{Açıklama:} İşverenin tedbir almama ısrarı durumunda çalışan sözleşmesini haklı nedenle tazminatlı olarak feshedebilir.
\end{testcard}

\begin{testcard}{Soru 50}
    İSG Kanunu Madde 13'e göre, çalışmaktan kaçınma hakkını kullanan bir işçi, işverenin gerekli tedbirleri almaması halinde iş sözleşmesini nasıl sonlandırabilir? (Örnek No: 34)
    \begin{itemize}
        \item[A)] Haklı nedenle derhal fesih hakkını kullanarak (tazminat hakkı saklı kalacak şekilde)
        \item[B)] İstifa ederek tazminat hakkından vazgeçerek
        \item[C)] Hiçbir şekilde sonlandıramaz
        \item[D)] Sadece arabulucu kararıyla
        \item[E)] Hiçbiri
    \end{itemize}
    \tcblower
    \textbf{Doğru Cevap: A} \\
    \textbf{Açıklama:} İşverenin tedbir almama ısrarı durumunda çalışan sözleşmesini haklı nedenle tazminatlı olarak feshedebilir.
\end{testcard}

\subsection{Ünite 5: Risk Yönetimi ve Kritik Tablolar}
\begin{testcard}{Soru 1}
    Sınav kağıdında verilecek olan Sayfa 135 Olasılık Tablosuna göre: Bir tehlikenin yılda bir kez veya daha az gerçekleşme olasılığı tablodaki hangi dereceye karşılık gelir?
    \begin{itemize}
        \item[A)] Çok Küçük (1) / Yılda bir veya daha az
        \item[B)] Küçük (2) / Ayda bir veya daha az
        \item[C)] Orta (3) / Haftada bir veya daha az
        \item[D)] Yüksek (4) / Her gün
        \item[E)] Hiçbiri
    \end{itemize}
    \tcblower
    \textbf{Doğru Cevap: A} \\
    \textbf{Açıklama:} Sayfa 135 Olasılık Tablosunda: Çok Küçük (1) derecesi 'yılda bir veya daha az' olasılığına karşılık gelir.
\end{testcard}

\begin{testcard}{Soru 2}
    Sınav kağıdında verilecek olan Sayfa 136 Şiddet (Derecelendirme) Tablosuna göre: 'Ölüm veya sürekli tam iş göremezlik' durumu hangi şiddet derecesine aittir?
    \begin{itemize}
        \item[A)] Hafif (2)
        \item[B)] Orta (3)
        \item[C)] Ciddi (4)
        \item[D)] Çok Ciddi (5)
        \item[E)] Hiçbiri
    \end{itemize}
    \tcblower
    \textbf{Doğru Cevap: D} \\
    \textbf{Açıklama:} Sayfa 136 Şiddet Tablosunda: Ölüm veya sürekli iş göremezlik durumları Çok Ciddi (5) derecesindedir. Ciddi (4) derecesi ise uzuv kaybı gibi sürekli kısmi iş göremezliğe denk gelir.
\end{testcard}

\begin{testcard}{Soru 3}
    Sınavda kağıt üzerinde verilecek olan Sayfa 135 Olasılık ve Sayfa 136 Şiddet tabloları kullanılarak yapılan bir risk analizinde; olasılık değeri 'Orta (3)' ve şiddet değeri 'Ciddi (4)' olarak belirlenmiştir. L-Tipi matrise göre risk puanı kaçtır?
    \begin{itemize}
        \item[A)] 7
        \item[B)] 9
        \item[C)] 12
        \item[D)] 15
        \item[E)] Hiçbiri
    \end{itemize}
    \tcblower
    \textbf{Doğru Cevap: C} \\
    \textbf{Açıklama:} L-Tipi matriste Risk = Olasılık x Şiddet formülü kullanılır. Puan: 3 (Orta) x 4 (Ciddi) = 12 olarak hesaplanır.
\end{testcard}

\begin{testcard}{Soru 4}
    Olasılık değeri 'Çok Yüksek (5)' ve şiddet değeri 'Çok Ciddi (5)' olan bir faaliyetin L-Tipi matris risk skoru kaçtır ve aksiyonu ne olmalıdır?
    \begin{itemize}
        \item[A)] 10 / Dikkate değer risk, önlemler alınmalıdır.
        \item[B)] 25 / Kabul edilemez risk, iş derhal durdurulmalıdır.
        \item[C)] 15 / Yüksek risk, 1 ay içinde önlem alınmalıdır.
        \item[D)] 20 / Kabul edilemez risk, rutin kontroller yeterlidir.
        \item[E)] Hiçbiri
    \end{itemize}
    \tcblower
    \textbf{Doğru Cevap: B} \\
    \textbf{Açıklama:} Risk skoru: 5 x 5 = 25'tir. Bu puan Sayfa 139 tablosuna göre 'Kabul Edilemez Risk'tir ve iş derhal durdurulmalıdır.
\end{testcard}

\begin{testcard}{Soru 5}
    Sınav kağıdındaki Sayfa 136 Şiddet tablosuna göre: 'İlkyardım gerektiren hafif yaralanmalar' hangi şiddet derecesine girmektedir?
    \begin{itemize}
        \item[A)] Çok Hafif (1)
        \item[B)] Hafif (2)
        \item[C)] Orta (3)
        \item[D)] Ciddi (4)
        \item[E)] Hiçbiri
    \end{itemize}
    \tcblower
    \textbf{Doğru Cevap: A} \\
    \textbf{Açıklama:} İlkyardım gerektiren iş günü kaybı olmayan yaralanmalar Çok Hafif (1) derecesindedir.
\end{testcard}

\begin{testcard}{Soru 6}
    Sınavda kağıtta verilmeyecek olan ve ezberlenmesi gereken Sayfa 139 aksiyon tablosuna göre, aşağıdaki puanlardan hangisi 'Dikkate Değer Risk' sınıfındadır?
    \begin{itemize}
        \item[A)] 5 Puan
        \item[B)] 10 Puan (veya 8, 9, 12)
        \item[C)] 16 Puan
        \item[D)] 25 Puan
        \item[E)] Hiçbiri
    \end{itemize}
    \tcblower
    \textbf{Doğru Cevap: B} \\
    \textbf{Açıklama:} Sayfa 139 tablosuna göre 8, 9, 10, 12 puanları Dikkate Değer Risk sınıfındadır.
\end{testcard}

\begin{testcard}{Soru 7}
    Ezberlenmesi gereken Sayfa 139 aksiyon tablosuna göre, hangi puan aralığındaki riskler 'Kabul Edilebilir Risk' olarak sınıflandırılmıştır?
    \begin{itemize}
        \item[A)] 1 ila 6 puan arası
        \item[B)] 8 ila 12 puan arası
        \item[C)] 15 ila 25 puan arası
        \item[D)] Sadece 25 puan
        \item[E)] Hiçbiri
    \end{itemize}
    \tcblower
    \textbf{Doğru Cevap: A} \\
    \textbf{Açıklama:} 1-6 puan aralığı Kabul Edilebilir Risktir. Acil tedbir gerektirmez, rutin kontroller yeterlidir.
\end{testcard}

\begin{testcard}{Soru 8}
    Sayfa 139 aksiyon tablosuna göre, 'Kabul Edilemez Risk' sınıfına giren puanlar hangileridir?
    \begin{itemize}
        \item[A)] 1, 2, 3
        \item[B)] 8, 9, 10
        \item[C)] 15, 16, 20, 25
        \item[D)] Sadece 12 puan
        \item[E)] Hiçbiri
    \end{itemize}
    \tcblower
    \textbf{Doğru Cevap: C} \\
    \textbf{Açıklama:} 15, 16, 20 ve 25 puanları Kabul Edilemez Risk sınıfında yer alır ve işin derhal durdurulmasını gerektirir.
\end{testcard}

\begin{testcard}{Soru 9}
    Sayfa 139 tablosuna göre 'Kabul Edilemez Risk' ile karşılaşıldığında yasal hüküm nedir?
    \begin{itemize}
        \item[A)] İş durdurulur, risk kabul edilebilir seviyeye çekilene kadar çalışmaya başlanmaz
        \item[B)] Rutin kontrollerle çalışmaya devam edilir
        \item[C)] Çalışanlara baret dağıtılıp çalışmaya devam edilir
        \item[D)] 3 ay içinde önlem alınır
        \item[E)] Hiçbiri
    \end{itemize}
    \tcblower
    \textbf{Doğru Cevap: A} \\
    \textbf{Açıklama:} Kabul edilemez risk en tehlikeli sınırdır. İş durdurulur ve risk düşürülene kadar çalışmaya izin verilmez.
\end{testcard}

\begin{testcard}{Soru 10}
    Sayfa 139 tablosuna göre 'Dikkate Değer Risk' tespit edildiğinde yapılması gereken eylem hangisidir?
    \begin{itemize}
        \item[A)] Mümkün olduğunca çabuk müdahale edilerek risk azaltılmalı ve ek denetimler konulmalıdır
        \item[B)] İşyeri tamamen kapatılmalıdır
        \item[C)] Herhangi bir önlem alınmasına gerek yoktur
        \item[D)] Sadece uzmana haber verilip rapor beklenir
        \item[E)] Hiçbiri
    \end{itemize}
    \tcblower
    \textbf{Doğru Cevap: A} \\
    \textbf{Açıklama:} Dikkate değer risk durumunda acil iş durdurma gerekmez ancak en kısa sürede iyileştirme yapılmalıdır.
\end{testcard}

\begin{testcard}{Soru 11}
    Sayfa 139 aksiyon tablosuna göre 'Kabul Edilebilir Risk' için hangisi doğrudur?
    \begin{itemize}
        \item[A)] Acil tedbir gerekmeyebilir ancak kontroller sürdürülmeli ve iyileştirmelere devam edilmelidir
        \item[B)] Risk tamamen sıfırlanmıştır, hiçbir denetime gerek yoktur
        \item[C)] Çalışanlar işi bırakmalıdır
        \item[D)] İşveren para cezası öder
        \item[E)] Hiçbiri
    \end{itemize}
    \tcblower
    \textbf{Doğru Cevap: A} \\
    \textbf{Açıklama:} Kabul edilebilir risk seviyesindeki tehlikeler sürekli gözetim altında tutularak durumun kötüleşmesi engellenir.
\end{testcard}

\begin{testcard}{Soru 12}
    9B ders notundaki sayfa 74, 114 ve 115'teki risk tablolarıyla ilgili hocanın uyarısı nedir?
    \begin{itemize}
        \item[A)] 9B'deki değerler hatalıdır, sınavda kesinlikle 9A sayfa 137 ve 139'daki doğru tablolar esas alınmalıdır
        \item[B)] 9B'deki tablolar daha doğrudur
        \item[C)] İki tablodan da sorulmayacaktır
        \item[D)] Tabloların hepsi aynıdır
        \item[E)] Hiçbiri
    \end{itemize}
    \tcblower
    \textbf{Doğru Cevap: A} \\
    \textbf{Açıklama:} Hoca 2.txt duyurusunda 9B'deki bazı değerlerin hatalı olduğunu, çalışırken 9A.pdf sayfa 137 ve 139'a bakılması gerektiğini belirtmiştir.
\end{testcard}

\begin{testcard}{Soru 13}
    L-Tipi risk değerlendirme matrisinde olasılık derecesi tablosuna (Sayfa 135) göre, bir tehlikenin haftada bir kez gerçekleşmesi hangi dereceye denk gelir? (Örnek No: 1)
    \begin{itemize}
        \item[A)] Orta (3) / Haftada bir veya daha az
        \item[B)] Yüksek (4) / Her gün
        \item[C)] Küçük (2) / Ayda bir
        \item[D)] Çok Küçük (1) / Yılda bir
        \item[E)] Hiçbiri
    \end{itemize}
    \tcblower
    \textbf{Doğru Cevap: A} \\
    \textbf{Açıklama:} Sayfa 135 Olasılık tablosunda 'Haftada bir veya daha az' olasılığı Orta (3) derecesindedir.
\end{testcard}

\begin{testcard}{Soru 14}
    L-Tipi risk değerlendirme matrisinde olasılık derecesi tablosuna (Sayfa 135) göre, bir tehlikenin haftada bir kez gerçekleşmesi hangi dereceye denk gelir? (Örnek No: 2)
    \begin{itemize}
        \item[A)] Orta (3) / Haftada bir veya daha az
        \item[B)] Yüksek (4) / Her gün
        \item[C)] Küçük (2) / Ayda bir
        \item[D)] Çok Küçük (1) / Yılda bir
        \item[E)] Hiçbiri
    \end{itemize}
    \tcblower
    \textbf{Doğru Cevap: A} \\
    \textbf{Açıklama:} Sayfa 135 Olasılık tablosunda 'Haftada bir veya daha az' olasılığı Orta (3) derecesindedir.
\end{testcard}

\begin{testcard}{Soru 15}
    L-Tipi risk değerlendirme matrisinde olasılık derecesi tablosuna (Sayfa 135) göre, bir tehlikenin haftada bir kez gerçekleşmesi hangi dereceye denk gelir? (Örnek No: 3)
    \begin{itemize}
        \item[A)] Orta (3) / Haftada bir veya daha az
        \item[B)] Yüksek (4) / Her gün
        \item[C)] Küçük (2) / Ayda bir
        \item[D)] Çok Küçük (1) / Yılda bir
        \item[E)] Hiçbiri
    \end{itemize}
    \tcblower
    \textbf{Doğru Cevap: A} \\
    \textbf{Açıklama:} Sayfa 135 Olasılık tablosunda 'Haftada bir veya daha az' olasılığı Orta (3) derecesindedir.
\end{testcard}

\begin{testcard}{Soru 16}
    L-Tipi risk değerlendirme matrisinde olasılık derecesi tablosuna (Sayfa 135) göre, bir tehlikenin haftada bir kez gerçekleşmesi hangi dereceye denk gelir? (Örnek No: 4)
    \begin{itemize}
        \item[A)] Orta (3) / Haftada bir veya daha az
        \item[B)] Yüksek (4) / Her gün
        \item[C)] Küçük (2) / Ayda bir
        \item[D)] Çok Küçük (1) / Yılda bir
        \item[E)] Hiçbiri
    \end{itemize}
    \tcblower
    \textbf{Doğru Cevap: A} \\
    \textbf{Açıklama:} Sayfa 135 Olasılık tablosunda 'Haftada bir veya daha az' olasılığı Orta (3) derecesindedir.
\end{testcard}

\begin{testcard}{Soru 17}
    L-Tipi risk değerlendirme matrisinde olasılık derecesi tablosuna (Sayfa 135) göre, bir tehlikenin haftada bir kez gerçekleşmesi hangi dereceye denk gelir? (Örnek No: 5)
    \begin{itemize}
        \item[A)] Orta (3) / Haftada bir veya daha az
        \item[B)] Yüksek (4) / Her gün
        \item[C)] Küçük (2) / Ayda bir
        \item[D)] Çok Küçük (1) / Yılda bir
        \item[E)] Hiçbiri
    \end{itemize}
    \tcblower
    \textbf{Doğru Cevap: A} \\
    \textbf{Açıklama:} Sayfa 135 Olasılık tablosunda 'Haftada bir veya daha az' olasılığı Orta (3) derecesindedir.
\end{testcard}

\begin{testcard}{Soru 18}
    L-Tipi risk değerlendirme matrisinde olasılık derecesi tablosuna (Sayfa 135) göre, bir tehlikenin haftada bir kez gerçekleşmesi hangi dereceye denk gelir? (Örnek No: 6)
    \begin{itemize}
        \item[A)] Orta (3) / Haftada bir veya daha az
        \item[B)] Yüksek (4) / Her gün
        \item[C)] Küçük (2) / Ayda bir
        \item[D)] Çok Küçük (1) / Yılda bir
        \item[E)] Hiçbiri
    \end{itemize}
    \tcblower
    \textbf{Doğru Cevap: A} \\
    \textbf{Açıklama:} Sayfa 135 Olasılık tablosunda 'Haftada bir veya daha az' olasılığı Orta (3) derecesindedir.
\end{testcard}

\begin{testcard}{Soru 19}
    L-Tipi risk değerlendirme matrisinde olasılık derecesi tablosuna (Sayfa 135) göre, bir tehlikenin haftada bir kez gerçekleşmesi hangi dereceye denk gelir? (Örnek No: 7)
    \begin{itemize}
        \item[A)] Orta (3) / Haftada bir veya daha az
        \item[B)] Yüksek (4) / Her gün
        \item[C)] Küçük (2) / Ayda bir
        \item[D)] Çok Küçük (1) / Yılda bir
        \item[E)] Hiçbiri
    \end{itemize}
    \tcblower
    \textbf{Doğru Cevap: A} \\
    \textbf{Açıklama:} Sayfa 135 Olasılık tablosunda 'Haftada bir veya daha az' olasılığı Orta (3) derecesindedir.
\end{testcard}

\begin{testcard}{Soru 20}
    L-Tipi risk değerlendirme matrisinde olasılık derecesi tablosuna (Sayfa 135) göre, bir tehlikenin haftada bir kez gerçekleşmesi hangi dereceye denk gelir? (Örnek No: 8)
    \begin{itemize}
        \item[A)] Orta (3) / Haftada bir veya daha az
        \item[B)] Yüksek (4) / Her gün
        \item[C)] Küçük (2) / Ayda bir
        \item[D)] Çok Küçük (1) / Yılda bir
        \item[E)] Hiçbiri
    \end{itemize}
    \tcblower
    \textbf{Doğru Cevap: A} \\
    \textbf{Açıklama:} Sayfa 135 Olasılık tablosunda 'Haftada bir veya daha az' olasılığı Orta (3) derecesindedir.
\end{testcard}

\begin{testcard}{Soru 21}
    L-Tipi risk değerlendirme matrisinde olasılık derecesi tablosuna (Sayfa 135) göre, bir tehlikenin haftada bir kez gerçekleşmesi hangi dereceye denk gelir? (Örnek No: 9)
    \begin{itemize}
        \item[A)] Orta (3) / Haftada bir veya daha az
        \item[B)] Yüksek (4) / Her gün
        \item[C)] Küçük (2) / Ayda bir
        \item[D)] Çok Küçük (1) / Yılda bir
        \item[E)] Hiçbiri
    \end{itemize}
    \tcblower
    \textbf{Doğru Cevap: A} \\
    \textbf{Açıklama:} Sayfa 135 Olasılık tablosunda 'Haftada bir veya daha az' olasılığı Orta (3) derecesindedir.
\end{testcard}

\begin{testcard}{Soru 22}
    L-Tipi risk değerlendirme matrisinde olasılık derecesi tablosuna (Sayfa 135) göre, bir tehlikenin haftada bir kez gerçekleşmesi hangi dereceye denk gelir? (Örnek No: 10)
    \begin{itemize}
        \item[A)] Orta (3) / Haftada bir veya daha az
        \item[B)] Yüksek (4) / Her gün
        \item[C)] Küçük (2) / Ayda bir
        \item[D)] Çok Küçük (1) / Yılda bir
        \item[E)] Hiçbiri
    \end{itemize}
    \tcblower
    \textbf{Doğru Cevap: A} \\
    \textbf{Açıklama:} Sayfa 135 Olasılık tablosunda 'Haftada bir veya daha az' olasılığı Orta (3) derecesindedir.
\end{testcard}

\begin{testcard}{Soru 23}
    L-Tipi risk değerlendirme matrisinde olasılık derecesi tablosuna (Sayfa 135) göre, bir tehlikenin haftada bir kez gerçekleşmesi hangi dereceye denk gelir? (Örnek No: 11)
    \begin{itemize}
        \item[A)] Orta (3) / Haftada bir veya daha az
        \item[B)] Yüksek (4) / Her gün
        \item[C)] Küçük (2) / Ayda bir
        \item[D)] Çok Küçük (1) / Yılda bir
        \item[E)] Hiçbiri
    \end{itemize}
    \tcblower
    \textbf{Doğru Cevap: A} \\
    \textbf{Açıklama:} Sayfa 135 Olasılık tablosunda 'Haftada bir veya daha az' olasılığı Orta (3) derecesindedir.
\end{testcard}

\begin{testcard}{Soru 24}
    L-Tipi risk değerlendirme matrisinde olasılık derecesi tablosuna (Sayfa 135) göre, bir tehlikenin haftada bir kez gerçekleşmesi hangi dereceye denk gelir? (Örnek No: 12)
    \begin{itemize}
        \item[A)] Orta (3) / Haftada bir veya daha az
        \item[B)] Yüksek (4) / Her gün
        \item[C)] Küçük (2) / Ayda bir
        \item[D)] Çok Küçük (1) / Yılda bir
        \item[E)] Hiçbiri
    \end{itemize}
    \tcblower
    \textbf{Doğru Cevap: A} \\
    \textbf{Açıklama:} Sayfa 135 Olasılık tablosunda 'Haftada bir veya daha az' olasılığı Orta (3) derecesindedir.
\end{testcard}

\begin{testcard}{Soru 25}
    L-Tipi risk değerlendirme matrisinde olasılık derecesi tablosuna (Sayfa 135) göre, bir tehlikenin haftada bir kez gerçekleşmesi hangi dereceye denk gelir? (Örnek No: 13)
    \begin{itemize}
        \item[A)] Orta (3) / Haftada bir veya daha az
        \item[B)] Yüksek (4) / Her gün
        \item[C)] Küçük (2) / Ayda bir
        \item[D)] Çok Küçük (1) / Yılda bir
        \item[E)] Hiçbiri
    \end{itemize}
    \tcblower
    \textbf{Doğru Cevap: A} \\
    \textbf{Açıklama:} Sayfa 135 Olasılık tablosunda 'Haftada bir veya daha az' olasılığı Orta (3) derecesindedir.
\end{testcard}

\begin{testcard}{Soru 26}
    L-Tipi risk değerlendirme matrisinde olasılık derecesi tablosuna (Sayfa 135) göre, bir tehlikenin haftada bir kez gerçekleşmesi hangi dereceye denk gelir? (Örnek No: 14)
    \begin{itemize}
        \item[A)] Orta (3) / Haftada bir veya daha az
        \item[B)] Yüksek (4) / Her gün
        \item[C)] Küçük (2) / Ayda bir
        \item[D)] Çok Küçük (1) / Yılda bir
        \item[E)] Hiçbiri
    \end{itemize}
    \tcblower
    \textbf{Doğru Cevap: A} \\
    \textbf{Açıklama:} Sayfa 135 Olasılık tablosunda 'Haftada bir veya daha az' olasılığı Orta (3) derecesindedir.
\end{testcard}

\begin{testcard}{Soru 27}
    L-Tipi risk değerlendirme matrisinde olasılık derecesi tablosuna (Sayfa 135) göre, bir tehlikenin haftada bir kez gerçekleşmesi hangi dereceye denk gelir? (Örnek No: 15)
    \begin{itemize}
        \item[A)] Orta (3) / Haftada bir veya daha az
        \item[B)] Yüksek (4) / Her gün
        \item[C)] Küçük (2) / Ayda bir
        \item[D)] Çok Küçük (1) / Yılda bir
        \item[E)] Hiçbiri
    \end{itemize}
    \tcblower
    \textbf{Doğru Cevap: A} \\
    \textbf{Açıklama:} Sayfa 135 Olasılık tablosunda 'Haftada bir veya daha az' olasılığı Orta (3) derecesindedir.
\end{testcard}

\begin{testcard}{Soru 28}
    L-Tipi risk değerlendirme matrisinde olasılık derecesi tablosuna (Sayfa 135) göre, bir tehlikenin haftada bir kez gerçekleşmesi hangi dereceye denk gelir? (Örnek No: 16)
    \begin{itemize}
        \item[A)] Orta (3) / Haftada bir veya daha az
        \item[B)] Yüksek (4) / Her gün
        \item[C)] Küçük (2) / Ayda bir
        \item[D)] Çok Küçük (1) / Yılda bir
        \item[E)] Hiçbiri
    \end{itemize}
    \tcblower
    \textbf{Doğru Cevap: A} \\
    \textbf{Açıklama:} Sayfa 135 Olasılık tablosunda 'Haftada bir veya daha az' olasılığı Orta (3) derecesindedir.
\end{testcard}

\begin{testcard}{Soru 29}
    L-Tipi risk değerlendirme matrisinde olasılık derecesi tablosuna (Sayfa 135) göre, bir tehlikenin haftada bir kez gerçekleşmesi hangi dereceye denk gelir? (Örnek No: 17)
    \begin{itemize}
        \item[A)] Orta (3) / Haftada bir veya daha az
        \item[B)] Yüksek (4) / Her gün
        \item[C)] Küçük (2) / Ayda bir
        \item[D)] Çok Küçük (1) / Yılda bir
        \item[E)] Hiçbiri
    \end{itemize}
    \tcblower
    \textbf{Doğru Cevap: A} \\
    \textbf{Açıklama:} Sayfa 135 Olasılık tablosunda 'Haftada bir veya daha az' olasılığı Orta (3) derecesindedir.
\end{testcard}

\begin{testcard}{Soru 30}
    L-Tipi risk değerlendirme matrisinde olasılık derecesi tablosuna (Sayfa 135) göre, bir tehlikenin haftada bir kez gerçekleşmesi hangi dereceye denk gelir? (Örnek No: 18)
    \begin{itemize}
        \item[A)] Orta (3) / Haftada bir veya daha az
        \item[B)] Yüksek (4) / Her gün
        \item[C)] Küçük (2) / Ayda bir
        \item[D)] Çok Küçük (1) / Yılda bir
        \item[E)] Hiçbiri
    \end{itemize}
    \tcblower
    \textbf{Doğru Cevap: A} \\
    \textbf{Açıklama:} Sayfa 135 Olasılık tablosunda 'Haftada bir veya daha az' olasılığı Orta (3) derecesindedir.
\end{testcard}

\begin{testcard}{Soru 31}
    L-Tipi risk değerlendirme matrisinde olasılık derecesi tablosuna (Sayfa 135) göre, bir tehlikenin haftada bir kez gerçekleşmesi hangi dereceye denk gelir? (Örnek No: 19)
    \begin{itemize}
        \item[A)] Orta (3) / Haftada bir veya daha az
        \item[B)] Yüksek (4) / Her gün
        \item[C)] Küçük (2) / Ayda bir
        \item[D)] Çok Küçük (1) / Yılda bir
        \item[E)] Hiçbiri
    \end{itemize}
    \tcblower
    \textbf{Doğru Cevap: A} \\
    \textbf{Açıklama:} Sayfa 135 Olasılık tablosunda 'Haftada bir veya daha az' olasılığı Orta (3) derecesindedir.
\end{testcard}

\begin{testcard}{Soru 32}
    L-Tipi risk değerlendirme matrisinde olasılık derecesi tablosuna (Sayfa 135) göre, bir tehlikenin haftada bir kez gerçekleşmesi hangi dereceye denk gelir? (Örnek No: 20)
    \begin{itemize}
        \item[A)] Orta (3) / Haftada bir veya daha az
        \item[B)] Yüksek (4) / Her gün
        \item[C)] Küçük (2) / Ayda bir
        \item[D)] Çok Küçük (1) / Yılda bir
        \item[E)] Hiçbiri
    \end{itemize}
    \tcblower
    \textbf{Doğru Cevap: A} \\
    \textbf{Açıklama:} Sayfa 135 Olasılık tablosunda 'Haftada bir veya daha az' olasılığı Orta (3) derecesindedir.
\end{testcard}

\begin{testcard}{Soru 33}
    L-Tipi risk değerlendirme matrisinde olasılık derecesi tablosuna (Sayfa 135) göre, bir tehlikenin haftada bir kez gerçekleşmesi hangi dereceye denk gelir? (Örnek No: 21)
    \begin{itemize}
        \item[A)] Orta (3) / Haftada bir veya daha az
        \item[B)] Yüksek (4) / Her gün
        \item[C)] Küçük (2) / Ayda bir
        \item[D)] Çok Küçük (1) / Yılda bir
        \item[E)] Hiçbiri
    \end{itemize}
    \tcblower
    \textbf{Doğru Cevap: A} \\
    \textbf{Açıklama:} Sayfa 135 Olasılık tablosunda 'Haftada bir veya daha az' olasılığı Orta (3) derecesindedir.
\end{testcard}

\begin{testcard}{Soru 34}
    L-Tipi risk değerlendirme matrisinde olasılık derecesi tablosuna (Sayfa 135) göre, bir tehlikenin haftada bir kez gerçekleşmesi hangi dereceye denk gelir? (Örnek No: 22)
    \begin{itemize}
        \item[A)] Orta (3) / Haftada bir veya daha az
        \item[B)] Yüksek (4) / Her gün
        \item[C)] Küçük (2) / Ayda bir
        \item[D)] Çok Küçük (1) / Yılda bir
        \item[E)] Hiçbiri
    \end{itemize}
    \tcblower
    \textbf{Doğru Cevap: A} \\
    \textbf{Açıklama:} Sayfa 135 Olasılık tablosunda 'Haftada bir veya daha az' olasılığı Orta (3) derecesindedir.
\end{testcard}

\begin{testcard}{Soru 35}
    L-Tipi risk değerlendirme matrisinde olasılık derecesi tablosuna (Sayfa 135) göre, bir tehlikenin haftada bir kez gerçekleşmesi hangi dereceye denk gelir? (Örnek No: 23)
    \begin{itemize}
        \item[A)] Orta (3) / Haftada bir veya daha az
        \item[B)] Yüksek (4) / Her gün
        \item[C)] Küçük (2) / Ayda bir
        \item[D)] Çok Küçük (1) / Yılda bir
        \item[E)] Hiçbiri
    \end{itemize}
    \tcblower
    \textbf{Doğru Cevap: A} \\
    \textbf{Açıklama:} Sayfa 135 Olasılık tablosunda 'Haftada bir veya daha az' olasılığı Orta (3) derecesindedir.
\end{testcard}

\begin{testcard}{Soru 36}
    L-Tipi risk değerlendirme matrisinde olasılık derecesi tablosuna (Sayfa 135) göre, bir tehlikenin haftada bir kez gerçekleşmesi hangi dereceye denk gelir? (Örnek No: 24)
    \begin{itemize}
        \item[A)] Orta (3) / Haftada bir veya daha az
        \item[B)] Yüksek (4) / Her gün
        \item[C)] Küçük (2) / Ayda bir
        \item[D)] Çok Küçük (1) / Yılda bir
        \item[E)] Hiçbiri
    \end{itemize}
    \tcblower
    \textbf{Doğru Cevap: A} \\
    \textbf{Açıklama:} Sayfa 135 Olasılık tablosunda 'Haftada bir veya daha az' olasılığı Orta (3) derecesindedir.
\end{testcard}

\begin{testcard}{Soru 37}
    L-Tipi risk değerlendirme matrisinde olasılık derecesi tablosuna (Sayfa 135) göre, bir tehlikenin haftada bir kez gerçekleşmesi hangi dereceye denk gelir? (Örnek No: 25)
    \begin{itemize}
        \item[A)] Orta (3) / Haftada bir veya daha az
        \item[B)] Yüksek (4) / Her gün
        \item[C)] Küçük (2) / Ayda bir
        \item[D)] Çok Küçük (1) / Yılda bir
        \item[E)] Hiçbiri
    \end{itemize}
    \tcblower
    \textbf{Doğru Cevap: A} \\
    \textbf{Açıklama:} Sayfa 135 Olasılık tablosunda 'Haftada bir veya daha az' olasılığı Orta (3) derecesindedir.
\end{testcard}

\begin{testcard}{Soru 38}
    L-Tipi risk değerlendirme matrisinde olasılık derecesi tablosuna (Sayfa 135) göre, bir tehlikenin haftada bir kez gerçekleşmesi hangi dereceye denk gelir? (Örnek No: 26)
    \begin{itemize}
        \item[A)] Orta (3) / Haftada bir veya daha az
        \item[B)] Yüksek (4) / Her gün
        \item[C)] Küçük (2) / Ayda bir
        \item[D)] Çok Küçük (1) / Yılda bir
        \item[E)] Hiçbiri
    \end{itemize}
    \tcblower
    \textbf{Doğru Cevap: A} \\
    \textbf{Açıklama:} Sayfa 135 Olasılık tablosunda 'Haftada bir veya daha az' olasılığı Orta (3) derecesindedir.
\end{testcard}

\begin{testcard}{Soru 39}
    L-Tipi risk değerlendirme matrisinde olasılık derecesi tablosuna (Sayfa 135) göre, bir tehlikenin haftada bir kez gerçekleşmesi hangi dereceye denk gelir? (Örnek No: 27)
    \begin{itemize}
        \item[A)] Orta (3) / Haftada bir veya daha az
        \item[B)] Yüksek (4) / Her gün
        \item[C)] Küçük (2) / Ayda bir
        \item[D)] Çok Küçük (1) / Yılda bir
        \item[E)] Hiçbiri
    \end{itemize}
    \tcblower
    \textbf{Doğru Cevap: A} \\
    \textbf{Açıklama:} Sayfa 135 Olasılık tablosunda 'Haftada bir veya daha az' olasılığı Orta (3) derecesindedir.
\end{testcard}

\begin{testcard}{Soru 40}
    L-Tipi risk değerlendirme matrisinde olasılık derecesi tablosuna (Sayfa 135) göre, bir tehlikenin haftada bir kez gerçekleşmesi hangi dereceye denk gelir? (Örnek No: 28)
    \begin{itemize}
        \item[A)] Orta (3) / Haftada bir veya daha az
        \item[B)] Yüksek (4) / Her gün
        \item[C)] Küçük (2) / Ayda bir
        \item[D)] Çok Küçük (1) / Yılda bir
        \item[E)] Hiçbiri
    \end{itemize}
    \tcblower
    \textbf{Doğru Cevap: A} \\
    \textbf{Açıklama:} Sayfa 135 Olasılık tablosunda 'Haftada bir veya daha az' olasılığı Orta (3) derecesindedir.
\end{testcard}

\begin{testcard}{Soru 41}
    L-Tipi risk değerlendirme matrisinde olasılık derecesi tablosuna (Sayfa 135) göre, bir tehlikenin haftada bir kez gerçekleşmesi hangi dereceye denk gelir? (Örnek No: 29)
    \begin{itemize}
        \item[A)] Orta (3) / Haftada bir veya daha az
        \item[B)] Yüksek (4) / Her gün
        \item[C)] Küçük (2) / Ayda bir
        \item[D)] Çok Küçük (1) / Yılda bir
        \item[E)] Hiçbiri
    \end{itemize}
    \tcblower
    \textbf{Doğru Cevap: A} \\
    \textbf{Açıklama:} Sayfa 135 Olasılık tablosunda 'Haftada bir veya daha az' olasılığı Orta (3) derecesindedir.
\end{testcard}

\begin{testcard}{Soru 42}
    L-Tipi risk değerlendirme matrisinde olasılık derecesi tablosuna (Sayfa 135) göre, bir tehlikenin haftada bir kez gerçekleşmesi hangi dereceye denk gelir? (Örnek No: 30)
    \begin{itemize}
        \item[A)] Orta (3) / Haftada bir veya daha az
        \item[B)] Yüksek (4) / Her gün
        \item[C)] Küçük (2) / Ayda bir
        \item[D)] Çok Küçük (1) / Yılda bir
        \item[E)] Hiçbiri
    \end{itemize}
    \tcblower
    \textbf{Doğru Cevap: A} \\
    \textbf{Açıklama:} Sayfa 135 Olasılık tablosunda 'Haftada bir veya daha az' olasılığı Orta (3) derecesindedir.
\end{testcard}

\begin{testcard}{Soru 43}
    L-Tipi risk değerlendirme matrisinde olasılık derecesi tablosuna (Sayfa 135) göre, bir tehlikenin haftada bir kez gerçekleşmesi hangi dereceye denk gelir? (Örnek No: 31)
    \begin{itemize}
        \item[A)] Orta (3) / Haftada bir veya daha az
        \item[B)] Yüksek (4) / Her gün
        \item[C)] Küçük (2) / Ayda bir
        \item[D)] Çok Küçük (1) / Yılda bir
        \item[E)] Hiçbiri
    \end{itemize}
    \tcblower
    \textbf{Doğru Cevap: A} \\
    \textbf{Açıklama:} Sayfa 135 Olasılık tablosunda 'Haftada bir veya daha az' olasılığı Orta (3) derecesindedir.
\end{testcard}

\begin{testcard}{Soru 44}
    L-Tipi risk değerlendirme matrisinde olasılık derecesi tablosuna (Sayfa 135) göre, bir tehlikenin haftada bir kez gerçekleşmesi hangi dereceye denk gelir? (Örnek No: 32)
    \begin{itemize}
        \item[A)] Orta (3) / Haftada bir veya daha az
        \item[B)] Yüksek (4) / Her gün
        \item[C)] Küçük (2) / Ayda bir
        \item[D)] Çok Küçük (1) / Yılda bir
        \item[E)] Hiçbiri
    \end{itemize}
    \tcblower
    \textbf{Doğru Cevap: A} \\
    \textbf{Açıklama:} Sayfa 135 Olasılık tablosunda 'Haftada bir veya daha az' olasılığı Orta (3) derecesindedir.
\end{testcard}

\begin{testcard}{Soru 45}
    L-Tipi risk değerlendirme matrisinde olasılık derecesi tablosuna (Sayfa 135) göre, bir tehlikenin haftada bir kez gerçekleşmesi hangi dereceye denk gelir? (Örnek No: 33)
    \begin{itemize}
        \item[A)] Orta (3) / Haftada bir veya daha az
        \item[B)] Yüksek (4) / Her gün
        \item[C)] Küçük (2) / Ayda bir
        \item[D)] Çok Küçük (1) / Yılda bir
        \item[E)] Hiçbiri
    \end{itemize}
    \tcblower
    \textbf{Doğru Cevap: A} \\
    \textbf{Açıklama:} Sayfa 135 Olasılık tablosunda 'Haftada bir veya daha az' olasılığı Orta (3) derecesindedir.
\end{testcard}

\begin{testcard}{Soru 46}
    L-Tipi risk değerlendirme matrisinde olasılık derecesi tablosuna (Sayfa 135) göre, bir tehlikenin haftada bir kez gerçekleşmesi hangi dereceye denk gelir? (Örnek No: 34)
    \begin{itemize}
        \item[A)] Orta (3) / Haftada bir veya daha az
        \item[B)] Yüksek (4) / Her gün
        \item[C)] Küçük (2) / Ayda bir
        \item[D)] Çok Küçük (1) / Yılda bir
        \item[E)] Hiçbiri
    \end{itemize}
    \tcblower
    \textbf{Doğru Cevap: A} \\
    \textbf{Açıklama:} Sayfa 135 Olasılık tablosunda 'Haftada bir veya daha az' olasılığı Orta (3) derecesindedir.
\end{testcard}

\begin{testcard}{Soru 47}
    L-Tipi risk değerlendirme matrisinde olasılık derecesi tablosuna (Sayfa 135) göre, bir tehlikenin haftada bir kez gerçekleşmesi hangi dereceye denk gelir? (Örnek No: 35)
    \begin{itemize}
        \item[A)] Orta (3) / Haftada bir veya daha az
        \item[B)] Yüksek (4) / Her gün
        \item[C)] Küçük (2) / Ayda bir
        \item[D)] Çok Küçük (1) / Yılda bir
        \item[E)] Hiçbiri
    \end{itemize}
    \tcblower
    \textbf{Doğru Cevap: A} \\
    \textbf{Açıklama:} Sayfa 135 Olasılık tablosunda 'Haftada bir veya daha az' olasılığı Orta (3) derecesindedir.
\end{testcard}

\begin{testcard}{Soru 48}
    L-Tipi risk değerlendirme matrisinde olasılık derecesi tablosuna (Sayfa 135) göre, bir tehlikenin haftada bir kez gerçekleşmesi hangi dereceye denk gelir? (Örnek No: 36)
    \begin{itemize}
        \item[A)] Orta (3) / Haftada bir veya daha az
        \item[B)] Yüksek (4) / Her gün
        \item[C)] Küçük (2) / Ayda bir
        \item[D)] Çok Küçük (1) / Yılda bir
        \item[E)] Hiçbiri
    \end{itemize}
    \tcblower
    \textbf{Doğru Cevap: A} \\
    \textbf{Açıklama:} Sayfa 135 Olasılık tablosunda 'Haftada bir veya daha az' olasılığı Orta (3) derecesindedir.
\end{testcard}

\begin{testcard}{Soru 49}
    L-Tipi risk değerlendirme matrisinde olasılık derecesi tablosuna (Sayfa 135) göre, bir tehlikenin haftada bir kez gerçekleşmesi hangi dereceye denk gelir? (Örnek No: 37)
    \begin{itemize}
        \item[A)] Orta (3) / Haftada bir veya daha az
        \item[B)] Yüksek (4) / Her gün
        \item[C)] Küçük (2) / Ayda bir
        \item[D)] Çok Küçük (1) / Yılda bir
        \item[E)] Hiçbiri
    \end{itemize}
    \tcblower
    \textbf{Doğru Cevap: A} \\
    \textbf{Açıklama:} Sayfa 135 Olasılık tablosunda 'Haftada bir veya daha az' olasılığı Orta (3) derecesindedir.
\end{testcard}

\begin{testcard}{Soru 50}
    L-Tipi risk değerlendirme matrisinde olasılık derecesi tablosuna (Sayfa 135) göre, bir tehlikenin haftada bir kez gerçekleşmesi hangi dereceye denk gelir? (Örnek No: 38)
    \begin{itemize}
        \item[A)] Orta (3) / Haftada bir veya daha az
        \item[B)] Yüksek (4) / Her gün
        \item[C)] Küçük (2) / Ayda bir
        \item[D)] Çok Küçük (1) / Yılda bir
        \item[E)] Hiçbiri
    \end{itemize}
    \tcblower
    \textbf{Doğru Cevap: A} \\
    \textbf{Açıklama:} Sayfa 135 Olasılık tablosunda 'Haftada bir veya daha az' olasılığı Orta (3) derecesindedir.
\end{testcard}

\begin{testcard}{Soru 51}
    L-Tipi risk değerlendirme matrisinde olasılık derecesi tablosuna (Sayfa 135) göre, bir tehlikenin haftada bir kez gerçekleşmesi hangi dereceye denk gelir? (Örnek No: 39)
    \begin{itemize}
        \item[A)] Orta (3) / Haftada bir veya daha az
        \item[B)] Yüksek (4) / Her gün
        \item[C)] Küçük (2) / Ayda bir
        \item[D)] Çok Küçük (1) / Yılda bir
        \item[E)] Hiçbiri
    \end{itemize}
    \tcblower
    \textbf{Doğru Cevap: A} \\
    \textbf{Açıklama:} Sayfa 135 Olasılık tablosunda 'Haftada bir veya daha az' olasılığı Orta (3) derecesindedir.
\end{testcard}

\begin{testcard}{Soru 52}
    L-Tipi risk değerlendirme matrisinde olasılık derecesi tablosuna (Sayfa 135) göre, bir tehlikenin haftada bir kez gerçekleşmesi hangi dereceye denk gelir? (Örnek No: 40)
    \begin{itemize}
        \item[A)] Orta (3) / Haftada bir veya daha az
        \item[B)] Yüksek (4) / Her gün
        \item[C)] Küçük (2) / Ayda bir
        \item[D)] Çok Küçük (1) / Yılda bir
        \item[E)] Hiçbiri
    \end{itemize}
    \tcblower
    \textbf{Doğru Cevap: A} \\
    \textbf{Açıklama:} Sayfa 135 Olasılık tablosunda 'Haftada bir veya daha az' olasılığı Orta (3) derecesindedir.
\end{testcard}

\begin{testcard}{Soru 53}
    L-Tipi risk değerlendirme matrisinde olasılık derecesi tablosuna (Sayfa 135) göre, bir tehlikenin haftada bir kez gerçekleşmesi hangi dereceye denk gelir? (Örnek No: 41)
    \begin{itemize}
        \item[A)] Orta (3) / Haftada bir veya daha az
        \item[B)] Yüksek (4) / Her gün
        \item[C)] Küçük (2) / Ayda bir
        \item[D)] Çok Küçük (1) / Yılda bir
        \item[E)] Hiçbiri
    \end{itemize}
    \tcblower
    \textbf{Doğru Cevap: A} \\
    \textbf{Açıklama:} Sayfa 135 Olasılık tablosunda 'Haftada bir veya daha az' olasılığı Orta (3) derecesindedir.
\end{testcard}

\begin{testcard}{Soru 54}
    L-Tipi risk değerlendirme matrisinde olasılık derecesi tablosuna (Sayfa 135) göre, bir tehlikenin haftada bir kez gerçekleşmesi hangi dereceye denk gelir? (Örnek No: 42)
    \begin{itemize}
        \item[A)] Orta (3) / Haftada bir veya daha az
        \item[B)] Yüksek (4) / Her gün
        \item[C)] Küçük (2) / Ayda bir
        \item[D)] Çok Küçük (1) / Yılda bir
        \item[E)] Hiçbiri
    \end{itemize}
    \tcblower
    \textbf{Doğru Cevap: A} \\
    \textbf{Açıklama:} Sayfa 135 Olasılık tablosunda 'Haftada bir veya daha az' olasılığı Orta (3) derecesindedir.
\end{testcard}

\begin{testcard}{Soru 55}
    L-Tipi risk değerlendirme matrisinde olasılık derecesi tablosuna (Sayfa 135) göre, bir tehlikenin haftada bir kez gerçekleşmesi hangi dereceye denk gelir? (Örnek No: 43)
    \begin{itemize}
        \item[A)] Orta (3) / Haftada bir veya daha az
        \item[B)] Yüksek (4) / Her gün
        \item[C)] Küçük (2) / Ayda bir
        \item[D)] Çok Küçük (1) / Yılda bir
        \item[E)] Hiçbiri
    \end{itemize}
    \tcblower
    \textbf{Doğru Cevap: A} \\
    \textbf{Açıklama:} Sayfa 135 Olasılık tablosunda 'Haftada bir veya daha az' olasılığı Orta (3) derecesindedir.
\end{testcard}

\begin{testcard}{Soru 56}
    L-Tipi risk değerlendirme matrisinde olasılık derecesi tablosuna (Sayfa 135) göre, bir tehlikenin haftada bir kez gerçekleşmesi hangi dereceye denk gelir? (Örnek No: 44)
    \begin{itemize}
        \item[A)] Orta (3) / Haftada bir veya daha az
        \item[B)] Yüksek (4) / Her gün
        \item[C)] Küçük (2) / Ayda bir
        \item[D)] Çok Küçük (1) / Yılda bir
        \item[E)] Hiçbiri
    \end{itemize}
    \tcblower
    \textbf{Doğru Cevap: A} \\
    \textbf{Açıklama:} Sayfa 135 Olasılık tablosunda 'Haftada bir veya daha az' olasılığı Orta (3) derecesindedir.
\end{testcard}

\begin{testcard}{Soru 57}
    L-Tipi risk değerlendirme matrisinde olasılık derecesi tablosuna (Sayfa 135) göre, bir tehlikenin haftada bir kez gerçekleşmesi hangi dereceye denk gelir? (Örnek No: 45)
    \begin{itemize}
        \item[A)] Orta (3) / Haftada bir veya daha az
        \item[B)] Yüksek (4) / Her gün
        \item[C)] Küçük (2) / Ayda bir
        \item[D)] Çok Küçük (1) / Yılda bir
        \item[E)] Hiçbiri
    \end{itemize}
    \tcblower
    \textbf{Doğru Cevap: A} \\
    \textbf{Açıklama:} Sayfa 135 Olasılık tablosunda 'Haftada bir veya daha az' olasılığı Orta (3) derecesindedir.
\end{testcard}

\begin{testcard}{Soru 58}
    L-Tipi risk değerlendirme matrisinde olasılık derecesi tablosuna (Sayfa 135) göre, bir tehlikenin haftada bir kez gerçekleşmesi hangi dereceye denk gelir? (Örnek No: 46)
    \begin{itemize}
        \item[A)] Orta (3) / Haftada bir veya daha az
        \item[B)] Yüksek (4) / Her gün
        \item[C)] Küçük (2) / Ayda bir
        \item[D)] Çok Küçük (1) / Yılda bir
        \item[E)] Hiçbiri
    \end{itemize}
    \tcblower
    \textbf{Doğru Cevap: A} \\
    \textbf{Açıklama:} Sayfa 135 Olasılık tablosunda 'Haftada bir veya daha az' olasılığı Orta (3) derecesindedir.
\end{testcard}

\begin{testcard}{Soru 59}
    L-Tipi risk değerlendirme matrisinde olasılık derecesi tablosuna (Sayfa 135) göre, bir tehlikenin haftada bir kez gerçekleşmesi hangi dereceye denk gelir? (Örnek No: 47)
    \begin{itemize}
        \item[A)] Orta (3) / Haftada bir veya daha az
        \item[B)] Yüksek (4) / Her gün
        \item[C)] Küçük (2) / Ayda bir
        \item[D)] Çok Küçük (1) / Yılda bir
        \item[E)] Hiçbiri
    \end{itemize}
    \tcblower
    \textbf{Doğru Cevap: A} \\
    \textbf{Açıklama:} Sayfa 135 Olasılık tablosunda 'Haftada bir veya daha az' olasılığı Orta (3) derecesindedir.
\end{testcard}

\begin{testcard}{Soru 60}
    L-Tipi risk değerlendirme matrisinde olasılık derecesi tablosuna (Sayfa 135) göre, bir tehlikenin haftada bir kez gerçekleşmesi hangi dereceye denk gelir? (Örnek No: 48)
    \begin{itemize}
        \item[A)] Orta (3) / Haftada bir veya daha az
        \item[B)] Yüksek (4) / Her gün
        \item[C)] Küçük (2) / Ayda bir
        \item[D)] Çok Küçük (1) / Yılda bir
        \item[E)] Hiçbiri
    \end{itemize}
    \tcblower
    \textbf{Doğru Cevap: A} \\
    \textbf{Açıklama:} Sayfa 135 Olasılık tablosunda 'Haftada bir veya daha az' olasılığı Orta (3) derecesindedir.
\end{testcard}

\subsection{Ünite 6: İş Hukuku ve Fesih Süreleri}
\begin{testcard}{Soru 1}
    4857 sayılı İş Kanunu'na göre, kıdemi fiilen 1 yıldır çalışan bir işçinin sözleşmesi feshedilecektir. Yasal ihbar süresi ne kadardır?
    \begin{itemize}
        \item[A)] 2 hafta
        \item[B)] 4 hafta
        \item[C)] 6 hafta
        \item[D)] 8 hafta
        \item[E)] Hiçbiri
    \end{itemize}
    \tcblower
    \textbf{Doğru Cevap: B} \\
    \textbf{Açıklama:} Kıdemi 6 ay ile 1.5 yıl arasında olan çalışanların ihbar (fesih bildirim) süresi 4 haftadır.
\end{testcard}

\begin{testcard}{Soru 2}
    4857 sayılı İş Kanunu'na göre, kıdemi 4 yıl olan bir işçinin uyması gereken yasal ihbar süresi ne kadardır?
    \begin{itemize}
        \item[A)] 2 hafta
        \item[B)] 4 hafta
        \item[C)] 6 hafta
        \item[D)] 8 hafta
        \item[E)] Hiçbiri
    \end{itemize}
    \tcblower
    \textbf{Doğru Cevap: D} \\
    \textbf{Açıklama:} Kıdemi 3 yıldan fazla sürmüş olan çalışanların ihbar süresi 8 haftadır.
\end{testcard}

\begin{testcard}{Soru 3}
    4857 sayılı İş Kanunu'na göre, kıdemi 2 yıl olan bir işçinin yasal ihbar süresi ne kadardır?
    \begin{itemize}
        \item[A)] 2 hafta
        \item[B)] 4 hafta
        \item[C)] 6 hafta
        \item[D)] 8 hafta
        \item[E)] Hiçbiri
    \end{itemize}
    \tcblower
    \textbf{Doğru Cevap: C} \\
    \textbf{Açıklama:} Kıdemi 1.5 yıl ile 3 yıl arasında olan çalışanlar için ihbar süresi 6 haftadır.
\end{testcard}

\begin{testcard}{Soru 4}
    4857 sayılı İş Kanunu'na göre, kıdemi 3 ay olan bir işçinin yasal ihbar süresi kaç haftadır?
    \begin{itemize}
        \item[A)] 2 hafta
        \item[B)] 4 hafta
        \item[C)] 6 hafta
        \item[D)] 8 hafta
        \item[E)] Hiçbiri
    \end{itemize}
    \tcblower
    \textbf{Doğru Cevap: A} \\
    \textbf{Açıklama:} Kıdemi 6 aydan az sürmüş olan çalışanların yasal ihbar süresi 2 haftadır.
\end{testcard}

\begin{testcard}{Soru 5}
    İşe iade davasını kazanan işçiyi işveren 1 ay içinde işe başlatmazsa ödemek zorunda kalacağı tazminat limiti nedir?
    \begin{itemize}
        \item[A)] En az 1, en çok 3 aylık ücreti tutarında
        \item[B)] En az 4, en çok 8 aylık ücreti tutarında
        \item[C)] En az 6, en çok 12 aylık ücreti tutarında
        \item[D)] Sabit 10 aylık ücret tutarı
        \item[E)] Hiçbiri
    \end{itemize}
    \tcblower
    \textbf{Doğru Cevap: B} \\
    \textbf{Açıklama:} Mahkeme kararı uyarınca işe başlatmama tazminatı en az 4 aylık, en çok 8 aylık ücret tutarında belirlenir.
\end{testcard}

\begin{testcard}{Soru 6}
    İş Kanunu'na göre, haftalık normal çalışma süresi en fazla kaç saattir?
    \begin{itemize}
        \item[A)] 40 saat
        \item[B)] 45 saat
        \item[C)] 48 saat
        \item[D)] 54 saat
        \item[E)] Hiçbiri
    \end{itemize}
    \tcblower
    \textbf{Doğru Cevap: B} \\
    \textbf{Açıklama:} 4857 sayılı İş Kanunu Madde 63 uyarınca haftalık çalışma süresi en çok 45 saattir.
\end{testcard}

\begin{testcard}{Soru 7}
    İş Kanunu'na göre, günlük çalışma süresi (fazla mesai dahil) hiçbir koşulda kaç saati aşamaz?
    \begin{itemize}
        \item[A)] 9 saat
        \item[B)] 11 saat
        \item[C)] 12 saat
        \item[D)] 15 saat
        \item[E)] Hiçbiri
    \end{itemize}
    \tcblower
    \textbf{Doğru Cevap: B} \\
    \textbf{Açıklama:} Yasa gereği günlük toplam çalışma süresi 11 saati aşamaz.
\end{testcard}

\begin{testcard}{Soru 8}
    İş Kanunu'na göre, bir işçiye yaptırılabilecek fazla çalışmanın (mesai) yıllık toplam sınırı kaç saattir?
    \begin{itemize}
        \item[A)] 180 saat
        \item[B)] 220 saat
        \item[C)] 270 saat
        \item[D)] 300 saat
        \item[E)] Hiçbiri
    \end{itemize}
    \tcblower
    \textbf{Doğru Cevap: C} \\
    \textbf{Açıklama:} Fazla çalışma süresinin toplamı bir yılda 270 saatten fazla olamaz.
\end{testcard}

\begin{testcard}{Soru 9}
    İş Kanunu'na göre günlük normal dinlenme saatleri dışında kalan gece çalışmaları en fazla kaç saati aşamaz?
    \begin{itemize}
        \item[A)] 6 saat
        \item[B)] 7.5 saat
        \item[C)] 9 saat
        \item[D)] 11 saat
        \item[E)] Hiçbiri
    \end{itemize}
    \tcblower
    \textbf{Doğru Cevap: B} \\
    \textbf{Açıklama:} İş Kanunu Madde 69'a göre, gece çalışmaları günde yedi buçuk (7.5) saati geçemez.
\end{testcard}

\begin{testcard}{Soru 10}
    İş Kanunu'na göre, 4 saatten fazla ve 7.5 saate kadar (7.5 saat dahil) süren işlerde verilmesi gereken ara dinlenmesi asgari ne kadardır?
    \begin{itemize}
        \item[A)] 15 dakika
        \item[B)] 30 dakika
        \item[C)] 45 dakika
        \item[D)] 60 dakika
        \item[E)] Hiçbiri
    \end{itemize}
    \tcblower
    \textbf{Doğru Cevap: B} \\
    \textbf{Açıklama:} Günde 4 ila 7.5 saat arası süren işlerde ara dinlenmesi (mola) asgari 30 dakikadır.
\end{testcard}

\begin{testcard}{Soru 11}
    İş Kanunu'na göre, 7.5 saatten fazla süren işlerde verilmesi gereken ara dinlenmesi asgari ne kadardır?
    \begin{itemize}
        \item[A)] 30 dakika
        \item[B)] 45 dakika
        \item[C)] 60 dakika (1 saat)
        \item[D)] 90 dakika
        \item[E)] Hiçbiri
    \end{itemize}
    \tcblower
    \textbf{Doğru Cevap: C} \\
    \textbf{Açıklama:} 7.5 saati aşan günlük çalışmalarda ara dinlenmesi asgari 1 saat olmak zorundadır.
\end{testcard}

\begin{testcard}{Soru 12}
    İş Kanunu'na göre, 4 saat veya daha kısa süren işlerde verilmesi gereken ara dinlenmesi asgari ne kadardır?
    \begin{itemize}
        \item[A)] 10 dakika
        \item[B)] 15 dakika
        \item[C)] 20 dakika
        \item[D)] Ara dinlenmesi gerekmez
        \item[E)] Hiçbiri
    \end{itemize}
    \tcblower
    \textbf{Doğru Cevap: B} \\
    \textbf{Açıklama:} 4 saat ve daha kısa süren işlerde ara dinlenmesi asgari 15 dakikadır.
\end{testcard}

\begin{testcard}{Soru 13}
    İSG Kanunu uyarınca, 'Çok Tehlikeli' sınıftaki işyerinde çalışanların periyodik sağlık raporları en geç kaç yılda bir yenilenmelidir?
    \begin{itemize}
        \item[A)] Yılda bir
        \item[B)] 2 yılda bir
        \item[C)] 3 yılda bir
        \item[D)] 5 yılda bir
        \item[E)] Hiçbiri
    \end{itemize}
    \tcblower
    \textbf{Doğru Cevap: A} \\
    \textbf{Açıklama:} Çok Tehlikeli sınıftaki işyerlerinde çalışanların periyodik sağlık raporları yılda en az bir kez yenilenmelidir.
\end{testcard}

\begin{testcard}{Soru 14}
    İSG Kanunu uyarınca, 'Tehlikeli' sınıftaki işyerinde çalışanların periyodik sağlık raporları en geç kaç yılda bir yenilenmelidir?
    \begin{itemize}
        \item[A)] Yılda bir
        \item[B)] 2 yılda bir
        \item[C)] 3 yılda bir
        \item[D)] 5 yılda bir
        \item[E)] Hiçbiri
    \end{itemize}
    \tcblower
    \textbf{Doğru Cevap: C} \\
    \textbf{Açıklama:} Tehlikeli sınıftaki işyerlerinde periyodik sağlık raporları en geç 3 yılda bir yenilenmek zorundadır.
\end{testcard}

\begin{testcard}{Soru 15}
    İSG Kanunu uyarınca, 'Az Tehlikeli' sınıftaki işyerinde çalışanların periyodik sağlık raporları en geç kaç yılda bir yenilenmelidir?
    \begin{itemize}
        \item[A)] Yılda bir
        \item[B)] 3 yılda bir
        \item[C)] 5 yılda bir
        \item[D)] 7 yılda bir
        \item[E)] Hiçbiri
    \end{itemize}
    \tcblower
    \textbf{Doğru Cevap: C} \\
    \textbf{Açıklama:} Az Tehlikeli sınıftaki işyerlerinde periyodik sağlık raporları en geç 5 yılda bir yenilenir.
\end{testcard}

\begin{testcard}{Soru 16}
    4857 sayılı İş Kanunu'na göre, temel çocuk çalıştırma yaş sınırı kaçtır?
    \begin{itemize}
        \item[A)] 12 yaşını doldurmamış olmak
        \item[B)] 15 yaşını doldurmamış çocukların çalıştırılması yasaktır
        \item[C)] 18 yaşını doldurana kadar her iş yasaktır
        \item[D)] Yaş sınırı yoktur
        \item[E)] Hiçbiri
    \end{itemize}
    \tcblower
    \textbf{Doğru Cevap: B} \\
    \textbf{Açıklama:} İş Kanunu md. 71 uyarınca 15 yaşını doldurmamış çocukların çalıştırılması yasaktır.
\end{testcard}

\begin{testcard}{Soru 17}
    Okula devam eden çocukların (14 yaşını doldurmuş) çalışma saatleri hakkında hangisi doğrudur?
    \begin{itemize}
        \item[A)] Günde en fazla 8 saat çalışabilirler
        \item[B)] Eğitim saatleri dışında olmak üzere günde en fazla 2 saat ve haftada 10 saat olabilir
        \item[C)] Gece mesailerine kalabilirler
        \item[D)] Okul zamanı çalıştırılmaları tamamen yasaktır
        \item[E)] Hiçbiri
    \end{itemize}
    \tcblower
    \textbf{Doğru Cevap: B} \\
    \textbf{Açıklama:} Eğitime devam eden çocukların ders dışındaki çalışma süreleri günde en fazla 2 saat, haftada ise 10 saat olabilir.
\end{testcard}

\begin{testcard}{Soru 18}
    İş Kanunu kapsamındaki çalışma sürelerine göre, haftalık normal çalışma süresi olan 45 saat, işyerlerinde haftanın çalışılan günlerine nasıl bölüştürülür? (Örnek No: 1)
    \begin{itemize}
        \item[A)] Aksi kararlaştırılmamışsa, çalışılan günlere eşit ölçüde bölünerek uygulanır
        \item[B)] Sadece günde 15 saat olarak 3 günde bitirilir
        \item[C)] İşçinin isteğine göre serbestçe dağıtılır
        \item[D)] Sadece gece mesaisi olarak yaptırılır
        \item[E)] Hiçbiri
    \end{itemize}
    \tcblower
    \textbf{Doğru Cevap: A} \\
    \textbf{Açıklama:} İş Kanunu uyarınca haftalık çalışma süresi aksi kararlaştırılmadıkça haftanın çalışılan günlerine eşit olarak dağıtılır.
\end{testcard}

\begin{testcard}{Soru 19}
    İş Kanunu kapsamındaki çalışma sürelerine göre, haftalık normal çalışma süresi olan 45 saat, işyerlerinde haftanın çalışılan günlerine nasıl bölüştürülür? (Örnek No: 2)
    \begin{itemize}
        \item[A)] Aksi kararlaştırılmamışsa, çalışılan günlere eşit ölçüde bölünerek uygulanır
        \item[B)] Sadece günde 15 saat olarak 3 günde bitirilir
        \item[C)] İşçinin isteğine göre serbestçe dağıtılır
        \item[D)] Sadece gece mesaisi olarak yaptırılır
        \item[E)] Hiçbiri
    \end{itemize}
    \tcblower
    \textbf{Doğru Cevap: A} \\
    \textbf{Açıklama:} İş Kanunu uyarınca haftalık çalışma süresi aksi kararlaştırılmadıkça haftanın çalışılan günlerine eşit olarak dağıtılır.
\end{testcard}

\begin{testcard}{Soru 20}
    İş Kanunu kapsamındaki çalışma sürelerine göre, haftalık normal çalışma süresi olan 45 saat, işyerlerinde haftanın çalışılan günlerine nasıl bölüştürülür? (Örnek No: 3)
    \begin{itemize}
        \item[A)] Aksi kararlaştırılmamışsa, çalışılan günlere eşit ölçüde bölünerek uygulanır
        \item[B)] Sadece günde 15 saat olarak 3 günde bitirilir
        \item[C)] İşçinin isteğine göre serbestçe dağıtılır
        \item[D)] Sadece gece mesaisi olarak yaptırılır
        \item[E)] Hiçbiri
    \end{itemize}
    \tcblower
    \textbf{Doğru Cevap: A} \\
    \textbf{Açıklama:} İş Kanunu uyarınca haftalık çalışma süresi aksi kararlaştırılmadıkça haftanın çalışılan günlerine eşit olarak dağıtılır.
\end{testcard}

\begin{testcard}{Soru 21}
    İş Kanunu kapsamındaki çalışma sürelerine göre, haftalık normal çalışma süresi olan 45 saat, işyerlerinde haftanın çalışılan günlerine nasıl bölüştürülür? (Örnek No: 4)
    \begin{itemize}
        \item[A)] Aksi kararlaştırılmamışsa, çalışılan günlere eşit ölçüde bölünerek uygulanır
        \item[B)] Sadece günde 15 saat olarak 3 günde bitirilir
        \item[C)] İşçinin isteğine göre serbestçe dağıtılır
        \item[D)] Sadece gece mesaisi olarak yaptırılır
        \item[E)] Hiçbiri
    \end{itemize}
    \tcblower
    \textbf{Doğru Cevap: A} \\
    \textbf{Açıklama:} İş Kanunu uyarınca haftalık çalışma süresi aksi kararlaştırılmadıkça haftanın çalışılan günlerine eşit olarak dağıtılır.
\end{testcard}

\begin{testcard}{Soru 22}
    İş Kanunu kapsamındaki çalışma sürelerine göre, haftalık normal çalışma süresi olan 45 saat, işyerlerinde haftanın çalışılan günlerine nasıl bölüştürülür? (Örnek No: 5)
    \begin{itemize}
        \item[A)] Aksi kararlaştırılmamışsa, çalışılan günlere eşit ölçüde bölünerek uygulanır
        \item[B)] Sadece günde 15 saat olarak 3 günde bitirilir
        \item[C)] İşçinin isteğine göre serbestçe dağıtılır
        \item[D)] Sadece gece mesaisi olarak yaptırılır
        \item[E)] Hiçbiri
    \end{itemize}
    \tcblower
    \textbf{Doğru Cevap: A} \\
    \textbf{Açıklama:} İş Kanunu uyarınca haftalık çalışma süresi aksi kararlaştırılmadıkça haftanın çalışılan günlerine eşit olarak dağıtılır.
\end{testcard}

\begin{testcard}{Soru 23}
    İş Kanunu kapsamındaki çalışma sürelerine göre, haftalık normal çalışma süresi olan 45 saat, işyerlerinde haftanın çalışılan günlerine nasıl bölüştürülür? (Örnek No: 6)
    \begin{itemize}
        \item[A)] Aksi kararlaştırılmamışsa, çalışılan günlere eşit ölçüde bölünerek uygulanır
        \item[B)] Sadece günde 15 saat olarak 3 günde bitirilir
        \item[C)] İşçinin isteğine göre serbestçe dağıtılır
        \item[D)] Sadece gece mesaisi olarak yaptırılır
        \item[E)] Hiçbiri
    \end{itemize}
    \tcblower
    \textbf{Doğru Cevap: A} \\
    \textbf{Açıklama:} İş Kanunu uyarınca haftalık çalışma süresi aksi kararlaştırılmadıkça haftanın çalışılan günlerine eşit olarak dağıtılır.
\end{testcard}

\begin{testcard}{Soru 24}
    İş Kanunu kapsamındaki çalışma sürelerine göre, haftalık normal çalışma süresi olan 45 saat, işyerlerinde haftanın çalışılan günlerine nasıl bölüştürülür? (Örnek No: 7)
    \begin{itemize}
        \item[A)] Aksi kararlaştırılmamışsa, çalışılan günlere eşit ölçüde bölünerek uygulanır
        \item[B)] Sadece günde 15 saat olarak 3 günde bitirilir
        \item[C)] İşçinin isteğine göre serbestçe dağıtılır
        \item[D)] Sadece gece mesaisi olarak yaptırılır
        \item[E)] Hiçbiri
    \end{itemize}
    \tcblower
    \textbf{Doğru Cevap: A} \\
    \textbf{Açıklama:} İş Kanunu uyarınca haftalık çalışma süresi aksi kararlaştırılmadıkça haftanın çalışılan günlerine eşit olarak dağıtılır.
\end{testcard}

\begin{testcard}{Soru 25}
    İş Kanunu kapsamındaki çalışma sürelerine göre, haftalık normal çalışma süresi olan 45 saat, işyerlerinde haftanın çalışılan günlerine nasıl bölüştürülür? (Örnek No: 8)
    \begin{itemize}
        \item[A)] Aksi kararlaştırılmamışsa, çalışılan günlere eşit ölçüde bölünerek uygulanır
        \item[B)] Sadece günde 15 saat olarak 3 günde bitirilir
        \item[C)] İşçinin isteğine göre serbestçe dağıtılır
        \item[D)] Sadece gece mesaisi olarak yaptırılır
        \item[E)] Hiçbiri
    \end{itemize}
    \tcblower
    \textbf{Doğru Cevap: A} \\
    \textbf{Açıklama:} İş Kanunu uyarınca haftalık çalışma süresi aksi kararlaştırılmadıkça haftanın çalışılan günlerine eşit olarak dağıtılır.
\end{testcard}

\begin{testcard}{Soru 26}
    İş Kanunu kapsamındaki çalışma sürelerine göre, haftalık normal çalışma süresi olan 45 saat, işyerlerinde haftanın çalışılan günlerine nasıl bölüştürülür? (Örnek No: 9)
    \begin{itemize}
        \item[A)] Aksi kararlaştırılmamışsa, çalışılan günlere eşit ölçüde bölünerek uygulanır
        \item[B)] Sadece günde 15 saat olarak 3 günde bitirilir
        \item[C)] İşçinin isteğine göre serbestçe dağıtılır
        \item[D)] Sadece gece mesaisi olarak yaptırılır
        \item[E)] Hiçbiri
    \end{itemize}
    \tcblower
    \textbf{Doğru Cevap: A} \\
    \textbf{Açıklama:} İş Kanunu uyarınca haftalık çalışma süresi aksi kararlaştırılmadıkça haftanın çalışılan günlerine eşit olarak dağıtılır.
\end{testcard}

\begin{testcard}{Soru 27}
    İş Kanunu kapsamındaki çalışma sürelerine göre, haftalık normal çalışma süresi olan 45 saat, işyerlerinde haftanın çalışılan günlerine nasıl bölüştürülür? (Örnek No: 10)
    \begin{itemize}
        \item[A)] Aksi kararlaştırılmamışsa, çalışılan günlere eşit ölçüde bölünerek uygulanır
        \item[B)] Sadece günde 15 saat olarak 3 günde bitirilir
        \item[C)] İşçinin isteğine göre serbestçe dağıtılır
        \item[D)] Sadece gece mesaisi olarak yaptırılır
        \item[E)] Hiçbiri
    \end{itemize}
    \tcblower
    \textbf{Doğru Cevap: A} \\
    \textbf{Açıklama:} İş Kanunu uyarınca haftalık çalışma süresi aksi kararlaştırılmadıkça haftanın çalışılan günlerine eşit olarak dağıtılır.
\end{testcard}

\begin{testcard}{Soru 28}
    İş Kanunu kapsamındaki çalışma sürelerine göre, haftalık normal çalışma süresi olan 45 saat, işyerlerinde haftanın çalışılan günlerine nasıl bölüştürülür? (Örnek No: 11)
    \begin{itemize}
        \item[A)] Aksi kararlaştırılmamışsa, çalışılan günlere eşit ölçüde bölünerek uygulanır
        \item[B)] Sadece günde 15 saat olarak 3 günde bitirilir
        \item[C)] İşçinin isteğine göre serbestçe dağıtılır
        \item[D)] Sadece gece mesaisi olarak yaptırılır
        \item[E)] Hiçbiri
    \end{itemize}
    \tcblower
    \textbf{Doğru Cevap: A} \\
    \textbf{Açıklama:} İş Kanunu uyarınca haftalık çalışma süresi aksi kararlaştırılmadıkça haftanın çalışılan günlerine eşit olarak dağıtılır.
\end{testcard}

\begin{testcard}{Soru 29}
    İş Kanunu kapsamındaki çalışma sürelerine göre, haftalık normal çalışma süresi olan 45 saat, işyerlerinde haftanın çalışılan günlerine nasıl bölüştürülür? (Örnek No: 12)
    \begin{itemize}
        \item[A)] Aksi kararlaştırılmamışsa, çalışılan günlere eşit ölçüde bölünerek uygulanır
        \item[B)] Sadece günde 15 saat olarak 3 günde bitirilir
        \item[C)] İşçinin isteğine göre serbestçe dağıtılır
        \item[D)] Sadece gece mesaisi olarak yaptırılır
        \item[E)] Hiçbiri
    \end{itemize}
    \tcblower
    \textbf{Doğru Cevap: A} \\
    \textbf{Açıklama:} İş Kanunu uyarınca haftalık çalışma süresi aksi kararlaştırılmadıkça haftanın çalışılan günlerine eşit olarak dağıtılır.
\end{testcard}

\begin{testcard}{Soru 30}
    İş Kanunu kapsamındaki çalışma sürelerine göre, haftalık normal çalışma süresi olan 45 saat, işyerlerinde haftanın çalışılan günlerine nasıl bölüştürülür? (Örnek No: 13)
    \begin{itemize}
        \item[A)] Aksi kararlaştırılmamışsa, çalışılan günlere eşit ölçüde bölünerek uygulanır
        \item[B)] Sadece günde 15 saat olarak 3 günde bitirilir
        \item[C)] İşçinin isteğine göre serbestçe dağıtılır
        \item[D)] Sadece gece mesaisi olarak yaptırılır
        \item[E)] Hiçbiri
    \end{itemize}
    \tcblower
    \textbf{Doğru Cevap: A} \\
    \textbf{Açıklama:} İş Kanunu uyarınca haftalık çalışma süresi aksi kararlaştırılmadıkça haftanın çalışılan günlerine eşit olarak dağıtılır.
\end{testcard}

\begin{testcard}{Soru 31}
    İş Kanunu kapsamındaki çalışma sürelerine göre, haftalık normal çalışma süresi olan 45 saat, işyerlerinde haftanın çalışılan günlerine nasıl bölüştürülür? (Örnek No: 14)
    \begin{itemize}
        \item[A)] Aksi kararlaştırılmamışsa, çalışılan günlere eşit ölçüde bölünerek uygulanır
        \item[B)] Sadece günde 15 saat olarak 3 günde bitirilir
        \item[C)] İşçinin isteğine göre serbestçe dağıtılır
        \item[D)] Sadece gece mesaisi olarak yaptırılır
        \item[E)] Hiçbiri
    \end{itemize}
    \tcblower
    \textbf{Doğru Cevap: A} \\
    \textbf{Açıklama:} İş Kanunu uyarınca haftalık çalışma süresi aksi kararlaştırılmadıkça haftanın çalışılan günlerine eşit olarak dağıtılır.
\end{testcard}

\begin{testcard}{Soru 32}
    İş Kanunu kapsamındaki çalışma sürelerine göre, haftalık normal çalışma süresi olan 45 saat, işyerlerinde haftanın çalışılan günlerine nasıl bölüştürülür? (Örnek No: 15)
    \begin{itemize}
        \item[A)] Aksi kararlaştırılmamışsa, çalışılan günlere eşit ölçüde bölünerek uygulanır
        \item[B)] Sadece günde 15 saat olarak 3 günde bitirilir
        \item[C)] İşçinin isteğine göre serbestçe dağıtılır
        \item[D)] Sadece gece mesaisi olarak yaptırılır
        \item[E)] Hiçbiri
    \end{itemize}
    \tcblower
    \textbf{Doğru Cevap: A} \\
    \textbf{Açıklama:} İş Kanunu uyarınca haftalık çalışma süresi aksi kararlaştırılmadıkça haftanın çalışılan günlerine eşit olarak dağıtılır.
\end{testcard}

\begin{testcard}{Soru 33}
    İş Kanunu kapsamındaki çalışma sürelerine göre, haftalık normal çalışma süresi olan 45 saat, işyerlerinde haftanın çalışılan günlerine nasıl bölüştürülür? (Örnek No: 16)
    \begin{itemize}
        \item[A)] Aksi kararlaştırılmamışsa, çalışılan günlere eşit ölçüde bölünerek uygulanır
        \item[B)] Sadece günde 15 saat olarak 3 günde bitirilir
        \item[C)] İşçinin isteğine göre serbestçe dağıtılır
        \item[D)] Sadece gece mesaisi olarak yaptırılır
        \item[E)] Hiçbiri
    \end{itemize}
    \tcblower
    \textbf{Doğru Cevap: A} \\
    \textbf{Açıklama:} İş Kanunu uyarınca haftalık çalışma süresi aksi kararlaştırılmadıkça haftanın çalışılan günlerine eşit olarak dağıtılır.
\end{testcard}

\begin{testcard}{Soru 34}
    İş Kanunu kapsamındaki çalışma sürelerine göre, haftalık normal çalışma süresi olan 45 saat, işyerlerinde haftanın çalışılan günlerine nasıl bölüştürülür? (Örnek No: 17)
    \begin{itemize}
        \item[A)] Aksi kararlaştırılmamışsa, çalışılan günlere eşit ölçüde bölünerek uygulanır
        \item[B)] Sadece günde 15 saat olarak 3 günde bitirilir
        \item[C)] İşçinin isteğine göre serbestçe dağıtılır
        \item[D)] Sadece gece mesaisi olarak yaptırılır
        \item[E)] Hiçbiri
    \end{itemize}
    \tcblower
    \textbf{Doğru Cevap: A} \\
    \textbf{Açıklama:} İş Kanunu uyarınca haftalık çalışma süresi aksi kararlaştırılmadıkça haftanın çalışılan günlerine eşit olarak dağıtılır.
\end{testcard}

\begin{testcard}{Soru 35}
    İş Kanunu kapsamındaki çalışma sürelerine göre, haftalık normal çalışma süresi olan 45 saat, işyerlerinde haftanın çalışılan günlerine nasıl bölüştürülür? (Örnek No: 18)
    \begin{itemize}
        \item[A)] Aksi kararlaştırılmamışsa, çalışılan günlere eşit ölçüde bölünerek uygulanır
        \item[B)] Sadece günde 15 saat olarak 3 günde bitirilir
        \item[C)] İşçinin isteğine göre serbestçe dağıtılır
        \item[D)] Sadece gece mesaisi olarak yaptırılır
        \item[E)] Hiçbiri
    \end{itemize}
    \tcblower
    \textbf{Doğru Cevap: A} \\
    \textbf{Açıklama:} İş Kanunu uyarınca haftalık çalışma süresi aksi kararlaştırılmadıkça haftanın çalışılan günlerine eşit olarak dağıtılır.
\end{testcard}

\begin{testcard}{Soru 36}
    İş Kanunu kapsamındaki çalışma sürelerine göre, haftalık normal çalışma süresi olan 45 saat, işyerlerinde haftanın çalışılan günlerine nasıl bölüştürülür? (Örnek No: 19)
    \begin{itemize}
        \item[A)] Aksi kararlaştırılmamışsa, çalışılan günlere eşit ölçüde bölünerek uygulanır
        \item[B)] Sadece günde 15 saat olarak 3 günde bitirilir
        \item[C)] İşçinin isteğine göre serbestçe dağıtılır
        \item[D)] Sadece gece mesaisi olarak yaptırılır
        \item[E)] Hiçbiri
    \end{itemize}
    \tcblower
    \textbf{Doğru Cevap: A} \\
    \textbf{Açıklama:} İş Kanunu uyarınca haftalık çalışma süresi aksi kararlaştırılmadıkça haftanın çalışılan günlerine eşit olarak dağıtılır.
\end{testcard}

\begin{testcard}{Soru 37}
    İş Kanunu kapsamındaki çalışma sürelerine göre, haftalık normal çalışma süresi olan 45 saat, işyerlerinde haftanın çalışılan günlerine nasıl bölüştürülür? (Örnek No: 20)
    \begin{itemize}
        \item[A)] Aksi kararlaştırılmamışsa, çalışılan günlere eşit ölçüde bölünerek uygulanır
        \item[B)] Sadece günde 15 saat olarak 3 günde bitirilir
        \item[C)] İşçinin isteğine göre serbestçe dağıtılır
        \item[D)] Sadece gece mesaisi olarak yaptırılır
        \item[E)] Hiçbiri
    \end{itemize}
    \tcblower
    \textbf{Doğru Cevap: A} \\
    \textbf{Açıklama:} İş Kanunu uyarınca haftalık çalışma süresi aksi kararlaştırılmadıkça haftanın çalışılan günlerine eşit olarak dağıtılır.
\end{testcard}

\begin{testcard}{Soru 38}
    İş Kanunu kapsamındaki çalışma sürelerine göre, haftalık normal çalışma süresi olan 45 saat, işyerlerinde haftanın çalışılan günlerine nasıl bölüştürülür? (Örnek No: 21)
    \begin{itemize}
        \item[A)] Aksi kararlaştırılmamışsa, çalışılan günlere eşit ölçüde bölünerek uygulanır
        \item[B)] Sadece günde 15 saat olarak 3 günde bitirilir
        \item[C)] İşçinin isteğine göre serbestçe dağıtılır
        \item[D)] Sadece gece mesaisi olarak yaptırılır
        \item[E)] Hiçbiri
    \end{itemize}
    \tcblower
    \textbf{Doğru Cevap: A} \\
    \textbf{Açıklama:} İş Kanunu uyarınca haftalık çalışma süresi aksi kararlaştırılmadıkça haftanın çalışılan günlerine eşit olarak dağıtılır.
\end{testcard}

\begin{testcard}{Soru 39}
    İş Kanunu kapsamındaki çalışma sürelerine göre, haftalık normal çalışma süresi olan 45 saat, işyerlerinde haftanın çalışılan günlerine nasıl bölüştürülür? (Örnek No: 22)
    \begin{itemize}
        \item[A)] Aksi kararlaştırılmamışsa, çalışılan günlere eşit ölçüde bölünerek uygulanır
        \item[B)] Sadece günde 15 saat olarak 3 günde bitirilir
        \item[C)] İşçinin isteğine göre serbestçe dağıtılır
        \item[D)] Sadece gece mesaisi olarak yaptırılır
        \item[E)] Hiçbiri
    \end{itemize}
    \tcblower
    \textbf{Doğru Cevap: A} \\
    \textbf{Açıklama:} İş Kanunu uyarınca haftalık çalışma süresi aksi kararlaştırılmadıkça haftanın çalışılan günlerine eşit olarak dağıtılır.
\end{testcard}

\begin{testcard}{Soru 40}
    İş Kanunu kapsamındaki çalışma sürelerine göre, haftalık normal çalışma süresi olan 45 saat, işyerlerinde haftanın çalışılan günlerine nasıl bölüştürülür? (Örnek No: 23)
    \begin{itemize}
        \item[A)] Aksi kararlaştırılmamışsa, çalışılan günlere eşit ölçüde bölünerek uygulanır
        \item[B)] Sadece günde 15 saat olarak 3 günde bitirilir
        \item[C)] İşçinin isteğine göre serbestçe dağıtılır
        \item[D)] Sadece gece mesaisi olarak yaptırılır
        \item[E)] Hiçbiri
    \end{itemize}
    \tcblower
    \textbf{Doğru Cevap: A} \\
    \textbf{Açıklama:} İş Kanunu uyarınca haftalık çalışma süresi aksi kararlaştırılmadıkça haftanın çalışılan günlerine eşit olarak dağıtılır.
\end{testcard}

\subsection{Yavuz Hoca'nın Özel İSG Notları}
\begin{testcard}{Soru 1}
    Yavuz Hoca'nın ders notunda belirttiği genel İSG kurallarına göre, işyerinde İSG kültürünün yerleşmesini engelleyen en temel davranış hangisidir?
    \begin{itemize}
        \item[A)] Rutin sağlık muayenelerine katılmak
        \item[B)] Tehlikeli durumları kurula rapor etmek
        \item[C)] 'Bana bir şey olmaz' algısıyla kuralları göz ardı etmek
        \item[D)] Kişisel koruyucu donanım kullanmayı talep etmek
        \item[E)] Hiçbiri
    \end{itemize}
    \tcblower
    \textbf{Doğru Cevap: C} \\
    \textbf{Açıklama:} Bireysel kanıksama ve 'bana bir şey olmaz' algısı, güvenlik kültürünün yerleşmesini engelleyen en yaygın davranış biçimidir.
\end{testcard}

\begin{testcard}{Soru 2}
    Yavuz Hoca'nın ders notuna göre, işyerinde İSG farkındalığı oluşturmak için kullanılan en etkili iletişim kanalı hangisidir?
    \begin{itemize}
        \item[A)] Sadece yazılı cezalar kesmek
        \item[B)] Karşılıklı diyalog, eğitimler ve çalışanların görüşlerinin alındığı toplantılar
        \item[C)] Telefon mesajları
        \item[D)] Panolara sadece ingilizce uyarılar asmak
        \item[E)] Hiçbiri
    \end{itemize}
    \tcblower
    \textbf{Doğru Cevap: B} \\
    \textbf{Açıklama:} Çift yönlü iletişim ve çalışanların katılımı, İSG kültürünü kalıcı hale getirmenin en sağlıklı yoludur.
\end{testcard}

\begin{testcard}{Soru 3}
    Yavuz Hoca'nın özel ders notunda önemle belirttiği üzere, işyerinde risk kültürünü geliştirmek için yöneticilerin sergilemesi gereken davranış hangisidir? (Örnek No: 1)
    \begin{itemize}
        \item[A)] Kurallara bizzat uyarak çalışanlara rol model olmak ve geri bildirimleri dikkate almak
        \item[B)] Sadece uzaktan izleyip ceza yazmak
        \item[C)] İSG bütçesini kısmak
        \item[D)] Baret takmayı isteğe bağlı kılmak
        \item[E)] Hiçbiri
    \end{itemize}
    \tcblower
    \textbf{Doğru Cevap: A} \\
    \textbf{Açıklama:} Yönetimin İSG'ye liderlik etmesi ve kurallara uyması, çalışanların da bu kültürü benimsemesini sağlayan en önemli faktördür.
\end{testcard}

\begin{testcard}{Soru 4}
    Yavuz Hoca'nın özel ders notunda önemle belirttiği üzere, işyerinde risk kültürünü geliştirmek için yöneticilerin sergilemesi gereken davranış hangisidir? (Örnek No: 2)
    \begin{itemize}
        \item[A)] Kurallara bizzat uyarak çalışanlara rol model olmak ve geri bildirimleri dikkate almak
        \item[B)] Sadece uzaktan izleyip ceza yazmak
        \item[C)] İSG bütçesini kısmak
        \item[D)] Baret takmayı isteğe bağlı kılmak
        \item[E)] Hiçbiri
    \end{itemize}
    \tcblower
    \textbf{Doğru Cevap: A} \\
    \textbf{Açıklama:} Yönetimin İSG'ye liderlik etmesi ve kurallara uyması, çalışanların da bu kültürü benimsemesini sağlayan en önemli faktördür.
\end{testcard}

\begin{testcard}{Soru 5}
    Yavuz Hoca'nın özel ders notunda önemle belirttiği üzere, işyerinde risk kültürünü geliştirmek için yöneticilerin sergilemesi gereken davranış hangisidir? (Örnek No: 3)
    \begin{itemize}
        \item[A)] Kurallara bizzat uyarak çalışanlara rol model olmak ve geri bildirimleri dikkate almak
        \item[B)] Sadece uzaktan izleyip ceza yazmak
        \item[C)] İSG bütçesini kısmak
        \item[D)] Baret takmayı isteğe bağlı kılmak
        \item[E)] Hiçbiri
    \end{itemize}
    \tcblower
    \textbf{Doğru Cevap: A} \\
    \textbf{Açıklama:} Yönetimin İSG'ye liderlik etmesi ve kurallara uyması, çalışanların da bu kültürü benimsemesini sağlayan en önemli faktördür.
\end{testcard}

\begin{testcard}{Soru 6}
    Yavuz Hoca'nın özel ders notunda önemle belirttiği üzere, işyerinde risk kültürünü geliştirmek için yöneticilerin sergilemesi gereken davranış hangisidir? (Örnek No: 4)
    \begin{itemize}
        \item[A)] Kurallara bizzat uyarak çalışanlara rol model olmak ve geri bildirimleri dikkate almak
        \item[B)] Sadece uzaktan izleyip ceza yazmak
        \item[C)] İSG bütçesini kısmak
        \item[D)] Baret takmayı isteğe bağlı kılmak
        \item[E)] Hiçbiri
    \end{itemize}
    \tcblower
    \textbf{Doğru Cevap: A} \\
    \textbf{Açıklama:} Yönetimin İSG'ye liderlik etmesi ve kurallara uyması, çalışanların da bu kültürü benimsemesini sağlayan en önemli faktördür.
\end{testcard}

\begin{testcard}{Soru 7}
    Yavuz Hoca'nın özel ders notunda önemle belirttiği üzere, işyerinde risk kültürünü geliştirmek için yöneticilerin sergilemesi gereken davranış hangisidir? (Örnek No: 5)
    \begin{itemize}
        \item[A)] Kurallara bizzat uyarak çalışanlara rol model olmak ve geri bildirimleri dikkate almak
        \item[B)] Sadece uzaktan izleyip ceza yazmak
        \item[C)] İSG bütçesini kısmak
        \item[D)] Baret takmayı isteğe bağlı kılmak
        \item[E)] Hiçbiri
    \end{itemize}
    \tcblower
    \textbf{Doğru Cevap: A} \\
    \textbf{Açıklama:} Yönetimin İSG'ye liderlik etmesi ve kurallara uyması, çalışanların da bu kültürü benimsemesini sağlayan en önemli faktördür.
\end{testcard}

\begin{testcard}{Soru 8}
    Yavuz Hoca'nın özel ders notunda önemle belirttiği üzere, işyerinde risk kültürünü geliştirmek için yöneticilerin sergilemesi gereken davranış hangisidir? (Örnek No: 6)
    \begin{itemize}
        \item[A)] Kurallara bizzat uyarak çalışanlara rol model olmak ve geri bildirimleri dikkate almak
        \item[B)] Sadece uzaktan izleyip ceza yazmak
        \item[C)] İSG bütçesini kısmak
        \item[D)] Baret takmayı isteğe bağlı kılmak
        \item[E)] Hiçbiri
    \end{itemize}
    \tcblower
    \textbf{Doğru Cevap: A} \\
    \textbf{Açıklama:} Yönetimin İSG'ye liderlik etmesi ve kurallara uyması, çalışanların da bu kültürü benimsemesini sağlayan en önemli faktördür.
\end{testcard}

\begin{testcard}{Soru 9}
    Yavuz Hoca'nın özel ders notunda önemle belirttiği üzere, işyerinde risk kültürünü geliştirmek için yöneticilerin sergilemesi gereken davranış hangisidir? (Örnek No: 7)
    \begin{itemize}
        \item[A)] Kurallara bizzat uyarak çalışanlara rol model olmak ve geri bildirimleri dikkate almak
        \item[B)] Sadece uzaktan izleyip ceza yazmak
        \item[C)] İSG bütçesini kısmak
        \item[D)] Baret takmayı isteğe bağlı kılmak
        \item[E)] Hiçbiri
    \end{itemize}
    \tcblower
    \textbf{Doğru Cevap: A} \\
    \textbf{Açıklama:} Yönetimin İSG'ye liderlik etmesi ve kurallara uyması, çalışanların da bu kültürü benimsemesini sağlayan en önemli faktördür.
\end{testcard}

\begin{testcard}{Soru 10}
    Yavuz Hoca'nın özel ders notunda önemle belirttiği üzere, işyerinde risk kültürünü geliştirmek için yöneticilerin sergilemesi gereken davranış hangisidir? (Örnek No: 8)
    \begin{itemize}
        \item[A)] Kurallara bizzat uyarak çalışanlara rol model olmak ve geri bildirimleri dikkate almak
        \item[B)] Sadece uzaktan izleyip ceza yazmak
        \item[C)] İSG bütçesini kısmak
        \item[D)] Baret takmayı isteğe bağlı kılmak
        \item[E)] Hiçbiri
    \end{itemize}
    \tcblower
    \textbf{Doğru Cevap: A} \\
    \textbf{Açıklama:} Yönetimin İSG'ye liderlik etmesi ve kurallara uyması, çalışanların da bu kültürü benimsemesini sağlayan en önemli faktördür.
\end{testcard}

\end{document}